% -----------------------------------------------------------
\section{1842}
% -----------------------------------------------------------

	% ----------------------
	\subsection{Joseph Smith}
	% ----------------------
		\label{EMS-1842-JS}

		% ----------------------
		\subsubsection{Introduction to 1842}
		% ----------------------
		
			sdf
			
		% -----
		\subsubsection{Letter to Edward Hunter, 5 January 1842}
		% -----
			\label{EMS-1842-JS-5-JAN-1842}
			% \label{EMS-1842-JS-DAY-3LETTERMONTH-YEAR}
			
			\begin{description}
				\item[Print Sources]
			\end{description}

			\begin{tabular}{p{0.5in}p{1in}p{0.5in}p{1in}}
				JSP	 	& ---		& JSR 		& --- \\
				DC	 	& ---		& HC 		& --- \\
				HDDC 	& ---		& TCHC 		& --- \\
			\end{tabular}
			% HDDC = woodford's DC diss; JSR = Marquardt's JS Revelations; TCHC = Vogel HC
			
			\begin{description}
				\item[Online Access] \href{https://www.josephsmithpapers.org/paper-summary/letter-to-edward-hunter-5-january-1842/1}{Here}
				% \item[Analysis] \hyperref[XX]{Here}
			\end{description}
			
			\begin{description}
				\item[Background]
			\end{description}
			
				% It might also be useful to give a two sentence context of the period: e.g., "This speech was given in Nauvoo to a small congregation one month prior to the destruction of the Nauvoo Expositor. These and other tensions may have been on the prophet's mind when he gave the speech."
			
			\begin{description}
				\item[Original Source]
			\end{description}
			
				% brief description of original source material, with reference to JSP or somewhere else for more info
				% Background material: it might be helpful to have a note that says whether a given document was presented as a speech, was written in a journal, etc.
				% Also a comment here about how reliable the source might be, or what shape it is in, if there are large omissions or parts that are hard to understand, and so on.
			
			\begin{description}
				\item[Text]
			\end{description}
			
				\begin{tightright}
						\hypertarget{EMS-1842-JS-5-JAN-1842-a}%
					Nauvoo
						\marginnote{a}%
					January 5[\supersu{th}] 1842
				\end{tightright}

				Mr Edward Hunter.

				Beloved Brother,

				I wrote you on the 21 ultimo, in reply to yours of the 27[\supersu{th}] of october, but lest by any means the letter should fail to reach you I will recapitulate very briefly some important items therein contained,

				The power of Attorney was duly executed by Mrs [Margaret] Smith. \& forwareded to the clerks office for seal of state. to be sent. from thence direct to you.

				The Goods are accepted and <will be> applied according to you[r] request,

					\hypertarget{EMS-1842-JS-5-JAN-1842-b}%
				I
					\marginnote{b}%
				have purchased 90 acres of woodland, a little up the River; have made proposals to [Hugh] McFall. but am yet waiting his answer, from his eastern correspondent.

				Steam Engines \& mills of any description will do well here, the more of such things you can bring, the better.--- for particulars on the foregoing I would refer you to my letter of the 21. ult which I hope you have received ere this---

				\sout{The}

					\hypertarget{EMS-1842-JS-5-JAN-1842-c}%
				I
					\marginnote{c}
				am happy that it is my privilige to say to you that the large New. Building. which I had commenced when you were here, is now completed, and the doors are opened this day. for the sale of goods for the first time. The foundations of the building is somewhat spacio[u]s, (as you will doubtless recollect,) for a country store, The principal part of the building below, which is nearly 10 feet high is devoted \sout{to} exclusively to Shelves. <\&> drawers, Except 1 door opening back into the space, on the left of which \sout{is} are the cellar \& chamber Stairs. \& on the Right the Counting Room;--- from the space at the top of the chamber stairs, opens a door into the Large front room, of the same size with the one below.---
					\hypertarget{EMS-1842-JS-5-JAN-1842-d}%
				the
					\marginnote{d}%
				walls lined with counters. coverd with reserve goods.--- in front of the stairs opens the door to my Private Office, or where I keep the sacred writings. with a window to the south. overlooking the River below. \& the opposite shore for a great distance, <which> together with the passage of boats in the season thereof, constitutes a peculiarly interesting situation, in prospect \& no less interesting from its retirement from the bustle \& confusion of the neighborhood \& city. and altogether is a place the Lord is pleased to bless.---

					\hypertarget{EMS-1842-JS-5-JAN-1842-e}%
				The
					\marginnote{e}%
				painting of the store has been ex[e]cuted by some of our English bethen [brethren].--- \& the counters, drawers--- \& pillars present a very respectble representation of Oak. Mahagony \& Marble--- for a back woods establishment,---

					\hypertarget{EMS-1842-JS-5-JAN-1842-f}%
				The
					\marginnote{f}%
				Lord has blessed our exertions in a wonderful manner, and although some individuals have succeded in detain[in]g goods. to a considerable amount for the time being, yet we have been enabled to s[e]cure goods in the building Sufficient to fill all the shelves. <\& soon as they were completed> \& have some in reserve, both in loft. \& cellar. Our assortment, is tolerably good--- very good considering the different purchases made by different individuals,--- at different times, and under circumstances which controuled their choice to some extent, but, I rejoice that we have been enabled to do as well as we have, for the hearts of many of the poor brethren \& sisters will be made glad, with those comforts which are now within their reach.
					\hypertarget{EMS-1842-JS-5-JAN-1842-g}%
				The
					\marginnote{g}%
				store has been filled to overflowing all day, \& I have stood behind the counter dealing out goods as steady as any clerk you ever Saw to oblige those who were compelled to go without their <usual> christmas \& New year, dinners. for the want of a little Sugar, Molasses, Rasions \&c. \&c,--- \& to please myself also for I love to wait upon the Saints, and be a servant to all hoping that I may be exalted in the due time of the Lord.

					\hypertarget{EMS-1842-JS-5-JAN-1842-h}%
				It
					\marginnote{h}%
				is highly necessary that the store be well supplied with merchandise from this time forward, both for the interest of the church generally--- \& the comfort of the brethren individually and as expenses have been incurred already to a great amount in building the Store, Temple, ``Nauvoo house'' \&c \&c--- a great many of the goods on hand will have to pass away on orders previous contracts, \&c. \& we shall be obliged to lean upon other resources, <to a great extent> rather than the profits of goods, this winter, to supply a new stock in the spring, \& for this reason as well as those before stated, \& also, for your gratification in learning of our prosperity.
					\hypertarget{EMS-1842-JS-5-JAN-1842-i}%
				I
					\marginnote{i}%
				write you this early. to disire you to have the money you are to get on the power of Attorney, of Mrs Smith ready for disposal in Philadelphia, as soon as the rivers shall open, \& I sincerely hope \& trust that nothing will prevent your getting the money as you expect, so that it may be ready in depositee at Philadelphia, or \sout{or} so that you can meet Mr Whitny [Newel K. Whitney] <at Phil\supersu{a}\supers{.}''> or someone who may go for the goods, at a time which may be appointed hereafter.---
					\hypertarget{EMS-1842-JS-5-JAN-1842-j}%
				\&c
					\marginnote{j}%
				that we may have an early supply, of a spring selection, as you are aware that the first opening of a new assortment would be much to the advantage of the establishment, and I wish you to give me the earliest information possible, of any thing new, in relation to this, matter, With Sentiments of high consideration I remain your Brother in Christ,

				\begin{tightright}
					Joseph Smith

					pr. W[illard] Richards, Scribe
				\end{tightright}

				<NAUVOO Ills. JAN 11>

				\begin{tightright}
					<25>
				\end{tightright}

				Edward Hunter Esq\supersu{r}\supers{.}

				West Nantmeal

				Chester County

				Pensylvania
				
		% -----
		\subsubsection{Revelation, 28 January 1842}
		% -----
			\label{EMS-1842-JS-28-JAN-1842}
			% \label{EMS-1842-JS-DAY-3LETTERMONTH-YEAR}
			
			\begin{description}
				\item[Print Sources]
			\end{description}

			\begin{tabular}{p{0.5in}p{1in}p{0.5in}p{1in}}
				JSP	 	& ---		& JSR 		& --- \\
				DC	 	& ---		& HC 		& --- \\
				HDDC 	& ---		& TCHC 		& --- \\
			\end{tabular}
			% HDDC = woodford's DC diss; JSR = Marquardt's JS Revelations; TCHC = Vogel HC
			
			\begin{description}
				\item[Online Access] \href{https://www.josephsmithpapers.org/paper-summary/revelation-28-january-1842/1}{Here}
				% \item[Analysis] \hyperref[XX]{Here}
			\end{description}
			
			\begin{description}
				\item[Background]
			\end{description}
			
				% It might also be useful to give a two sentence context of the period: e.g., "This speech was given in Nauvoo to a small congregation one month prior to the destruction of the Nauvoo Expositor. These and other tensions may have been on the prophet's mind when he gave the speech."
			
			\begin{description}
				\item[Original Source]
			\end{description}
			
				% brief description of original source material, with reference to JSP or somewhere else for more info
				% Background material: it might be helpful to have a note that says whether a given document was presented as a speech, was written in a journal, etc.
				% Also a comment here about how reliable the source might be, or what shape it is in, if there are large omissions or parts that are hard to understand, and so on.
			
			\begin{description}
				\item[Text]
			\end{description}
			
				1842
				
				January 28.
				
				\begin{tightcenter}
					A Revelation \sout{of} <to> the twelve concrning the Times and Seasons.
				\end{tightcenter}
				
				Verily thus saith the Lord unto you my servant Joseph. go and say unto the Twelve That it is my will to have them take in hand the Editorial department of the Times and Seasons according to that manifestation. which Shall be given unto them by the Power of My Holy Spirit in the midst of their counsel Saith the Lord. Amen
				
		% -----
		\subsubsection{Discourse, 30 January 1842}
		% -----
			\label{EMS-1842-JS-30-JAN-1842}
			% \label{EMS-1842-JS-DAY-3LETTERMONTH-YEAR}
			
			% --
			\paragraph{As Reported by Wilford Woodruff}
			% --
				\label{EMS-1842-JS-30-JAN-1842-WW}
			
				\begin{description}
					\item[Print Sources]
				\end{description}

				\begin{tabular}{p{0.5in}p{1in}p{0.5in}p{1in}}
					JSP	 	& ---		& JSR 		& --- \\
					DC	 	& ---		& HC 		& --- \\
					HDDC 	& ---		& TCHC 		& --- \\
				\end{tabular}
				% HDDC = woodford's DC diss; JSR = Marquardt's JS Revelations; TCHC = Vogel HC
				
				\begin{description}
					\item[Online Access] \href{https://www.josephsmithpapers.org/paper-summary/discourse-30-january-1842/1}{Here}
					% \item[Analysis] \hyperref[XX]{Here}
				\end{description}
				
				\begin{description}
					\item[Background]
				\end{description}
				
					% It might also be useful to give a two sentence context of the period: e.g., "This speech was given in Nauvoo to a small congregation one month prior to the destruction of the Nauvoo Expositor. These and other tensions may have been on the prophet's mind when he gave the speech."
				
				\begin{description}
					\item[Original Source]
				\end{description}
				
					% brief description of original source material, with reference to JSP or somewhere else for more info
					% Background material: it might be helpful to have a note that says whether a given document was presented as a speech, was written in a journal, etc.
					% Also a comment here about how reliable the source might be, or what shape it is in, if there are large omissions or parts that are hard to understand, and so on.
				
				\begin{description}
					\item[Text]
				\end{description}
				
						\hypertarget{EMS-1842-JS-30-JAN-1842-WW-a}%
					Jan
						\marginnote{a}%
					30[\supersu{th}] [1842] Joseph the Seer taught the following principles that the God \& father of our Lord Jesus Christ was once the same as the Son or Holy Ghost bothaving [both having] redeemed a world became the eternal God of that world
						\hypertarget{EMS-1842-JS-30-JAN-1842-WW-b}%
					he
						\marginnote{b}%
					had a son Jesus Christ who redeemed this earth the same as his father had a world which made them equal \& the Holy Ghost would do the same in his turn \& so would all the Saints who inherited a Celestial glory so their would be Gods many \& Lords many
						\hypertarget{EMS-1842-JS-30-JAN-1842-WW-c}%
					their
						\marginnote{c}%
					were many mansions even 12 from the abode of Devils to the Celestial glory All Spirits that have bodies have power over those that have not hence men have power over Devils \&c
					
		% -----
		\subsubsection{``Church History,'' 1 March 1842}
		% -----
			\label{EMS-1842-JS-1-MAR-1842}
			% \label{EMS-1842-JS-DAY-3LETTERMONTH-YEAR}
			
			\begin{description}
				\item[Print Sources]
			\end{description}

			\begin{tabular}{p{0.5in}p{1in}p{0.5in}p{1in}}
				JSP	 	& ---		& JSR 		& --- \\
				DC	 	& ---		& HC 		& --- \\
				HDDC 	& ---		& TCHC 		& --- \\
			\end{tabular}
			% HDDC = woodford's DC diss; JSR = Marquardt's JS Revelations; TCHC = Vogel HC
			
			\begin{description}
				\item[Online Access] \href{https://www.josephsmithpapers.org/paper-summary/church-history-1-march-1842/1}{Here}
				% \item[Analysis] \hyperref[XX]{Here}
			\end{description}
			
			\begin{description}
				\item[Background]
			\end{description}
			
				% It might also be useful to give a two sentence context of the period: e.g., "This speech was given in Nauvoo to a small congregation one month prior to the destruction of the Nauvoo Expositor. These and other tensions may have been on the prophet's mind when he gave the speech."
			
			\begin{description}
				\item[Original Source]
			\end{description}
			
				% brief description of original source material, with reference to JSP or somewhere else for more info
				% Background material: it might be helpful to have a note that says whether a given document was presented as a speech, was written in a journal, etc.
				% Also a comment here about how reliable the source might be, or what shape it is in, if there are large omissions or parts that are hard to understand, and so on.
			
			\begin{description}
				\item[Text]
			\end{description}
			
				\begin{tightcenter}
						\hypertarget{EMS-1842-JS-1-MAR-1842-a}%
					CHURCH
						\marginnote{a}%
					HISTORY.
				\end{tightcenter}

				At the request of Mr. John Wentworth, Editor, and Proprietor of the ``Chicago Democrat,'' I have written the following sketch of the rise, progress, persecution, and faith of the Latter-Day Saints, of which I have the honor, under God, of being the founder. Mr. Wentworth says, that he wishes to furnish Mr. Bastow [George Barstow], a friend of his, who is writing the history of New Hampshire, with this document. As Mr. Bastow has taken the proper steps to obtain correct information all that I shall ask at his hands, is, that he publish the account entire, ungarnished, and without misrepresentation.

					\hypertarget{EMS-1842-JS-1-MAR-1842-b}%
				I
					\marginnote{b}%
				was born in the town of Sharon Windsor co., Vermont, on the 23d of December, A. D. 1805. When ten years old my parents removed to Palmyra New York, where we resided about four years, and from thence we removed to the town of Manchester.

					\hypertarget{EMS-1842-JS-1-MAR-1842-c}%
				My
					\marginnote{c}%
				father was a farmer and taught me the art of husbandry. When about fourteen years of age I began to reflect upon the importance of being prepared for a future state, and upon enquiring the plan of salvation I found that there was a great clash in religious sentiment; if I went to one society they referred me to one plan, and another to another; each one pointing to his own particular creed as the summum bonum of perfection: considering that all could not be right, and that God could not be the author of so much confusion I determined to investigate the subject more fully, believing that if God had a church it would not be split up into factions, and that if he taught one society to worship one way, and administer in one set of ordinances, he would not teach another principles which were diametrically opposed.
					\hypertarget{EMS-1842-JS-1-MAR-1842-d}%
				Believing
					\marginnote{d}%
				the word of God I had confidence in the declaration of James; ``If any man lack wisdom let him ask of God who giveth to all men liberally and upbraideth not and it shall be given him,'' I retired to a secret place in a grove and began to call upon the Lord, while fervently engaged in supplication my mind was taken away from the objects with which I was surrounded, and I was enwrapped in a heavenly vision and saw two glorious personages who exactly resembled each other in features, and likeness, surrounded with a brilliant light which eclipsed the sun at noon-day.
					\hypertarget{EMS-1842-JS-1-MAR-1842-e}%
				They
					\marginnote{e}%
				told me that all religious denominations were believing in incorrect doctrines, and that none of them was acknowledged of God as his church and kingdom. And I was expressly commanded to ``go not after them,'' at the same time receiving a promise that the fulness of the gospel should at some future time be made known unto me.

					\hypertarget{EMS-1842-JS-1-MAR-1842-f}%
				On
					\marginnote{f}%
				the evening of the 21st of September, A. D. 1823, while I was praying unto God, and endeavoring to exercise faith in the precious promises of scripture on a sudden a light like that of day, only of a far purer and more glorious appearance, and brightness burst into the room, indeed the first sight was as though the house was filled with consuming fire; the appearance produced a shock that affected the whole body; in a moment a personage stood before me surrounded with a glory yet greater than that with which I was already surrounded.
					\hypertarget{EMS-1842-JS-1-MAR-1842-g}%
				This
					\marginnote{g}%
				messenger proclaimed himself to be an angel of God sent to bring the joyful tidings, that the covenant which God made with ancient Israel was at hand to be fulfilled, that the preparatory work for the second coming of the Messiah was speedily to commence; that the time was at hand for the gospel, in all its fulness to be preached in power, unto all nations that a people might be prepared for the millennial reign.

				I was informed that I was chosen to be an instrument in the hands of God to bring about some of his purposes in this glorious dispensation.

					\hypertarget{EMS-1842-JS-1-MAR-1842-h}%
				I
					\marginnote{h}%
				was also informed concerning the aboriginal inhabitants of this country, and shown who they were, and from whence they came; a brief sketch of their origin, progress, civilization, laws, governments, of their righteousness and iniquity, and the blessings of God being finally withdrawn from them as a people was made known unto me: I was also told where there was deposited some plates on which were engraven an abridgement of the records of the ancient prophets that had existed on this continent.
					\hypertarget{EMS-1842-JS-1-MAR-1842-i}%
				The
					\marginnote{i}%
				angel appeared to me three times the same night and unfolded the same things. After having received many visits from the angels of God unfolding the majesty, and glory of the events that should transpire in the last days, on the morning of the 22d of September A. D. 1827, the angel of the Lord delivered the records into my hands.

					\hypertarget{EMS-1842-JS-1-MAR-1842-j}%
				These
					\marginnote{j}%
				records were engraven on plates which had the appearance of gold, each plate was six inches wide and eight inches long and not quite so thick as common tin. They were filled with engravings, in Egyptian characters and bound together in a volume, as the leaves of a book with three rings running through the whole. The volume was something near six inches in thickness, a part of which was sealed. The characters on the unsealed part were small, and beautifully engraved.
					\hypertarget{EMS-1842-JS-1-MAR-1842-k}%
				The
					\marginnote{k}%
				whole book exhibited many marks of antiquity in its construction and much skill in the art of engraving. With the records was found a curious instrument which the ancients called ``Urim and Thummim,'' which consisted of two transparent stones set in the rim of a bow fastened to a breastplate.

				Through the medium of the Urim and Thummim I translated the record by the gift, and power of God.

					\hypertarget{EMS-1842-JS-1-MAR-1842-l}%
				In
					\marginnote{l}%
				this important and interesting book the history of ancient America is unfolded, from its first settlement by a colony that came from the tower of Babel, at the confusion of languages to the beginning of the fifth century of the Christian era. We are informed by these records that America in ancient times has been inhabited by two distinct races of people. The first were called Jaredites and came directly from the tower of Babel. The second race came directly from the city of Jerusalem, about six hundred years before Christ.
					\hypertarget{EMS-1842-JS-1-MAR-1842-m}%
				They
					\marginnote{m}%
				were principally Israelites, of the descendants of Joseph. The Jaredites were destroyed about the time that the Israelites came from Jerusalem, who succeeded them in the inheritance of the country. The principal nation of the second race fell in battle towards the close of the fourth century. The remnant are the Indians that now inhabit this country. This book also tells us that our Saviour made his appearance upon this continent after his resurrection, that he planted the gospel here in all its fulness, and richness, and power, and blessing; that they had apostles, prophets, pastors, teachers and evangelists;
					\hypertarget{EMS-1842-JS-1-MAR-1842-n}%
				the
					\marginnote{n}%
				same order, the same priesthood, the same ordinances, gifts, powers, and blessing, as was enjoyed on the eastern continent, that the people were cut off in consequence of their transgressions, that the last of their prophets who existed among them was commanded to write an abridgement of their prophesies, history \&c., and to hide it up in the earth, and that it should come forth and be united with the bible for the accomplishment of the purposes of God in the last days. For a more particular account I would refer to the Book of Mormon, which can be purchased at Nauvoo, or from any of our travelling elders.

					\hypertarget{EMS-1842-JS-1-MAR-1842-o}%
				As
					\marginnote{o}%
				soon as the news of this discovery was made known, false reports, misrepresentation and slander flew as on the wings of the wind in every direction, the house was frequently beset by mobs, and evil designing persons, several times I was shot at, and very narrowly escaped, and every device was made use of to get the plates away from me, but the power and blessing of God attended me, and several began to believe my testimony.

					\hypertarget{EMS-1842-JS-1-MAR-1842-p}%
				On
					\marginnote{p}%
				the 6th of April, 1830, the ``Church of Jesus Christ of Latter-Day Saints,'' was first organized in the town of Manchester, Ontario co., state of New York. Some few were called and ordained by the spirit of revelation, and prophesy, and began to preach as the spirit gave them utterance, and though weak, yet were they strengthened by the power of God, and many were brought to repentance, were immersed in the water, and were filled with the Holy Ghost by the laying on of hands. They saw visions and prophesied, devils were cast out and the sick healed by the laying on of hands. From that time the work rolled forth with astonishing rapidity, and churches were soon formed in the states of New York, Pennsylvania, Ohio, Indiana, Illinois and Missouri; in the last named state a considerable settlement was formed in Jackson co.;
					\hypertarget{EMS-1842-JS-1-MAR-1842-q}%
				numbers
					\marginnote{q}%
				joined the church and we were increasing rapidly; we made large purchases of land, our farms teemed with plenty, and peace and happiness was enjoyed in our domestic circle and throughout our neighborhood; but as we could not associate with our neighbors who were many of them of the basest of men and had fled from the face of civilized society, to the frontier country to escape the hand of justice, in their midnight revels, their sabbath breaking, horseracing, and gambling, they commenced at first ridicule, then to persecute, and finally an organized mob assembled and burned our houses, tarred, and feathered, and whipped many of our brethren and finally drove them from their habitations;
					\hypertarget{EMS-1842-JS-1-MAR-1842-r}%
				who
					\marginnote{r}%
				houseless, and homeless, contrary to law, justice and humanity, had to wander on the bleak prairies till the children left the tracks of their blood on the prairie, this took place in the month of November, and they had no other covering but the canopy of heaven, in this inclement season of the year; this proceeding was winked at by the government and although we had warrantee deeds for our land, and had violated no law we could obtain no redress.

					\hypertarget{EMS-1842-JS-1-MAR-1842-s}%
				There
					\marginnote{s}%
				were many sick, who were thus inhumanly driven from their houses, and had to endure all this abuse and to seek homes where they could be found. The result was, that a great many of them being deprived of the comforts of life, and the necessary attendances, died; many children were left orphans; wives, widows; and husbands widowers.---Our farms were taken possession of by the mob, many thousands of cattle, sheep, horses, and hogs, were taken and our household goods, store goods, and printing press, and type were broken, taken, or otherwise destroyed.

					\hypertarget{EMS-1842-JS-1-MAR-1842-t}%
				Many
					\marginnote{t}%
				of our brethren removed to Clay where they continued until 1836, three years; there was no violence offered but there were threatnings of violence. But in the summer of 1836, these threatnings began to assume a more serious form; from threats, public meetings were called, resolutions were passed, vengeance and destruction were threatened, and affairs again assumed a fearful attitude, Jackson county was a sufficient precedent, and as the authorities in that county did not interfere, they boasted that they would not in this, which on application to the authorities we found to be too true, and after much violence, privation and loss of property we were again driven from our homes.

					\hypertarget{EMS-1842-JS-1-MAR-1842-u}%
				We
					\marginnote{u}%
				next settled in Caldwell, and Davies[s] counties, where we made large and extensive settlements, thinking to free ourselves from the power of oppression, by settling in new counties, with very few inhabitants in them; but here we were not allowed to live in peace, but in 1838 we were again attacked by mobs an exterminating order was issued by Gov. [Lilburn W.] Boggs, and under the sanction of law an organized banditti ranged through the country, robbed us of our cattle, sheep, horses, hogs \&c., many of our people were murdered in cold blood, the chastity of our women was violated, and we were forced to sign away our property at the point of the sword,
					\hypertarget{EMS-1842-JS-1-MAR-1842-v}%
				and
					\marginnote{v}%
				after enduring every indignity that could be heaped upon us by an inhuman, ungodly band of maurauders, from twelve to fifteen thousand souls men, women, and children were driven from their own fire sides, and from lands that they had warrantee deeds of, houseless, friendless, and homeless (in the depth of winter,) to wander as exiles on the earth or to seek an asylum in a more genial clime, and among a less barbarous people.

					\hypertarget{EMS-1842-JS-1-MAR-1842-w}%
				Many
					\marginnote{w}%
				sickened and died, in consequence of the cold, and hardships they had to endure; many wives were left widows, and children orphans, and destitute. It would take more time than is allotted me here to describe the injustice, the wrongs, the murders, the bloodshed, the theft, misery and woe that has been caused by the barbarous, inhuman, and lawless, proceedings of the state of Missouri.

					\hypertarget{EMS-1842-JS-1-MAR-1842-x}%
				In
					\marginnote{x}%
				the situation before alluded to we arrived in the state of Illinois in 1839, where we found a hospitable people and a friendly home; a people who were willing to be governed by the principles of law and humanity. We have commenced to build a city called ``Nauvoo'' in Hancock co., we number from six to eight thousand here besides vast numbers in the county around and in almost every county of the state. We have a city charter granted us and a charter for a legion the troops of which now number 1500. We have also a charter for a university, for an agricultural and manufacturing society, have our own laws and administrators, and possess all the privileges that other free and enlightened citizens enjoy.

					\hypertarget{EMS-1842-JS-1-MAR-1842-y}%
				Persecution
					\marginnote{y}%
				has not stopped the progress of truth, but has only added fuel to the flame, it has spread with increasing rapidity, proud of the cause which they have espoused and conscious of their innocence and of the truth of their system amidst calumny and reproach have the elders of this church gone forth, and planted the gospel in almost every state in the Union; it has penetrated our cities, it has spread over our villages, and has caused thousands of our intelligent, noble, and patriotic citizens to obey its divine mandates, and be governed by its sacred truths. It has also spread into England, Ireland, Scotland and Wales: in the year of 1839 where a few of our missionaries were sent over five thousand joined the standard of truth, there are numbers now joining in every land.

					\hypertarget{EMS-1842-JS-1-MAR-1842-z}%
				Our
					\marginnote{z}%
				missionaries are going forth to different nations, and in Germany, Palestine, New Holland, the East Indies, and other places, the standard of truth has been erected: no unhallowed hand can stop the work from progressing, persecutions may rage, mobs may combine, armies may assemble, calumny may defame, but the truth of God will go forth boldly, nobly, and independent till it has penetrated every continent, visited every clime, swept every country, and sounded in every ear, till the purposes of God shall be accomplished and the great Jehovah shall say the work is done.

					\hypertarget{EMS-1842-JS-1-MAR-1842-aa}%
				We
					\marginnote{aa}%
				believe in God the Eternal Father, and in his son Jesus Christ, and in the Holy Ghost.

				We believe that men will be punished for their own sins and not for Adam's transgression.

				We believe that through the atonement of Christ all mankind may be saved by obedience to the laws and ordinances of the Gospel.

				We believe that these ordinances are 1st, Faith in the Lord Jesus Christ; 2d, Repentance; 3d, Baptism by immersion for the remission of sins; 4th, Laying on of hands for the gift of the Holy Ghost.

					\hypertarget{EMS-1842-JS-1-MAR-1842-bb}%
				We
					\marginnote{bb}%
				believe that a man must be called of God by ``prophesy, and by laying on of hands'' by those who are in authority to preach the gospel and administer in the ordinances thereof.

				We believe in the same organization that existed in the primitive church, viz: apostles, prophets, pastors, teachers, evangelists \&c.

				We believe in the gift of tongues, prophesy, revelation, visions, healing, interpretation of tongues \&c.

				We believe the bible to be the word of God as far as it is translated correctly; we also believe the Book of Mormon to be the word of God.

				We believe all that God has revealed, all that he does now reveal, and we believe that he will yet reveal many great and important things pertaining to the kingdom of God.

					\hypertarget{EMS-1842-JS-1-MAR-1842-cc}%
				We
					\marginnote{cc}%
				believe in the literal gathering of Israel and in the restoration of the Ten Tribes. That Zion will be built upon this continent. That Christ will reign personally upon the earth, and that the earth will be renewed and receive its paradasaic glory.

				We claim the privilege of worshipping Almighty God according to the dictates of our conscience, and allow all men the same privilege let them worship how, where, or what they may.

				We believe in being subject to kings, presidents, rulers, and magistrates, in obeying, honoring and sustaining the law.

					\hypertarget{EMS-1842-JS-1-MAR-1842-dd}%
				We
					\marginnote{dd}%
				believe in being honest, true, chaste, benevolent, virtuous, and in doing good to \textit{all men}; indeed we may say that we follow the admonition of Paul ``we believe all things we hope all things,'' we have endured many things and hope to be able to endure all things. If there is any thing virtuous, lovely, or of good report or praise worthy we seek after these things. Respectfully \&c.,

				\begin{tightright}
					JOSEPH SMITH.
				\end{tightright}
				
		% -----
		\subsubsection{Notice, circa 1 March 1842}
		% -----
			\label{EMS-1842-JS-CIR-1-MAR-1842}
			% \label{EMS-1842-JS-DAY-3LETTERMONTH-YEAR}
			
			\begin{description}
				\item[Print Sources]
			\end{description}

			\begin{tabular}{p{0.5in}p{1in}p{0.5in}p{1in}}
				JSP	 	& ---		& JSR 		& --- \\
				DC	 	& ---		& HC 		& --- \\
				HDDC 	& ---		& TCHC 		& --- \\
			\end{tabular}
			% HDDC = woodford's DC diss; JSR = Marquardt's JS Revelations; TCHC = Vogel HC
			
			\begin{description}
				\item[Online Access] \href{https://www.josephsmithpapers.org/paper-summary/notice-circa-1-march-1842/1}{Here}
				% \item[Analysis] \hyperref[XX]{Here}
			\end{description}
			
			\begin{description}
				\item[Background]
			\end{description}
			
				% It might also be useful to give a two sentence context of the period: e.g., "This speech was given in Nauvoo to a small congregation one month prior to the destruction of the Nauvoo Expositor. These and other tensions may have been on the prophet's mind when he gave the speech."
			
			\begin{description}
				\item[Original Source]
			\end{description}
			
				% brief description of original source material, with reference to JSP or somewhere else for more info
				% Background material: it might be helpful to have a note that says whether a given document was presented as a speech, was written in a journal, etc.
				% Also a comment here about how reliable the source might be, or what shape it is in, if there are large omissions or parts that are hard to understand, and so on.
			
			\begin{description}
				\item[Text]
			\end{description}
			
				\begin{tightcenter}
						\hypertarget{EMS-1842-JS-CIR-1-MAR-1842-a}%
					TO
						\marginnote{a}%
					THE BRETHREN IN NAUVOO CITY, GREETING:---
				\end{tightcenter}
				
				It is highly important, for the forwarding of the Temple, that an equal distribution of labor should be made, in relation to time; as a superabundance of hands one week, and none the next, tends to retard the progress of the work; therefore, every brother is requested to be particular to labor on the day set apart for the same, in his ward; and to remember that he that sows sparingly shall also reap sparingly,---so that if the brethren want a plentiful harvest, they will do well to be at the place of labor in good season in the morning, bringing all necessary tools, according to their occupation; and those who have teams bring them also, unless otherwise advised by the temple committee.
				
					\hypertarget{EMS-1842-JS-CIR-1-MAR-1842-b}%
				Should
					\marginnote{b}%
				any one be detained from his labor by unavoidable circumstances, on the day appointed, let him labor the next day, or the first day possible.
				
				N. B.---The captains of the respective wards are particularly requested to be at the place of labor on their respective days, and keep an accurate account of each man's work, and be ready to exhibit a list of the same when called for.
				
				The heart of the trustee is daily made to rejoice in the good feelings of the brethren, made manifest in their exertion to carry forward the work of the Lord, and rear his temple; and it is hoped that neither planting, sowing or reaping will hereafter be made to interfere with the regulations hinted at above.
				
				\begin{tightright}
					JOSEPH SMITH.
				
					Trustee in Trust.
				\end{tightright}
				
		% -----
		\subsubsection{Editorial, circa 1 March 1842, Draft}
		% -----
			\label{EMS-1842-JS-CIR-1-MAR-1842-DR}
			% \label{EMS-1842-JS-DAY-3LETTERMONTH-YEAR}
			
			\begin{description}
				\item[Print Sources]
			\end{description}

			\begin{tabular}{p{0.5in}p{1in}p{0.5in}p{1in}}
				JSP	 	& ---		& JSR 		& --- \\
				DC	 	& ---		& HC 		& --- \\
				HDDC 	& ---		& TCHC 		& --- \\
			\end{tabular}
			% HDDC = woodford's DC diss; JSR = Marquardt's JS Revelations; TCHC = Vogel HC
			
			\begin{description}
				\item[Online Access] \href{https://www.josephsmithpapers.org/paper-summary/editorial-circa-1-march-1842-draft/1}{Here}
				% \item[Analysis] \hyperref[XX]{Here}
			\end{description}
			
			\begin{description}
				\item[Background]
			\end{description}
			
				% It might also be useful to give a two sentence context of the period: e.g., "This speech was given in Nauvoo to a small congregation one month prior to the destruction of the Nauvoo Expositor. These and other tensions may have been on the prophet's mind when he gave the speech."
			
			\begin{description}
				\item[Original Source]
			\end{description}
			
				% brief description of original source material, with reference to JSP or somewhere else for more info
				% Background material: it might be helpful to have a note that says whether a given document was presented as a speech, was written in a journal, etc.
				% Also a comment here about how reliable the source might be, or what shape it is in, if there are large omissions or parts that are hard to understand, and so on.
			
			\begin{description}
				\item[Text]
			\end{description}
			
				\begin{tightcenter}
						\hypertarget{EMS-1842-JS-CIR-1-MAR-1842-DR-a}%
					Times
						\marginnote{a}%
					\& Seasons
				\end{tightcenter}

				A. considerable quantity of the matter in the last paper. was in type, before the establishment come into \sout{our} <My> hands,--- Some of which went to press. without \sout{our} <my> recivecd. \sout{or} knowledge \sout{The} and a multiplicity of business--- while enteri[n]g on the additional care of the editorial departmet of the Times \& Seasons. mu[s]t be my apology for what is past.

					\hypertarget{EMS-1842-JS-CIR-1-MAR-1842-DR-b}%
				In
					\marginnote{b}%
				future. I design to furnish much original matter, which will be found of enestimable adventage to the saints,--- \& to all who--- desire a knowledge of the kingdom of God.--- and as it is not practicable to bring forthe the new translation. of the Scriptures. \& varioes records of ancint date. \& great worth to this gen[e]ration in \sout{book} <the usual> form. by books. I shall prenit [print] specimens of the same in the Times \& Seasons as fast. as time \& space will admit. so that the honest in heart may be cheerd \& comforted and go on their way rejoi[ci]ng.--- as their souls become exp[an]ded.--- \& their undestandig [understanding] enlightend, by a knowledg of \sout{what} Gods work through the fathers. in former days, as well as what He is about to do in Latter Days--- To fulfil the words of the fathers.---

					\hypertarget{EMS-1842-JS-CIR-1-MAR-1842-DR-c}%
				In
					\marginnote{c}%
				the penst [present] no. will be found the Commencmet of the Records discoverd in Egypt. some time since. as penend by the hand. of Father Abraham. which I shall contin[u]e to translate \& publish as fast as possible till the whole is completed.--- and as the saints have long been anxious to obtain a copy of these re[c]ords, those are now taking this times \& Seasons. will confer a sp[e]cial favor on their brethren, who do not take the paper, by infor[m]ing them that. they can now obtain their hearts [\textit{3/4 page blank}]
				
		% -----
		\subsubsection{Letter to John C. Bennett, 7 March 1842}
		% -----
			\label{EMS-1842-JS-7-MAR-1842}
			% \label{EMS-1842-JS-DAY-3LETTERMONTH-YEAR}
			
			\begin{description}
				\item[Print Sources]
			\end{description}

			\begin{tabular}{p{0.5in}p{1in}p{0.5in}p{1in}}
				JSP	 	& ---		& JSR 		& --- \\
				DC	 	& ---		& HC 		& --- \\
				HDDC 	& ---		& TCHC 		& --- \\
			\end{tabular}
			% HDDC = woodford's DC diss; JSR = Marquardt's JS Revelations; TCHC = Vogel HC
			
			\begin{description}
				\item[Online Access] \href{https://www.josephsmithpapers.org/paper-summary/letter-to-john-c-bennett-7-march-1842/1}{Here}
				% \item[Analysis] \hyperref[XX]{Here}
			\end{description}
			
			\begin{description}
				\item[Background]
			\end{description}
			
				% It might also be useful to give a two sentence context of the period: e.g., "This speech was given in Nauvoo to a small congregation one month prior to the destruction of the Nauvoo Expositor. These and other tensions may have been on the prophet's mind when he gave the speech."
			
			\begin{description}
				\item[Original Source]
			\end{description}
			
				% brief description of original source material, with reference to JSP or somewhere else for more info
				% Background material: it might be helpful to have a note that says whether a given document was presented as a speech, was written in a journal, etc.
				% Also a comment here about how reliable the source might be, or what shape it is in, if there are large omissions or parts that are hard to understand, and so on.
			
			\begin{description}
				\item[Text]
			\end{description}
			
				\begin{tightright}
					Editor's Office, Nauvoo, Ill.,

					March 7th, 1842.
				\end{tightright}

				\textsc{General} [John C.] \textsc{Bennett};

				\textit{Respected Brother}:---I have just been perusing your correspondence with Doctor [Charles V.] Dyer on the subject of American Slavery, and the students of the Quincy Mission Institute, and it makes my blood boil within me to reflect upon the injustice, cruelty, and oppression, of the rulers of the people---when will these things cease to be, and the Constitution and the Laws again bear rule? I fear for my beloved country---mob violence, injustice, and cruelty, appear to be the darling attributes of Missouri, and no man taketh it to heart! \textit{O, tempora! O, mores!} What think you should be done?

				\begin{tightright}
					Your friend,

					JOSEPH SMITH.
				\end{tightright}
				
		% -----
		\subsubsection{Letter to Edward Hunter, 9 and 11 March 1842}
		% -----
			\label{EMS-1842-JS-9-11-MAR-1842}
			% \label{EMS-1842-JS-DAY-3LETTERMONTH-YEAR}
			
			\begin{description}
				\item[Print Sources]
			\end{description}

			\begin{tabular}{p{0.5in}p{1in}p{0.5in}p{1in}}
				JSP	 	& ---		& JSR 		& --- \\
				DC	 	& ---		& HC 		& --- \\
				HDDC 	& ---		& TCHC 		& --- \\
			\end{tabular}
			% HDDC = woodford's DC diss; JSR = Marquardt's JS Revelations; TCHC = Vogel HC
			
			\begin{description}
				\item[Online Access] \href{https://www.josephsmithpapers.org/paper-summary/letter-to-edward-hunter-9-and-11-march-1842/1}{Here}
				% \item[Analysis] \hyperref[XX]{Here}
			\end{description}
			
			\begin{description}
				\item[Background]
			\end{description}
			
				% It might also be useful to give a two sentence context of the period: e.g., "This speech was given in Nauvoo to a small congregation one month prior to the destruction of the Nauvoo Expositor. These and other tensions may have been on the prophet's mind when he gave the speech."
			
			\begin{description}
				\item[Original Source]
			\end{description}
			
				% brief description of original source material, with reference to JSP or somewhere else for more info
				% Background material: it might be helpful to have a note that says whether a given document was presented as a speech, was written in a journal, etc.
				% Also a comment here about how reliable the source might be, or what shape it is in, if there are large omissions or parts that are hard to understand, and so on.
			
			\begin{description}
				\item[Text]
			\end{description}
			
				...I am now very busily engaged in Translating, and therefore cannot give as much time to Public matters as I could wish, but will nevertheless do what I Can to forward your affairs...
				
		% -----
		\subsubsection{Minutes, 15--16 March 1842}
		% -----
			\label{EMS-1842-JS-15-16-MAR-1842}
			% \label{EMS-1842-JS-DAY-3LETTERMONTH-YEAR}
			
			\begin{description}
				\item[Print Sources]
			\end{description}

			\begin{tabular}{p{0.5in}p{1in}p{0.5in}p{1in}}
				JSP	 	& ---		& JSR 		& --- \\
				DC	 	& ---		& HC 		& --- \\
				HDDC 	& ---		& TCHC 		& --- \\
			\end{tabular}
			% HDDC = woodford's DC diss; JSR = Marquardt's JS Revelations; TCHC = Vogel HC
			
			\begin{description}
				\item[Online Access] \href{https://www.josephsmithpapers.org/paper-summary/minutes-15-16-march-1842/1}{Here}
				% \item[Analysis] \hyperref[XX]{Here}
			\end{description}
			
			\begin{description}
				\item[Background]
			\end{description}
			
				% It might also be useful to give a two sentence context of the period: e.g., "This speech was given in Nauvoo to a small congregation one month prior to the destruction of the Nauvoo Expositor. These and other tensions may have been on the prophet's mind when he gave the speech."
			
			\begin{description}
				\item[Original Source]
			\end{description}
			
				% brief description of original source material, with reference to JSP or somewhere else for more info
				% Background material: it might be helpful to have a note that says whether a given document was presented as a speech, was written in a journal, etc.
				% Also a comment here about how reliable the source might be, or what shape it is in, if there are large omissions or parts that are hard to understand, and so on.
			
			\begin{description}
				\item[Text]
			\end{description}
			
				\begin{tightcenter}
						\hypertarget{EMS-1842-JS-15-16-MAR-1842-a}%
					Tuesday,
						\marginnote{a}%
					March 15th, A. L. 5842,%
						\footnote{%
							% \label{}%
							JSP note: ```5842' refers to the year in the Masonic dating system, representing the Gregorian year plus 4,000. Dates in the Masonic system are often preceded by \textit{Anno Lucis}, Latin for `year of light.' (`Masonic Computation of Time,' 129--131.).'' The source is : ```Masonic Computation of Time.' \textit{Freemasons' Monthly Magazine} 9, no. 5 (1 Mar. 1850): 129--131.''
							}
					AD. 1842,

					9 o'clock, A. M.
				\end{tightcenter}

				Special Communication, by order of Grand Master [Abraham] Jonas. Lodge met pursuant to previous notice. Present--- W. M. [Worshipful Master] A. Jonas, G. M. of the G. L. [Grand Lodge] of Illinois; W. D. McCann, D. G. S. [Deputy Grand Secretary]; Asahel Perry, W. M. p. t. [pro tem]; Hyrum Smith, S. W. [Senior Warden]; L[ucius] N. Scovil, J. W. [Junior Warden]; N[ewel] K. Whitney, T[reasurer]; John C. Bennett, S[ecretary]; Charles Allen, S. D. [Senior Deacon]; Heber C. Kimball, J. D. [Junior Deacon]; S[amuel] Rolfe, T[yler]; Wm. Felshaw, and Hiram Clark, S[tewards]; Hezekiah Peck, Josiah Arnold, Joshua Smith, Lyman Leonard, David Pettegrew, Christopher Williams, John Patten, Noah Rogers, Elijah Fordham, Samuel Henderson, Benjamin Brown, Stephen Chase, George Montague, Noble Rogers, and Daniel S. Miles, Members
					\hypertarget{EMS-1842-JS-15-16-MAR-1842-b}%
				and
					\marginnote{b}%
				L. B. Adams, Franklin, 22, Ill.; M. Plumb, Franklin, 22, Ill.; Henry King, Sylvan, 229, N.Y.; A. C. Graves, Harmony, 11, Ia.; J. Rose, St. Johns, 21, N.Y.; A. Lambert, Sciota, 28, O[hio]; G. Heberling, La Fayette, Pa.; D. Hibbard, New England, , O.; James Cummings, Maine, Me.; E. Jennings, Mt. Moriah, N.Y.; Samuel Miles, Rainbow, Vt.; Thomas C. King, Bodley, 1, Ill.; Caleb Baldwin, Concord, O.; S. Comer, Urbanna, O.; \& L. Evans, Urbanna, O.; visiting brethren. The M. W. [Most Worshipful] Grand Master presided, and opened the lodge in due form, in the 3rd degree of masonry.%
					\footnote{%
						% \label{}%
						Here the JSP gives a note to this source: ``Reynolds, John C. \textit{History of the M. W. Grand Lodge of Illinois, Ancient, Free, and Accepted Masons, From the Organization of the First Lodge Within the Present Limits of the State, Up to and Including 1850}. Springfield, IL: H. G. Reynolds, 1869.''
						}
				He, likewise, organised, and opened a Grand Lodge according to ancient usage, in the 3rd degree. The Grand and Subordinate Lodges then called off from labor to refreshment, until 1 o'clock, P. M.

				\begin{tightcenter}
						\hypertarget{EMS-1842-JS-15-16-MAR-1842-c}%
					1
						\marginnote{c}%
					o'clock, P. M.
				\end{tightcenter}

				The Lodges called to labor. Present as before. The Grand Master then directed the Grand Marshal, John C. Bennett, to form a procession according to ancient usage, in order to proceed to the grove, near the Temple, for installation, which was accordingly done. At the grove, after the ceremonies of installation, the Grand Master delivered a highly creditable and finished address on the subject of Ancient York Masonry, after which the lodges returned to the lodge room in masonic orders. The Grand Lodge then closed in due form, without day. The subordinate Lodge then adopted the by-laws which they had informally prepared on the 30th of December, as amended at the suggestion of the Grand Master. On motion, the Lodges requested a copy of the address of the Grand Master for publication--- carried unanimously, and a copy promised. The Lodge then closed to meet again at early candle lighting.

				\begin{tightcenter}
						\hypertarget{EMS-1842-JS-15-16-MAR-1842-d}%
					7
						\marginnote{d}%
					o'clock, P. M.
				\end{tightcenter}

				Lodge met pursuant to adjournment, and opened in the 1st degree of masonry. Present as before. A special dispensation from the Grand Master was then presented in the words following; to wit:

				\begin{quote}
					\begin{tightright}
							\footnote{%
								% \label{}%
								Beginning here and proceeding to Jonas's signature a few pararaphs below, this is a copy of an authorization created by Abraham Jonas, to begin the Nauvoo Lodge. See the original copy of the authorization \href{https://www.josephsmithpapers.org/paper-summary/authorization-from-abraham-jonas-15-march-1842/1}{here}.
								}%
						City of Nauvoo, March 15--- 5842.
					\end{tightright}

					Know all ye brethren to whom come these presents--- That I, Abraham Jonas, Grand Master of the Grand Lodge of the State of Illinois--- in virtue of the power and authority in me vested as Grand Master aforesaid do hereby, by these letters of Dispensation, authorize the Brethren of Nauvoo Lodge under dispensation, to receive the petitions of Joseph Smith, and Sidney Rigdon--- and act on the same \uline{instanter}--- and should the ballot be unanimous in favor of said Smith and Rigdon at a full meeting of said Nauvoo Lodge--- then in that case---
						\hypertarget{EMS-1842-JS-15-16-MAR-1842-e}%
					the
						\marginnote{e}%
					said Lodge is authorized to confer the three several degrees of Ancient York Masonry on the said Joseph Smith and Sidney Rigdon--- as speedily as the nature of the case will admit--- Provided, however, that nothing herein contained shall be deemed as authority by the said Nauvoo Lodge for violating any of the ancient landmarks of the order, or of acting contrary to the provisions of their By-Laws--- except in the cases herein authorized---

					Given under my hand--- as Grand Master, the day and date above named---

					A. Jonas, G. M. G. L. Ill.
				\end{quote}

				Petitions were then presented from Joseph Smith, and Sidney Rigdon, the persons above named, and on motion the Lodge proceeded to the ballot in each case separately, in order of record--- the ballot was found clear, and they were duly initiated--- Entered Apprentice Masons. The Lodge was then closed, to stand closed until to-morrow morning at 9 o'clock.

				John C. Bennett, Secretary.

				\begin{tightcenter}
						\hypertarget{EMS-1842-JS-15-16-MAR-1842-f}%
					Wednesday,
						\marginnote{f}%
					March 16th, A.L. 5842, AD. 1842.

					9 o'clock, A. M.
				\end{tightcenter}

				Lodge met pursuant to adjournment. Present--- W. M. A. Jonas, G. M. presiding; W. D. McCann, D. G. S.; H. Smith, S. W.; L. N. Scovil, J. W.; J. C. Bennett, S.; N. K. Whitney, T.; C. Allen, S. D.; H. C. Kimball, J. D.; Lyman Leonard, T. p. t.; W. Felshaw, \& H. Clark, S.; H. Peck, J. Arnold, Joshua Smith, A[ustin] Cowles, D. Pettegrew, S. Rolfe, J. Patten, Noah Rogers, E. Fordham, S. Henderson, S. Chase, Noble Rogers, Ormond Butler, Members; J. Smith \& S. Rigdon, E. A. Masons,--- and L. B. Adams, M. Plum, H. King, A. C. Graves, J. Rose, A. Lambert, D. Hibbard, \& S. Miles, visiting brethren. An E. A. Lodge was then opened in due form, and the minutes of the preceding meeting meeting read.
					\hypertarget{EMS-1842-JS-15-16-MAR-1842-g}%
				No
					\marginnote{g}%
				further business appearing in this degree, the E. A. were requested to withdraw, and a F. C. [Fellow Craft] Lodge was then opened. Joseph Smith, and Sidney Rigdon, signified their desire to be advanced in masonry, and their proficiency in the preceding degree was vouched for--- after which they were balloted for, the ballot found clear, and they duly received and passed to the degree of F. C. Masons. The Lodge then called off until 2 o'clock, P.M.

				\begin{tightcenter}
						\hypertarget{EMS-1842-JS-15-16-MAR-1842-h}%
					2
						\marginnote{h}%
					o'clock, P.M.
				\end{tightcenter}

				Lodge called to labor. Present as before. A Master's Lodge was then opened and Joseph Smith applied for the third and sublime degree of M. M. [Master Mason]--- the ballot was found clear, and he was duly raised to said degree, and signed the By-Laws. Lodge called off until 7 o'clock, P.M.

				\begin{tightcenter}
						\hypertarget{EMS-1842-JS-15-16-MAR-1842-i}%
					7
						\marginnote{i}%
					o'clock, P.M.
				\end{tightcenter}

				Lodge called to labor. Present as before. Sidney Rigdon applied for the sublime degree of Master Mason--- the ballot was found clear, and he was accordingly raised. The S. W. was then called to the chair, and the following resolutions unanimously adopted; to wit:

				Resolved, That the thanks of this Lodge be presented to G. M. Jonas, \& D. G. S. McCann, for their attendance with us, and the instructions given in the various degrees of masonry.

					\hypertarget{EMS-1842-JS-15-16-MAR-1842-j}%
				Resolved,
					\marginnote{j}%
				That they be requested to visit the Lodge whenever their circumstances will permit, at the expense of the Lodge.

				Resolved, That G. M. Jonas be requested to deliver an oration on masonry, in this city, on the 24th of June next; but, if not possible so to do, to give the lodge timely notice.

					\hypertarget{EMS-1842-JS-15-16-MAR-1842-k}%
				Resolved,
					\marginnote{k}%
				That Columbus Lodge, No. 6., \& Bodley Lodge, No. 1., be requested to meet with Nauvoo Lodge on the 24th of June--- and that all brethren in the vicinity be solicited to participate.

				G. M. Jonas, \& D. G. S. McCann, then made very neat speeches, expressive of their entire satisfaction with the proceedings of this lodge, and the kind treatment they had received at the hands of our citizens generally. The lodge then closed, in great cordiality of feelings, to stand closed until to-morrow evening.

				John C. Bennett, Secretary.
				
		% -----
		\subsubsection{Minutes and Discourses, 17 March 1842}
		% -----
			\label{EMS-1842-JS-17-MAR-1842}
			% \label{EMS-1842-JS-DAY-3LETTERMONTH-YEAR}
			
			\begin{description}
				\item[Print Sources]
			\end{description}

			\begin{tabular}{p{0.5in}p{1in}p{0.5in}p{1in}}
				JSP	 	& ---		& JSR 		& --- \\
				DC	 	& ---		& HC 		& --- \\
				HDDC 	& ---		& TCHC 		& --- \\
			\end{tabular}
			% HDDC = woodford's DC diss; JSR = Marquardt's JS Revelations; TCHC = Vogel HC
			
			\begin{description}
				\item[Online Access] \href{https://www.josephsmithpapers.org/paper-summary/minutes-and-discourses-17-march-1842/1}{Here}
				% \item[Analysis] \hyperref[XX]{Here}
			\end{description}
			
			\begin{description}
				\item[Background]
			\end{description}
			
				% It might also be useful to give a two sentence context of the period: e.g., "This speech was given in Nauvoo to a small congregation one month prior to the destruction of the Nauvoo Expositor. These and other tensions may have been on the prophet's mind when he gave the speech."
			
			\begin{description}
				\item[Original Source]
			\end{description}
			
				% brief description of original source material, with reference to JSP or somewhere else for more info
				% Background material: it might be helpful to have a note that says whether a given document was presented as a speech, was written in a journal, etc.
				% Also a comment here about how reliable the source might be, or what shape it is in, if there are large omissions or parts that are hard to understand, and so on.
			
			\begin{description}
				\item[Text]
			\end{description}
			
				\begin{tightright}
						\hypertarget{EMS-1842-JS-17-MAR-1842-a}%
					Nauvoo
						\marginnote{a}%
					Lodge Room

					March 17\supers{th} 1842.
				\end{tightright}

				Present--- President Joseph Smith, John Taylor, Willard Richards, Emma Smith and others.

				Elder John Taylor was call'd to the chair by Pres\supers{t.} Smith, and elder W. Richards appointed Secretary,

				Meeting commenced by singing ``The spirit of God like a fire is burning'' \&c.--- Prayer by elder Taylor.

					\hypertarget{EMS-1842-JS-17-MAR-1842-b}%
				When
					\marginnote{b}%
				it was mov'd by Pres\supers{t.} Smith and seconded by Mrs. [Sarah Kingsley] Cleveland, that a vote be taken to know if all are satisfied with each female present; and are willing to acknowledge them in full fellowship, and admit them to the privileges of the Institution about to be formed.

				The names of those present were then taken as follows

				\begin{tightcenter}
					Mrs Emma Smith
				\end{tightcenter}
				
				\begin{tabular}{p{4in}p{3in}}
				Mrs. Sarah M. Cleveland					&	Bathsheba W. [Bigler] Smith		\\
				Phebe Ann [Baldwin] Hawkes				&	Phebe M. [Bartholomew] Wheeler  \\
				Elizabeth [Hughes] Jones				&	Elvira A. Co[w]les              \\
				Sophia [Bundy] Packard					&	Margaret A [Norris] Cook        \\
				Philinda [Eldredge] Merrick				&	Athalia [Rigdon] Robinson       \\
				Martha Knights [Martha McBride Knight]	&	Sarah M. [Granger] Kimball      \\
				Desdemona Fu[l]lmer						&	Eliza R. Snow                   \\
				Elizabeth Ann [Smith] Whitney			&	Sophia Robinson                 \\
				Leonora [Cannon] Taylor					&	Nancy Rigdon                    \\
				\end{tabular}

				\begin{tightcenter}
					Sophia R. Marks
				\end{tightcenter}

					\hypertarget{EMS-1842-JS-17-MAR-1842-c}%
				Pres\supers{t.}
					\marginnote{c}%
				Smith, \& Elders Taylor and Richa[r]ds withdrew while the females went into an investigation of the motion, and decided that all present, be admitted according to the motion, and that

				Mrs. Sarah [Ward] Higbee

				Thirza [Stiles] Cahoon

				Kezia A. [Voorhees] Morrison

				Miranda N. [Marinda Nancy Johnson] Hyde

				Abigail [McMurtrey] Allred

				Mary [Heron] Snider

				Sarah [Stiles] Granger [\textit{4 lines blank}]

				should be admitted; whose names were presented by Pres\supers{t.} Smith.

					\hypertarget{EMS-1842-JS-17-MAR-1842-d}%
				Pres\supers{t.} Smith,
					\marginnote{d}%
				\& Elders Taylor \& Richards return'd and the meeting was address'd by Prest. Smith, to illustrate the object of the Society--- that the Society of Sisters might provoke the brethren to good works in looking to the wants of the poor--- searching after objects of charity, and in administering to their wants--- to assist; by correcting the morals and strengthening the virtues of the female community, and save the Elders the trouble of rebuking; that they may give their time to other duties \&c. in their public teaching.

					\hypertarget{EMS-1842-JS-17-MAR-1842-e}%
				Pres\supers{t.}
					\marginnote{e}%
				Smith further remark'd that an organization to show them how to go to work would be sufficient. He propos'd that the Sisters elect a presiding officer to preside over them, and let that presiding officer choose two Counsellors to assist in the duties of her Office--- that he would ordain them to preside over the Society--- and let them preside just as the Presidency, preside over the church; and if they need his instruction--- ask him, he will give it from time to time.

					\hypertarget{EMS-1842-JS-17-MAR-1842-f}%
				Let
					\marginnote{f}%
				this Presidency serve as a constitution--- all their decisions be considered law; and acted upon as such.

				If any Officers are wanted to carry out the designs of the Institution, let them be appointed and set apart, as Deacons, Teachers \&c. are among us.

				The minutes of your meetings will be precedents for you to act upon--- your Constitutio[n] and law.

				He then suggested the propriety of electing a Presidency to continue in office during good behavior, or so long as they shall continue to fill the office with dignity \&c. like the first Presidency of the church.---

					\hypertarget{EMS-1842-JS-17-MAR-1842-g}%
				Motioned
					\marginnote{g}%
				by Sister Whitney and seconded by Sister Packard that Mrs. Emma Smith be chosen President--- passed unanimously---

				Mov'd by Pres\supers{t.} Smith, that Mrs. Smith proceed to choose her Counsellors, that they may be ordain'd to preside over this Society, in taking care of the poor--- administering to their wants, and attending to the various affairs of this Institution.

				The Presidentess Elect, then made choice of Mrs. Sarah M. Cleveland and Mrs. Elizabeth Ann Whitney for Counsellors---

					\hypertarget{EMS-1842-JS-17-MAR-1842-h}%
				President
					\marginnote{h}%
				Smith read the Revelation to Emma Smith, from the book of Doctrine and Covenants; and stated that she was ordain'd at the time, the Revelation was given, to expound the scriptures to all; and to teach the female part of community; and that not she alone, but others, may attain to the same blessings.---

				The 2\supers{d} Epistle of John, 1\supers{st} verse, was then read to show that respect was then had to the same thing; and that why she was called an Elect lady is because, elected to preside.

					\hypertarget{EMS-1842-JS-17-MAR-1842-i}%
				Elder
					\marginnote{i}%
				Taylor was then appointed to ordain the Counsellors--- he laid his hands on the head of Mrs Cleveland and ordain'd her to be a Counsellor to the Elect Lady, even Mrs. Emma Smith, to counsel, and assist her in all things pertaining to her office \&c.

				Elder T. then laid his hands on the head of Mrs. Whitney and ordain'd her to be a Counsellor to Mrs. Smith, the Prest. of the Institutio[n]--- with all the privileges pertaining to the office \&c.

					\hypertarget{EMS-1842-JS-17-MAR-1842-j}%
				He
					\marginnote{j}%
				hen laid his hands on the head of Mrs. Smith and blessed her, and confirm'd upon her all the blessings which have been confer'd on her, that she might be a mother in Israel and look to the wants of the needy, and be a pattern of virtue; and possess all the qualifications necessary for her to stand and preside and dignify her Office, to teach the females those principles requisite for their future usefulness.

					\hypertarget{EMS-1842-JS-17-MAR-1842-k}%
				Pres\supers{t.}
					\marginnote{k}%
				Smith then resumed his remarks and gave instruction how to govern themselves in their meetings--- when one wishes to speak, address the chair--- and the chairman responds to the address.

				Should two speak at once, the Chair shall decide who speaks first, if any one is dissatisfied, she appeals to the house---

				When one has the floor, occupies as long as she pleases.

				Proper manner of address is Mrs. Chairman or President and not Mr. Chairman \&c.

				A question can never be put until it has a second

				When the subject for discussion has been fairly investigated; the Chairman will say, are you ready for the question? \&c.

				Whatever the majority of the house decide upon becomes a law to the Society.

					\hypertarget{EMS-1842-JS-17-MAR-1842-l}%
				Pres\supers{t.}
					\marginnote{l}%
				Smith proceeded to give counsel--- do not injure the character of any one--- if members of the Society shall conduct improperly, deal with them, and keep all your doings within your own bosoms, and hold all characters sacred---

				It was then propos'd that Elder Taylor vacate the chair.

				Pres\supers{t.} Emma Smith and her Counsellors took the chair, and

				Elder Taylor mov'd--- sec\supers{d} by Pres\supers{t.} J. Smith that we go into an investigation respecting what this Society shall be call'd--- which was

				carried unanimously

					\hypertarget{EMS-1842-JS-17-MAR-1842-m}%
				Pres\supers{t.}
					\marginnote{m}%
				Smith continued instructions to the Chair to suggest to the members anything the chair might wish, and which it might not be proper for the chair to put, or move \&c.

				Mov'd by Counsellor Cleveland, and secon'd by Counsellor Whitney, that this Society be called \uline{The} \uline{Nauvoo} \uline{Female} \uline{Relief} Society.

				Elder Taylor offered an amendment, that it be called \uline{The} \uline{Nauvoo} \uline{Female} \uline{Benevolent} \uline{Society} which would give a more definite and extended idea of the Institution--- that Relief be struck out and Benevolent inserted.

					\hypertarget{EMS-1842-JS-17-MAR-1842-n}%
				Pres\supers{t.}
					\marginnote{n}%
				Smith offer'd instruction on votes---

				The motion was seconded by Counsellor Cleveland and unanimously carried, on the amendment by Elder Taylor.

				The Pres\supers{t.} then suggested that she would like an argument with Elder Taylor on the words Relief and Benevolence.

				Pres\supers{t.} J. Smith mov'd that the vote for amendment, be rescinded, which was carried---

				Motion for adjournment by Elder Richards and objected by Pres\supers{t.} J. Smith.---

					\hypertarget{EMS-1842-JS-17-MAR-1842-o}%
				Pres\supers{t.}
					\marginnote{o}%
				J. Smith--- Benevolent is a popular term--- and the term \uline{Relief} is not known among popular Societies--- \uline{Relief} is more extended in its signification than \uline{Benevolent} and might extend to the liberation of the culprit--- and might be wrongly construed by our enemies to say that the Society was to relieve criminals from punishment \&c. \&c--- to relieve a murderer, which would not be a benevolent act---

					\hypertarget{EMS-1842-JS-17-MAR-1842-p}%
				Pres\supers{t.}
					\marginnote{p}%
				Emma Smith, said the \uline{popularity} of the word benevolent is one great objection--- no person can think of the word as associated with public Institutions, without thinking of the Washingtonian Benevolent Society which was one of the most corrupt Institutions of the day--- do not wish to have it call'd after other Societies in the world---

				Pres\supers{t.} J. Smith arose to state that he had no objection to the word \uline{Relief}--- that on question they ought to deliberate candidly and investigate all subjects.

					\hypertarget{EMS-1842-JS-17-MAR-1842-q}%
				Counsellor
					\marginnote{q}%
				Cleveland arose to remark concerning the question before the house, that we should not regard the idle speech of our enemies--- we design to act in the name of the Lord--- to relieve the wants of the distressed, and do all the good we can.---

					\hypertarget{EMS-1842-JS-17-MAR-1842-r}%
				Eliza
					\marginnote{r}%
				R. Snow arose and said that she felt to concur with the President, with regard to the word Benevolent, that many Societies with which it had been associated, were corrupt,--- that the popular Institutions of the day should not be our guide--- that as daughters of Zion, we should set an example for all the world, rather than confine ourselves to the course which had been heretofore pursued--- one objection to the word \uline{Relief} is, that the idea associated with it is that of some great calamity--- that we intend appropriating on some extraordinary occasions instead of meeting the common occurrences---

					\hypertarget{EMS-1842-JS-17-MAR-1842-s}%
				Pres\supers{t.}
					\marginnote{s}%
				Emma Smith remark'd--- we are going to do something \uline{extraordinary}--- when a boat is stuck on the rapids with a multitude of Mormons on board we shall consider \uline{that} a loud call for \uline{relief}--- we expect extraordinary occasions and pressing calls---

				Elder Taylor arose and said--- I shall have to concede the point--- your arguments are so potent I cannot stand before them--- I shall have to give way---

				Pres\supers{t.} J. S. said I also shall have to concede the point, all I shall have to give to the poor, I shall give to this Society---

					\hypertarget{EMS-1842-JS-17-MAR-1842-t}%
				Counsellor
					\marginnote{t}
				Whitney mov'd, that this Society be call'd \uline{The} \uline{Nauvoo} \uline{Female} \uline{Relief} \uline{Society}--- secon\supers{d.} by Counsellor Cleveland---

				E. R. Snow offer'd an amendment by way of transposition of words, instead of The \uline{Nauvoo} \uline{Female} \uline{Relief} \uline{Society}, it shall be call'd The \uline{Female} \uline{Relief} \uline{Society} of Nauvoo--- Seconded by Pres\supers{t.} J. Smith and carried---

				The previous question was then put--- Shall this Society be call'd \uline{The} \uline{Female} \uline{Relief} \uline{Society} of \uline{Nauvoo}?--- carried unanimously.---

					\hypertarget{EMS-1842-JS-17-MAR-1842-u}%
				Pres\supers{t.}
					\marginnote{u}%
				J. Smith--- I now declare this Society organiz'd with President and Counsellors \&c. according to Parliamentary usages--- and all who shall hereafter be admitted into this Society must be free from censure and receiv'd by vote---

				\begin{tabular}{p{6in}p{1in}}
				Pres\supers{t.} J. Smith offered & \$5.00
				\end{tabular}

				in gold piece to commence the funds of the Institution.

				Pres\supers{t.} Emma Smith requested that the gentlemen withdraw before they proceed to the choice of Secretary and Treasurer, as was mov'd by Pres\supers{t.} J. Smith---

				\begin{tightright}
					\uline{Willard} \uline{Richards.} \uline{Sec}\supersu{ty}\supers{.}
				\end{tightright}

					\hypertarget{EMS-1842-JS-17-MAR-1842-v}%
				The
					\marginnote{v}%
				gentlemen withdrew when it was Motioned and secon\supers{d.} and unanimously pass'd that Eliza R. Snow be appointed Secretary, and Phebe M. Wheeler, Assistant Secretary------

				Motioned, secon\supers{d.} and carried unanim\supers{ly.} that Elvira A. Coles be appointed Treasurer---

				Pres\supers{t.} E. Smith then arose and proceeded to make appropriate remarks on the object of the Society--- its duties to others also its relative duties to each other Viz. to seek out and relieve the distressed--- that each member should be ambitious to do good--- that the members should deal frankly with each other--- to watch over the morals--- and be very careful <of> the character and reputation--- of the members of the Institution \&c.

					\hypertarget{EMS-1842-JS-17-MAR-1842-w}%
				P.
					\marginnote{w}%
				A. Hawkes--- Question--- What shall we reply to interrogatories relative to the object of this Society?

				Pres\supers{t.} E. Smith replied--- for charitable purposes.

				Mov'd and pass'd that Cynthia Ann [Howlett] Eldridge be admitted as a member of this Society---

				\begin{tabular}{p{6in}p{1in}}
				Cou\supers{lr.} Sarah M. Cleveland donated to the fund of the Society & \$12$\frac{1}{2}$ \\
				Sarah M. Kimball d[itt]o & 1.00 \\
				Pres\supers{t.} Emma Smith do & 1.00 \\
				Coun\supers{lr.} E. A. Whitney do & 50 \\
				\end{tabular}

					\hypertarget{EMS-1842-JS-17-MAR-1842-x}%
				Pres\supers{t.}
					\marginnote{x}%
				E. Smith said that Mrs. Merrick is a widow--- is industrious--- performs her work well, therefore recommend her to the patronage of such as wish to hire needlework--- those who hire widows must be prompt to pay and inasmuch as some have defrauded the laboring widow of her wages, we must be upright and deal justly---

				The business of the Society concluded--- the gentlemen before mentioned return'd---,

				\begin{tabular}{p{6in}p{1in}}
				Elder Richards appropriated to the fund of the Society, the sum of & \$1,00 \\
				Elder Taylor do	& 2.00 \\
				\end{tabular}

					\hypertarget{EMS-1842-JS-17-MAR-1842-y}%
				Elder
					\marginnote{y}%
				T. then arose and address'd the Society by saying that he is much gratified in seeing a meeting of this kind in Nauvoo--- his heart rejoices when he sees the most distinguished characters, stepping forth in such a cause, which is calculated to bring into exercise every virtue and give scope to the benevolent feelings of the female heart--- he rejoices to see this Institution organiz'd according to the law of Heaven--- according to a revelation previously given to Mrs E. Smith appointing her to this important calling--- and to see all things moving forward in such a glorious manner--- his prayer is that the blessings of God and the peace of heaven may rest on this Institution henceforth------

					\hypertarget{EMS-1842-JS-17-MAR-1842-z}%
				The
					\marginnote{z}%
				Choir then sang ``Come let us rejoice in the day of salvation[''] \&c.

				Motion'd, that this meeting adjourn to next week, thursday, ten o'clock--- A M.

				The meeting then arose and was dismiss'd by prayer by Elder Taylor.---
				
		% -----
		\subsubsection{Discourse, 20 March 1842}
		% -----
			\label{EMS-1842-JS-20-MAR-1842}
			% \label{EMS-1842-JS-DAY-3LETTERMONTH-YEAR}
			
			% --
			\paragraph{As Reported by Wilford Woodruff}
			% --
				\label{EMS-1842-JS-20-MAR-1842-WW}
			
				\begin{description}
					\item[Print Sources]
				\end{description}

				\begin{tabular}{p{0.5in}p{1in}p{0.5in}p{1in}}
					JSP	 	& ---		& JSR 		& --- \\
					DC	 	& ---		& HC 		& --- \\
					HDDC 	& ---		& TCHC 		& --- \\
				\end{tabular}
				% HDDC = woodford's DC diss; JSR = Marquardt's JS Revelations; TCHC = Vogel HC
				
				\begin{description}
					\item[Online Access] \href{https://www.josephsmithpapers.org/paper-summary/discourse-20-march-1842-as-reported-by-wilford-woodruff/1}{Here}
					% \item[Analysis] \hyperref[XX]{Here}
				\end{description}
				
				\begin{description}
					\item[Background]
				\end{description}
				
					% It might also be useful to give a two sentence context of the period: e.g., "This speech was given in Nauvoo to a small congregation one month prior to the destruction of the Nauvoo Expositor. These and other tensions may have been on the prophet's mind when he gave the speech."
				
				\begin{description}
					\item[Original Source]
				\end{description}
				
					% brief description of original source material, with reference to JSP or somewhere else for more info
					% Background material: it might be helpful to have a note that says whether a given document was presented as a speech, was written in a journal, etc.
					% Also a comment here about how reliable the source might be, or what shape it is in, if there are large omissions or parts that are hard to understand, and so on.
				
				\begin{description}
					\item[Text]
				\end{description}
				
						\hypertarget{EMS-1842-JS-20-MAR-1842-WW-a}%
					The
						\marginnote{a}%
					Speaker read the 14 \uline{ch} Revelations. And sayes ``we have again the warning voice sounded in our midst which shows the uncertainty of human life. And in my leasure moments I have meditated upon the subject, \& asked the question Why is it that, infant innocent children are taken away from us, esspecially those that seem to be most intelligent beings'' Answer ``This world is a vary wicked world \& it is a proverb that the world grow weaker \& wiser, but if it is the case the world grows more wicke[d] \& corrupt.
						\hypertarget{EMS-1842-JS-20-MAR-1842-WW-b}%
					In
						\marginnote{b}%
					the early ages of the world A richeoos [righteous] man \& a man of God \& intell[i]gence had a better chance to do good to be received \& believed than at the present day, but in \sout{this} these days such a man is much opposed \& persecuted by most of the inhabitants of the earth \& he has much sorrow to pass through, hence the Lord takes many away even in infancy that they may escape the envy of man, the sorrows \& evils of this present world \& they were two pure \& to[o] lov[e]ly to live on Earth,
						\hypertarget{EMS-1842-JS-20-MAR-1842-WW-c}%
					Therefore
						\marginnote{c}%
					if rightly considerd \sout{we have}, instead of mo[u]rning we have reason to rejoice, as they are deliverd from evil \& we shall soon have them again, What chanc[e] is their for infidelity when we are parting with our friends almost daily none at all The infidel will grasp at evry straw for help untill death stares him in the face \& then his infidelity takes its flight for the realities of the eternal world are resting upon him in mighty power \& when evry earthly support \& prop fails him, he then sensibly feels the eternal truths of the immortality of the soul.
						\hypertarget{EMS-1842-JS-20-MAR-1842-WW-d}%
					Also
						\marginnote{d}%
					the doctrin of Baptizing Children or sprinkling them or they must welter in Hell is a doctrin not true not supported in Holy writ \& is not consistant with the character of god The moment that children leave this world they are taken to the bosom of Abraham The ownly difference between the old \& young dying is one lives longer in heaven \& Eternal light \& glory than the other \& was freed a little sooner from this miserable wicked world Notwithstanding all this glory we for a moment loose sight of it \& mourn the loss but we do not mourn as those without hope.
						\hypertarget{EMS-1842-JS-20-MAR-1842-WW-e}%
					(We
						\marginnote{e}%
					should take warning \& not wait for the deathbed to repent) As we se[e] the infant taken away by death, so may the youth \& middle as well as the infant suddenly be called into <a great> eternity, Let this then proove as a warning to all not to procrastinate repentance or wait till a death bed, for it is the will of God that man should repent \& serve him in health \& in the strength \& power of his mind in order to secure his blessings \& not wait untill he is called to die.

						\hypertarget{EMS-1842-JS-20-MAR-1842-WW-f}%
					``My
						\marginnote{f}%
					intention (says the speaker) was to have treated upon the subject of Baptism, But having a case of death before, us I thought it proper to refer to that subject, I will now however say a few words upon \sout{that subject} <Baptism> as intended God has made certain decreas which are fixed \& unalterable for instance God [has] set the sun, the moon, the stars in the heavens, \& gi[v]en them their laws conditions, \& bounds which they cannot pass except by his command, they all move in perfect harmony in there sphere \& order \& are as wonders, lights \& signs unto us,
						\hypertarget{EMS-1842-JS-20-MAR-1842-WW-g}%
					The
						\marginnote{g}%
					sea also has its bounds which it cannot pass God has set many signs in the earth as well as in heaven for instance the oaks of the forest, the fruit of the tree, the herd of the field all bear a sign that seed hath been planted there, for it is a decree of the Lord that evry tree fruit or herb bearing seed, should bring forth after its kind \& cannot come forth after any other law or principle.
						\hypertarget{EMS-1842-JS-20-MAR-1842-WW-h}%
					upon
						\marginnote{h}%
					the same principle do I contend that Baptism is a sign, ordained of God for the believer in Christ to take upon himself in order to enter into the kingdom of God, ``for except you are born of the water \& the spirit you cannot enter into the kingdom of God,[''] Saith the Savior, as It is a sign of command which God hath set for man to enter into \sout{the} this \sout{Kingdom of God}%
						\footnote{%
							% \label{}%
							JSP note: ```\sout{of God} is circled.''
							}
					those who seek to enter in any other \sout{will} way will seek in vain,
						\hypertarget{EMS-1842-JS-20-MAR-1842-WW-i}%
					for
						\marginnote{i}%
					God will not receive those neither will the angels acknowledge their works as accepted, for they have not taken upon themselves those ordinances \& signs which God ordained for man to receive in order to receive a celestial glory, \& God had decreed that all \sout{that} <who> will not obey his voice shall not escape the damnation of hell,
						\hypertarget{EMS-1842-JS-20-MAR-1842-WW-j}%
					what
						\marginnote{j}%
					is the damnation of hell, to go with that society who have not obeyed his commands Baptism is a sign to God to Angels to heaven that we do the will of God \& their is no other way beneath the heavens whareby God hath ordained for man to come \sout{to God} \& any other cource is in vain. God hath decreed \& ordained that man should repent of all his sins \& Be Baptized for the remission of his sins then he can come to God in the name of Jesus Christ in faith then we have the promise of the Holy Ghost

						\hypertarget{EMS-1842-JS-20-MAR-1842-WW-k}%
					What
						\marginnote{k}%
					is the sign of the healing of the sick? the laying on of hands, is the sign or way marked out by James \& the custom of the ancient saints as ordered by the Lord \& we should not obtain the blessing by persuing any other course except the way which God has markd out.

						\hypertarget{EMS-1842-JS-20-MAR-1842-WW-l}%
					What
						\marginnote{l}%
					if we should attempt to get the Holy Ghost through any other means except the sign or way which God hath appointed, should we obtain it certainly not all other means would fail The Lord says do so \& so \& I will bless so \& so their is certain key words \& signs belonging to the priesthood which much be observed in order to obtaine the Blessings

						\hypertarget{EMS-1842-JS-20-MAR-1842-WW-m}%
					The
						\marginnote{m}%
					sign of Peter was, to repent \& be baptized for the remission of Sins, \sout{in it} with the promise of the gift of the Holy Ghost \& in no other way is the gift of the Holy Ghost obtained, Their is a difference between the Holy Ghost \& the gift of the Holy Ghost, Cornelius received the Holy Ghost before he was Baptized which was the convincing power of God unto him of the truth of the gospel,
						\hypertarget{EMS-1842-JS-20-MAR-1842-WW-n}%
					but
						\marginnote{n}%
					he could not receive the gift of the Holy Ghost untill after he was Baptized, \& had he not taken this sign or [or]dinances upon him the Holy Ghost which convinced him of the truth of God would have left him, untill he obeyed these ordinances \& received the gift of the Holy Ghost by the laying on of hands according to the order of God he could not have healed the sick or command[ed] an evil spirit to come out of a man \& it obey him for the spirit might say to him as he did to the sons of Scava[,] Peter I know \& Christ I know but who are ye

						\hypertarget{EMS-1842-JS-20-MAR-1842-WW-o}%
					It
						\marginnote{o}%
					mat[t]ereth not whither we live long or short after we come to a knowled[g]e of the<se> principles \& obey them, I know that all men will be damned if they do not come in the way which God has appointed

						\hypertarget{EMS-1842-JS-20-MAR-1842-WW-p}%
					As
						\marginnote{p}%
					concerning the resurrection I will mer[e]ly say that all men will come from the grave as they \sout{lay} lie down, whether old or young their will not be added unto the[i]r stature one cubit neither taken from it, All being raised by the power of God having the Spirit of God in their bodies \& not Blood
						\hypertarget{EMS-1842-JS-20-MAR-1842-WW-q}%
					children
						\marginnote{q}%
					will be enthroned in the presence of God \& the Lamb with bodies of the same stature that were on earth, Having been redeemed by the Blood of the Lamb they will there enjoy a fulness of that light Glory \& intelligence which is received in the celestial kingdom of God ``Blessed are the dead who die in the Lord, for they rest from their labours \& their works do follow them['']

						\hypertarget{EMS-1842-JS-20-MAR-1842-WW-r}%
					The
						\marginnote{r}%
					Speaker before closing, called upon the vast assembly before him to humble themselves in faith before God, \& in mighty Prayer \& fasting to call upon his Holy name, untill the elements were purified over our heads \& the earth sanctified under our feet that the inhabitants of this city may escape the power of the disease, pestilences \& destroyer that rideth upon the face of the earth \& that the holy spirit of God may rest upon this vast multitude
					
			% --
			\paragraph{As Published in \textit{Times and Seasons}}
			% --
				\label{EMS-1842-JS-20-MAR-1842-TS}
			
				\begin{description}
					\item[Print Sources]
				\end{description}

				\begin{tabular}{p{0.5in}p{1in}p{0.5in}p{1in}}
					JSP	 	& ---		& JSR 		& --- \\
					DC	 	& ---		& HC 		& --- \\
					HDDC 	& ---		& TCHC 		& --- \\
				\end{tabular}
				% HDDC = woodford's DC diss; JSR = Marquardt's JS Revelations; TCHC = Vogel HC
				
				\begin{description}
					\item[Online Access] \href{https://www.josephsmithpapers.org/paper-summary/discourse-20-march-as-published-in-times-and-seasons/1}{Here}
					% \item[Analysis] \hyperref[XX]{Here}
				\end{description}
				
				\begin{description}
					\item[Background]
				\end{description}
				
					% It might also be useful to give a two sentence context of the period: e.g., "This speech was given in Nauvoo to a small congregation one month prior to the destruction of the Nauvoo Expositor. These and other tensions may have been on the prophet's mind when he gave the speech."
				
				\begin{description}
					\item[Original Source]
				\end{description}
				
					% brief description of original source material, with reference to JSP or somewhere else for more info
					% Background material: it might be helpful to have a note that says whether a given document was presented as a speech, was written in a journal, etc.
					% Also a comment here about how reliable the source might be, or what shape it is in, if there are large omissions or parts that are hard to understand, and so on.
				
				\begin{description}
					\item[Text]
				\end{description}
				
						\hypertarget{EMS-1842-JS-20-MAR-1842-TS-a}%
					President
						\marginnote{a}%
					Smith read the 14th chap. of Rev. and said,

					\begin{quote}
						``We have again the warning voice sounded in our midst which shows the uncertainty of human life; and in my leisure moments I have meditated upon the subject, and asked the question, Why it is that infants, innocent children are taken away from us? especially those that seem to be the most intelligent and interesting? and the strongest reasons that present themselves to my mind are these;---This world is a very wicked world; and it is a proverb that the `world grows weaker and wiser' if it is the case, the world grows more wicked and corrupt.
							\hypertarget{EMS-1842-JS-20-MAR-1842-TS-b}%
						In
							\marginnote{b}%
						the early ages of the world, a righteous man, and a man of God, and of intelligence had a better chance to do good, to be believed and received, than at the present day; but in these days such a man is much opposed and persecuted by most of the inhabitants of the earth; and he has much sorrow to pass through here, the Lord takes many away even in infancy that they may escape the envy of man, and the sorrows and evils of this present world; they were too pure, too lovely, to live on earth;
							\hypertarget{EMS-1842-JS-20-MAR-1842-TS-c}%
						therefore
							\marginnote{c}%
						if rightly considered instead of mourning we have reason to rejoice as they are delivered from evil, and we shall soon have them again.

						``What chance is there for infidelity when we are parting with our friends almost daily? None at all. The infidel will grasp at every straw for help until death stares him in the face, and then his infidelity takes its flight, for the realities of the eternal world are resting upon him in mighty power; and when every earthly support and prop fails him, he then sensibly feels the eternal truths of the immortality of the soul. We should take warning and not wait for the death bed to repent, as we see the infant taken away by death, so may the youth and middle aged, as well as the infant suddenly be called into eternity. Let this then prove as a warning to all, not to procrastinate repentance, or wait till a death bed; for it is the will of God that man should repent, and serve him in health, and in the strength, and power of his mind, in order to secure his blessing; and not wait until he is called to die.
							\hypertarget{EMS-1842-JS-20-MAR-1842-TS-d}%
						Also
							\marginnote{d}%
						the doctrine of Baptizing children, or sprinkling them, or they must welter in hell is a doctrine not true, not supported in Holy writ, and is not consistent with the character of God. All children are redeemed by the blood of Jesus Christ, and the moment that children leave this world they are taken to the bosom of Abraham. The only difference between the old and young dying, is, one lives longer in heaven, and eternal light and glory than the other, and is freed a little sooner from this miserable wicked world.--- Notwithstanding all this glory, we for a moment lose sight of it, and mourn the loss; but we do not mourn as those without hope.

							\hypertarget{EMS-1842-JS-20-MAR-1842-TS-e}%
						``My
							\marginnote{e}%
						intention was, to have spoken upon the subject of baptism, but having a case of death before us I thought proper to refer to that subject. I will now however say a few words upon baptism, as I intended. God has made certain decreees which are fixed, and immovable, for instance; God set the sun, the moon, and the stars in the heavens; and gave them their laws, conditions, and bounds which they cannot pass, except by his commandments; they all move in perfect harmony in their sphere, and order, and are as lights, wonders, and signs unto us.
							\hypertarget{EMS-1842-JS-20-MAR-1842-TS-f}%
						The
							\marginnote{f}%
						sea also has its bounds which it cannot pass. God has set many signs on the earth, as well as in the heavens, for instance; the oak of the forest, the fruit of the tree, the herb of the field; all bear a sign that seed hath been planted there; for it is a decree of the Lord that every tree, plant, and herb, bearing seed, should bring forth of its kind, and cannot come forth after any other law, or principle. 
							\hypertarget{EMS-1842-JS-20-MAR-1842-TS-g}%
						Upon
							\marginnote{g}%
						the same principle do I contend that baptism is a sign ordained of God, for the believer in Christ to take upon himself in order to enter into the kingdom of God, ``for except ye are born of water, and of the spirit ye cannot enter into the kingdom of God,'' saith the Saviour. It is a sign, and commandment which God has set for man to enter into his Kingdom. Those who seek to enter in any other way will seek in vain;
							\hypertarget{EMS-1842-JS-20-MAR-1842-TS-h}%
						and
							\marginnote{h}%
						God will not receive them, neither will the angels acknowledge their works as accepted: for they have not obeyed the ordinances, nor attended to the signs which God ordained for the salvation of man, to prepare him for; and give him a title to a celestial glory; and God has decreed that all who will not obey his voice shall not escape the damnation of hell.
							\hypertarget{EMS-1842-JS-20-MAR-1842-TS-i}%
						What
							\marginnote{i}%
						is the damnation of hell? to go with that society who have not obeyed his commands. Baptism is a sign to God, to Angels, and to heaven that we do the will of God: and there is no other way beneath the heavens whereby God hath ordained for man to come to him, to be saved, and enter into the kingdom of God, except faith in Jesus Christ; repentance, and baptism for the remission of sins, and any other course is in vain; then you have the promise of the gift of the Holy Ghost.
							\hypertarget{EMS-1842-JS-20-MAR-1842-TS-j}%
						What
							\marginnote{j}%
						is the sign of the healing of the sick? the laying on of hands is the sign, or way marked out by James; and the custom of the ancient Saints as ordered by the Lord; and we can not obtain the blessing by pursuing any other course, except the way marked out by the Lord.

							\hypertarget{EMS-1842-JS-20-MAR-1842-TS-k}%
						What
							\marginnote{k}%
						if we should attempt to get the gift of the Holy Ghost through any other means, except the signs, or way which God hath appointed? should we obtain it? certainly not; all other means would fail. The Lord says do so, and so, and I will bless, so, and so.

						There are certain key-words, and signs belonging to the priesthood, which must be observed in order to obtain the blessing,
							\hypertarget{EMS-1842-JS-20-MAR-1842-TS-l}%
						the
							\marginnote{l}%
						sign of Peter was to repent, and be baptized for the remission of sins, with the promise of the gift of the Holy Ghost: and in no other way is the gift of the Holy Ghost obtained. There is a difference between the Holy Ghost, and the gift of the Holy Ghost. Cornelius received the Holy Ghost before he was baptized; which was the convincing power of God unto him of the truth of the gospel;
							\hypertarget{EMS-1842-JS-20-MAR-1842-TS-m}%
						but
							\marginnote{m}%
						he could not receive the gift of the Holy Ghost until after he was 	baptized. Had he not taken this sign, or ordinance upon him, the Holy Ghost which convinced him of the truth of God, would have left him. Until he obeyed these ordinances and received the gift of the Holy Ghost, by the laying on of hands, according to the order of God, he could not have healed the sick, or commanded an evil spirit to come out of a man, and it obey him; for the spirits might say unto him, as they did to the sons of Sceva;---`Paul we know; and Jesus we know, but who are ye!'
							\hypertarget{EMS-1842-JS-20-MAR-1842-TS-n}%
						It
							\marginnote{n}%
						mattereth not whether we live long or short on the earth after we come to a knowledge of these principles and obey them unto the end. I know that all men will be damned if they do not come in the way which he hath opened; and this is the way marked out by the word of the Lord.

							\hypertarget{EMS-1842-JS-20-MAR-1842-TS-o}%
						``As
							\marginnote{o}%
						concerning the resurrection I will merely say that all men will come from the grave as they lie down, whether old or young, there will not be `added unto their stature one cubit;' neither taken from it; all will be raised by the power of God, having spirit in their bodies, and not blood.
							\hypertarget{EMS-1842-JS-20-MAR-1842-TS-p}%
						Children
							\marginnote{p}%
						will be enthroned in the presence of God, and the Lamb; with bodies of the same stature that they had on earth; having been redeemed by the blood of the Lamb, they will there enjoy the fulness of that light glory, and intelligence which is prepared in the Celestial kingdom: `Blessed are the dead who die in the Lord; for they rest from their labors and their works do follow them.'
					\end{quote}

						\hypertarget{EMS-1842-JS-20-MAR-1842-TS-q}%
					The
						\marginnote{q}%
					speaker before closing called upon the assembly before him, to humble themselves in faith before God, and in mighty prayer and fasting to call upon the name of the Lord, until the elements were purified over our heads, and the earth sanctified under our feet; that the inhabitants of this city may escape the power of disease and pestilence, and the destroyer that rideth upon the face of the earth; and that the Holy Spirit of God may rest upon this vast multitude.
					
		% -----
		\subsubsection{Discourse, 27 March 1842}
		% -----
			\label{EMS-1842-JS-27-MAR-1842}
			% \label{EMS-1842-JS-DAY-3LETTERMONTH-YEAR}
			
			% --
			\paragraph{As Reported by Wilford Woodruff}
			% --
				\label{EMS-1842-JS-27-MAR-1842-WW}
			
				\begin{description}
					\item[Print Sources]
				\end{description}

				\begin{tabular}{p{0.5in}p{1in}p{0.5in}p{1in}}
					JSP	 	& ---		& JSR 		& --- \\
					DC	 	& ---		& HC 		& --- \\
					HDDC 	& ---		& TCHC 		& --- \\
				\end{tabular}
				% HDDC = woodford's DC diss; JSR = Marquardt's JS Revelations; TCHC = Vogel HC
				
				\begin{description}
					\item[Online Access] \href{https://www.josephsmithpapers.org/paper-summary/discourse-27-march-1842/1}{Here}
					% \item[Analysis] \hyperref[XX]{Here}
				\end{description}
				
				\begin{description}
					\item[Background]
				\end{description}
				
					% It might also be useful to give a two sentence context of the period: e.g., "This speech was given in Nauvoo to a small congregation one month prior to the destruction of the Nauvoo Expositor. These and other tensions may have been on the prophet's mind when he gave the speech."
				
				\begin{description}
					\item[Original Source]
				\end{description}
				
					% brief description of original source material, with reference to JSP or somewhere else for more info
					% Background material: it might be helpful to have a note that says whether a given document was presented as a speech, was written in a journal, etc.
					% Also a comment here about how reliable the source might be, or what shape it is in, if there are large omissions or parts that are hard to understand, and so on.
				
				\begin{description}
					\item[Text]
				\end{description}
				
						\hypertarget{EMS-1842-JS-27-MAR-1842-a}%
					He
						\marginnote{a}%
					\phantom{ }
					[Amasa Lyman] was followed by Joseph the Seer who made some edifying remarks concerning Baptism for the dead, He said the Bible supported the doctrin, ``why are ye Baptized for the dead if the dead rise not \&c'' if their is one word of the Lord that supports the doctrin it is enough to mik [make] it a true doctrin Again
						\hypertarget{EMS-1842-JS-27-MAR-1842-b}%
					if
						\marginnote{b}%
					we can baptize a man in the Name of the Father of the Son \& of the Holy Ghost for the remission of sins it is Just as much our privilege to act as an agent \& be baptized for the remission of sins for \& in behalf of our dead kindred who have not herd the gospel or fulness of it
					
		% -----
		\subsubsection{Letter to Emma Smith and the Relief Society, 31 March 1842}
		% -----
			\label{EMS-1842-JS-31-MAR-1842}
			% \label{EMS-1842-JS-DAY-3LETTERMONTH-YEAR}
			
			\begin{description}
				\item[Print Sources]
			\end{description}

			\begin{tabular}{p{0.5in}p{1in}p{0.5in}p{1in}}
				JSP	 	& ---		& JSR 		& --- \\
				DC	 	& ---		& HC 		& --- \\
				HDDC 	& ---		& TCHC 		& --- \\
			\end{tabular}
			% HDDC = woodford's DC diss; JSR = Marquardt's JS Revelations; TCHC = Vogel HC
			
			\begin{description}
				\item[Online Access] \href{https://www.josephsmithpapers.org/paper-summary/letter-to-emma-smith-and-the-relief-society-31-march-1842/1}{Here}
				% \item[Analysis] \hyperref[XX]{Here}
			\end{description}
			
			\begin{description}
				\item[Background]
			\end{description}
			
				% It might also be useful to give a two sentence context of the period: e.g., "This speech was given in Nauvoo to a small congregation one month prior to the destruction of the Nauvoo Expositor. These and other tensions may have been on the prophet's mind when he gave the speech."
			
			\begin{description}
				\item[Original Source]
			\end{description}
			
				% brief description of original source material, with reference to JSP or somewhere else for more info
				% Background material: it might be helpful to have a note that says whether a given document was presented as a speech, was written in a journal, etc.
				% Also a comment here about how reliable the source might be, or what shape it is in, if there are large omissions or parts that are hard to understand, and so on.
			
			\begin{description}
				\item[Text]
			\end{description}
			
					\hypertarget{EMS-1842-JS-31-MAR-1842-a}%
				To
					\marginnote{a}%
				the President of the F R So N [Female Relief Society of Nauvoo] Greeting

				Can the ``Female Releif Society of Nauvoo'' be Trusted with some important matters that ought actually to belong to them to see to which men have been under the necessity of seeing to to their chagrin \& Mortification in order to prevent iniquitous characters from carrying their iniquity into effect such as, for instance a man who may be aspiring after power \& authority and yet without principle; regardles[s] of god, man or the Devil or the interest or welfare of men, or the virtue or innocence of women?
					\hypertarget{EMS-1842-JS-31-MAR-1842-b}%
				Shall
					\marginnote{b}%
				the credulity, good faith, \& steadfast feelings of our Sisters for the cause of God or truth be imposed upon by believing such men because they say they have authority from Joseph or the first Presidency or any other Presidency of the church and thus with a lie in their mouth deceive \& debauch the innocent under the assumption that they are authorized from these sources! \uline{May} \uline{God} \uline{forbid!}

					\hypertarget{EMS-1842-JS-31-MAR-1842-c}%
				A
					\marginnote{c}%
				knowledge of some such thing having come to our ears we imp[r]ove this favorable opportunity wherein so goodly number of you may be informed that no such authority ever has, ever can, or ever will be given to any man \& if any man has been guilty of any such thing let him be treated with utter contempt \& let the curse of God fall on his head, \& let him be turned out of Society as unworthy of a place among men, \& \sout{renounced} denounced as the blackest \& the most unprincipled wretch \& finally let him be \uline{damned}.

					\hypertarget{EMS-1842-JS-31-MAR-1842-d}%
				We
					\marginnote{d}%
				have been informed that some unprincipled men whose names we will not mention at present have been guilty of such crimes: we do not mention their names, not knowing but what there may be some among you who are not sufficiently skilled in Masonry as to keep a secret, therefore suffice it to say there are those \& we therefore warn you \& forewarn you in the name of the Lord to check and destroy any faith that any innocent person may have in any such character
					\hypertarget{EMS-1842-JS-31-MAR-1842-e}%
				for
					\marginnote{e}%
				we don't want any body to believe any thing as coming from us contrary to the old established morals \& virtues \& scriptural laws regulating the habits customs \& conduct of Society unless it be by message del[iv]ered to you by our own mouth, by actual revelation \& commandment. and all persons pretending to be authorized by us or having any permit or sanction from us are \& will be liars \& base imposters \& you are authorized on the very first intimation of the kind to denounce them as such \& \sout{fly} \sout{from} <shun> them as the fiery flying serpents,
					\hypertarget{EMS-1842-JS-31-MAR-1842-f}%
				whether
					\marginnote{f}%
				they are prophets, seers, or Revelators, patriarchs, Twelve apostles, Elders, Priests. \sout{or what not,} Mayors, Generals, \sout{or what not,} city council alderman, Marshall, Police, Lord Mayor or the Devil, are alike culpable. \& shall be damned for such evil practices; \& if you \sout{yourself} yourselves \sout{hear} adhere to any thing of the kind you, also shall be damned.

					\hypertarget{EMS-1842-JS-31-MAR-1842-g}%
				Now
					\marginnote{g}%
				beloved Sisters do not believe for a moment that we wish to impose upon you, we actualy do know that such things have existed in the church \& <are> sorry that we are obliged to make mention of any such thing \& we want a stop put to them, \& we want you to do your part \& we will do ours \sout{part} for we wish to to keep the commandments of God in all things <as> given <directly from heaven> to us \sout{from} \sout{heaven}, living by every word that proceedeth out of the mouth of the Lord.---

					\hypertarget{EMS-1842-JS-31-MAR-1842-h}%
				May
					\marginnote{h}%
				God add his blessings upon your head \& lead you in all the paths of virtue piety and peace that you may be an ornament unto those to whom you belong \& arise up and crown them with power \& by so doing you shall be crowned with honor in heaven \& shall sit upon throne jud[g]ing them whom you are placed in authority \sout{of} over in the world and shall be judged of God for all the responsbilitys that are conferred upon you

					\hypertarget{EMS-1842-JS-31-MAR-1842-i}%
				At
					\marginnote{i}%
				a more convenient \& appropriate season we will give you further information upon this subject

				We are your humble servants in the bonds of the new \& Everlasting covenant

				Let that epistle be had as a private matter in your society \& then we shall learn whether you are good masons---

				\begin{tightright}
					Joseph Smith P. C. J. C. L.

					B[righam] Young Prst Twelve.
				\end{tightright}

				P.S. if the Lord be God serve him \& if Baal \textit{[2 words illegible]}
				
		% -----
		\subsubsection{Minutes and Discourse, 31 March 1842}
		% -----
			\label{EMS-1842-JS-31-MAR-1842-MD}
			% \label{EMS-1842-JS-DAY-3LETTERMONTH-YEAR}
			
			\begin{description}
				\item[Print Sources]
			\end{description}

			\begin{tabular}{p{0.5in}p{1in}p{0.5in}p{1in}}
				JSP	 	& ---		& JSR 		& --- \\
				DC	 	& ---		& HC 		& --- \\
				HDDC 	& ---		& TCHC 		& --- \\
			\end{tabular}
			% HDDC = woodford's DC diss; JSR = Marquardt's JS Revelations; TCHC = Vogel HC
			
			\begin{description}
				\item[Online Access] \href{https://www.josephsmithpapers.org/paper-summary/minutes-and-discourse-31-march-1842/1}{Here}
				% \item[Analysis] \hyperref[XX]{Here}
			\end{description}
			
			\begin{description}
				\item[Background]
			\end{description}
			
				See the \textit{Relief Society} documents book important notes.
			
				% It might also be useful to give a two sentence context of the period: e.g., "This speech was given in Nauvoo to a small congregation one month prior to the destruction of the Nauvoo Expositor. These and other tensions may have been on the prophet's mind when he gave the speech."
			
			\begin{description}
				\item[Original Source]
			\end{description}
			
				% brief description of original source material, with reference to JSP or somewhere else for more info
				% Background material: it might be helpful to have a note that says whether a given document was presented as a speech, was written in a journal, etc.
				% Also a comment here about how reliable the source might be, or what shape it is in, if there are large omissions or parts that are hard to understand, and so on.
			
			\begin{description}
				\item[Text]
			\end{description}
			
				\begin{tightright}
						\hypertarget{EMS-1842-JS-31-MAR-1842-MD-a}%
					Lodge
						\marginnote{a}%
					Room March 30\supers{th} [31st] 1842.
				\end{tightright}

				Meeting opened with Singing---

				Prayer by Pres\supers{t.} Joseph Smith---

				The house full to overflowing.

					\hypertarget{EMS-1842-JS-31-MAR-1842-MD-b}%
				Prest.
					\marginnote{b}%
				J. Smith arose--- spoke of the organization of the Society--- said he was deeply interested that it might be built up to the Most High in an acceptable manner--- that its rules must be observed--- that none should be received into the Society but those who were worthy--- propos'd that the Society go into a close examination of every candidate--- that they were going too fast--- that the Society should grow up by degrees--- should commence with a few individuals---
					\hypertarget{EMS-1842-JS-31-MAR-1842-MD-c}%
				thus
					\marginnote{c}%
				have a select Society of the virtuous and those who will walk circumspectly--- commended them for their zeal but said sometimes their zeal was not according to knowledge--- One principal object of the Institution, was to purge out iniquity--- said they must be extremely careful in all their examinations or the consequences would be serious

					\hypertarget{EMS-1842-JS-31-MAR-1842-MD-d}%
				Said
					\marginnote{d}%
				all difficulties which might \& would cross our way must be surmounted, though the soul be \uline{tried}, the heart faint, and hands hang down--- must not retrace our steps--- that there must be decision of character aside from sympathy--- that when instructed we must obey that voice, observe the Constitution that the blessings of heaven may rest down upon us--- all must act in concert or nothing can be done---
					\hypertarget{EMS-1842-JS-31-MAR-1842-MD-e}%
				that
					\marginnote{e}%
				the Society should move according to the ancient Priesthood, hence there should be a select Society separate from all the evils of the world, choice, virtuou[s] and holy--- Said he was going to make of this Society a kingdom of priests an in Enoch's day--- as in Pauls day--- that it is the privilege of each member to live long and enjoy health--- Pres\supers{t.} Smith propos'd that the ladies [gentlemen] withdraw, that the Society might proceed to business--- that those wishing to join should have their names presented at the next meeting.------

					\hypertarget{EMS-1842-JS-31-MAR-1842-MD-f}%
				Pres\supers{t.}
					\marginnote{f}%
				J. Smith withdrew---...
				
				Pres\supers{t.} E. S. said we were going to learn new things--- our way was straight--- said we wanted none in this society but those who \uline{could} and \uline{would} walk straight and were determined to do good and not evil---...
				
				After appropriate remarks Pres\supers{t.} E. Smith, said she had an Article to read which would test the ability of the members in keeping secrets; as it was for the benefit of the Society, and that alone------
				
				Read the Article.--- Then gave strictures on female propriety and dignity \&c. \&c.
				
					\hypertarget{EMS-1842-JS-31-MAR-1842-MD-g}%
				Mother
					\marginnote{g}%
				\phantom{ }
				[Lucy Mack] Smith rose and said she was glad the time had come that iniquity could be detected and reproach thrown off from the heads of the church--- We come into the church to be sav'd--- that we may live in peace and sit down in the kingdom of heaven--- If we listen to, and circulate every evil report we shall idly spend the time which should be appropriated to the reading the Scriptures, the Book of Mormon--- we must remember the words of Alma--- pray much at morning, noon and night--- feed the poor \&c.---
					\hypertarget{EMS-1842-JS-31-MAR-1842-MD-h}%
				She
					\marginnote{h}%
				said she was \uline{old}--- could not meet with the Society but few times more--- and wish'd to leave her testimony that the Book of Mormon is the book of God--- that Joseph is a man of God, a prophet of the Lord set apart to lead the people--- If we observe his words it will be well with us; if we live righteously on earth, it will be well with us in Eternity------...
				
		% -----
		\subsubsection{Discourses, 6--8 April 1842}
		% -----
			\label{EMS-1842-JS-6-8-APR-1842}
			% \label{EMS-1842-JS-DAY-3LETTERMONTH-YEAR}
			
			% --
			\paragraph{As Reported in \textit{Times and Seasons}}
			% --
				\label{EMS-1842-JS-6-8-APR-1842-TS}
			
				\begin{description}
					\item[Print Sources]
				\end{description}

				\begin{tabular}{p{0.5in}p{1in}p{0.5in}p{1in}}
					JSP	 	& ---		& JSR 		& --- \\
					DC	 	& ---		& HC 		& --- \\
					HDDC 	& ---		& TCHC 		& --- \\
				\end{tabular}
				% HDDC = woodford's DC diss; JSR = Marquardt's JS Revelations; TCHC = Vogel HC
				
				\begin{description}
					\item[Online Access] \href{https://www.josephsmithpapers.org/paper-summary/minutes-and-discourses-6-8-april-1842/1}{Here}
					% \item[Analysis] \hyperref[XX]{Here}
				\end{description}
				
				\begin{description}
					\item[Background]
				\end{description}
				
					% It might also be useful to give a two sentence context of the period: e.g., "This speech was given in Nauvoo to a small congregation one month prior to the destruction of the Nauvoo Expositor. These and other tensions may have been on the prophet's mind when he gave the speech."
				
				\begin{description}
					\item[Original Source]
				\end{description}
				
					% brief description of original source material, with reference to JSP or somewhere else for more info
					% Background material: it might be helpful to have a note that says whether a given document was presented as a speech, was written in a journal, etc.
					% Also a comment here about how reliable the source might be, or what shape it is in, if there are large omissions or parts that are hard to understand, and so on.
				
				\begin{description}
					\item[Text]
				\end{description}
				
						\hypertarget{EMS-1842-JS-6-8-APR-1842-TS-a}%
					...Conference
						\marginnote{a}%
					met, Pres't. Joseph Smith had the several quorums put in order, and seated: he then made some very appropriate remarks concerning the duties of the church, the necessity of unity of purpose in regard to the building of the houses, and the blessings connected with doing the will of God; and the inconsistency folly and danger of murmuring against the dispensations of Jehovah.
					
						\hypertarget{EMS-1842-JS-6-8-APR-1842-TS-b}%
					He
						\marginnote{b}%
					said that the principal object of the meeting was to bring the case of Elder Page before them, and that another object was to choose young men, and ordain them, and send them out to preach, that they may have an opportunity of proving themselves, and of enduring the tarring and feathering and such things as those of us who have gone before them, have had to endure...
					
						\hypertarget{EMS-1842-JS-6-8-APR-1842-TS-c}%
					Pres't.
						\marginnote{c}%
					J. Smith then arose and stated that it was wrong to make the covenant referred to by him; that it created a lack of confidence for two men to covenant to reveal all acts of secrecy or otherwise to each other---and Elder Page showed a little grannyism. He said that no two men when they agreed to go together ought to separate, that the prophets of old would not
						\hypertarget{EMS-1842-JS-6-8-APR-1842-TS-d}%
					and
						\marginnote{d}%
					quoted the circumstance of Elijah and Elisha iii Kings 2 chap. when about to go to Gilgal, also when about to go to Jericho, and to Jordan, that Elisha could not get clear of Elijah, that he clung to his garment until he was taken to heaven and that Elder Page should have stuck by Elder Hyde, and he might have gone to Jerusalem, that there is nothing very bad in it, but by the experience let us profit; again, the Lord made use of Elder Page as a scape-goat to procure funds for Elder Hyde.
					
						\hypertarget{EMS-1842-JS-6-8-APR-1842-TS-e}%
					When
						\marginnote{e}%
					Elder Hyde returns we will reconsider the matter, and perhaps send them back to Jerusalem, we will fellowship Elder Page until Elder Hyde comes, and we will then weld them together and make them one. A vote was then put, and carried that we hold Elder Page in full fellowship...
					
						\hypertarget{EMS-1842-JS-6-8-APR-1842-TS-f}%
					Pres't.
						\marginnote{f}%
					J. Smith spoke upon the subject of the stories respecting Elder Kimball and others, showing the folly and inconsistency of spending any time in conversing about such stories or hearkening to them, for there is no person that is acquainted with our principles would believe such lies, except [Thomas] Sharp the editor of the ``Warsaw Signal.'' Baptisms for the dead, and for the healing of the body must be in the font, those coming into the church and those rebaptized may be done in the river.
					
						\hypertarget{EMS-1842-JS-6-8-APR-1842-TS-g}%
					A
						\marginnote{g}%
					box should be prepared for the use of the font, that the clerk may be paid, and a book procured by the monies to be put therein by those baptized' the remainder to go to the use of the Temple.---Sung a hymn. Ordinations to take place to-morrow morning---Baptisms in the font also---There were 275 ordained to the office of Elder under the hands of the Twelve during the Conference...
					
					Pres't. J. Smith said the baptisms would be attended to, also the ordinations...
					
		% -----
		\subsubsection{Discourse, 9 April 1842}
		% -----
			\label{EMS-1842-JS-9-APR-1842}
			% \label{EMS-1842-JS-DAY-3LETTERMONTH-YEAR}
			
			% --
			\paragraph{As Reported by Wilford Woodruff}
			% --
				\label{EMS-1842-JS-9-APR-1842-WW}
			
				\begin{description}
					\item[Print Sources]
				\end{description}

				\begin{tabular}{p{0.5in}p{1in}p{0.5in}p{1in}}
					JSP	 	& ---		& JSR 		& --- \\
					DC	 	& ---		& HC 		& --- \\
					HDDC 	& ---		& TCHC 		& --- \\
				\end{tabular}
				% HDDC = woodford's DC diss; JSR = Marquardt's JS Revelations; TCHC = Vogel HC
				
				\begin{description}
					\item[Online Access] \href{https://www.josephsmithpapers.org/paper-summary/discourse-9-april-1842-as-reported-by-wilford-woodruff/1}{Here}
					% \item[Analysis] \hyperref[XX]{Here}
				\end{description}
				
				\begin{description}
					\item[Background]
				\end{description}
				
					% It might also be useful to give a two sentence context of the period: e.g., "This speech was given in Nauvoo to a small congregation one month prior to the destruction of the Nauvoo Expositor. These and other tensions may have been on the prophet's mind when he gave the speech."
				
				\begin{description}
					\item[Original Source]
				\end{description}
				
					% brief description of original source material, with reference to JSP or somewhere else for more info
					% Background material: it might be helpful to have a note that says whether a given document was presented as a speech, was written in a journal, etc.
					% Also a comment here about how reliable the source might be, or what shape it is in, if there are large omissions or parts that are hard to understand, and so on.
				
				\begin{description}
					\item[Text]
				\end{description}
				
						\hypertarget{EMS-1842-JS-9-APR-1842-WW-a}%
					President
						\marginnote{a}%
					Joseph Smith spoke upon the occasion with much feelings \& interest among his remarks he said it is a vary solomn \& awful time I never felt more solomn it calles to mind the death of my oldest Brother who died in New York \& my Youngest Brother Carloss [Don Carlos] Smith who died in Nauvoo It has been hard for me to live on earth \& see those young men upon whome we have leaned upon as a support \& comfort taken from us in the midst of their youth,
						\hypertarget{EMS-1842-JS-9-APR-1842-WW-b}%
					yes
						\marginnote{b}%
					it has been hard to be reconciled to these things I have sometimes felt that I should have felt more reconciled to have been called myself if it could have been the will of God, yet I know we ought to be still \& know it is of God \& be reconciled all is right it will be but a short time before we shall all in like manner be called, It may be the case with me as well as you Some has supposed that Br Joseph could not die but this is a mistake
						\hypertarget{EMS-1842-JS-9-APR-1842-WW-c}%
					it
						\marginnote{c}%
					is true their has been times when I have had the promise of my life to accomplish such \& such things, but having accomplish[ed] those things I have not at present any lease of my life \& am as liable to die as other men I can say in my heart that I have not done any thing against Ephraim Marks that I am sorry for \& I would ask any of his Companions if they have done any thing against him that they are sorry for.
						\hypertarget{EMS-1842-JS-9-APR-1842-WW-d}%
					or
						\marginnote{d}%
					that they would not like to meet at the bar of God if so let it prove as a warning to all men to deal justly before God \& with all men then we shall be clear in the day of judgment When we loose a near \& dear friend upon whom we have set our hearts we can never feel the same afterwards Knowing that if we set our hearts upon others they may in like manner be taken from us President Smith made many other interesting remarks[.]
					
		% -----
		\subsubsection{Discourse, 10 April 1842}
		% -----
			\label{EMS-1842-JS-10-APR-1842}
			% \label{EMS-1842-JS-DAY-3LETTERMONTH-YEAR}
			
			% --
			\paragraph{As Reported by Wilford Woodruff}
			% --
				\label{EMS-1842-JS-10-APR-1842-WW}
			
				\begin{description}
					\item[Print Sources]
				\end{description}

				\begin{tabular}{p{0.5in}p{1in}p{0.5in}p{1in}}
					JSP	 	& ---		& JSR 		& --- \\
					DC	 	& ---		& HC 		& --- \\
					HDDC 	& ---		& TCHC 		& --- \\
				\end{tabular}
				% HDDC = woodford's DC diss; JSR = Marquardt's JS Revelations; TCHC = Vogel HC
				
				\begin{description}
					\item[Online Access] \href{https://www.josephsmithpapers.org/paper-summary/discourse-10-april-1842-as-reported-by-wilford-woodruff/1}{Here}
					% \item[Analysis] \hyperref[XX]{Here}
				\end{description}
				
				\begin{description}
					\item[Background]
				\end{description}
				
					% It might also be useful to give a two sentence context of the period: e.g., "This speech was given in Nauvoo to a small congregation one month prior to the destruction of the Nauvoo Expositor. These and other tensions may have been on the prophet's mind when he gave the speech."
				
				\begin{description}
					\item[Original Source]
				\end{description}
				
					% brief description of original source material, with reference to JSP or somewhere else for more info
					% Background material: it might be helpful to have a note that says whether a given document was presented as a speech, was written in a journal, etc.
					% Also a comment here about how reliable the source might be, or what shape it is in, if there are large omissions or parts that are hard to understand, and so on.
				
				\begin{description}
					\item[Text]
				\end{description}
				
						\hypertarget{EMS-1842-JS-10-APR-1842-WW-a}%
					Then
						\marginnote{a}%
					Joseph the Seer arose in the power of God reproved \& rebuked wickedness \sout{in} before the people in the name of the Lord God He wished to say a few words to suit the condition of the general mass And I shall speak with authority of the priesthood in the name of the Lord God, which shall prove a savior of life unto life or of death unto death,
					
						\hypertarget{EMS-1842-JS-10-APR-1842-WW-b}%
					Notwithstanding
						\marginnote{b}%
					this congregation profess to be Saints yet I stand in the midst of all characters and classes of men If you wish to go whare God is you must be like God or possess the principles which God possesses for if we are not drawing towards God in principle we are going from him \& drawing towards the devil, yes I am standing in the midst of all kinds of people search your hearts \& see if you are like God, I have searched mine \& feel to repent of all my sins,
						\hypertarget{EMS-1842-JS-10-APR-1842-WW-c}%
					We
						\marginnote{c}%
					have theives among us Adulterers, liars, hypocritts, if God should speak from Heaven he would Command \sout{us} you not to steal, not to commit Adultery, nor to covet, nor deceive but be faithful over a few things
						\hypertarget{EMS-1842-JS-10-APR-1842-WW-d}%
					As
						\marginnote{d}%
					far as we degenerate from God we desend to the devil \& looses knowledge `\& without Knowledge we cannot be saved \& while our hearts are filled with evil \& we are studying evil their is no room in our hearts for good or studying good, is not God good, Yea then you be good, if he is faithful then you be faithful Add to your faith virtue to virtue knowledge. \& seek for evry good thing' the Church must be cleansed \& I proclaim against all iniquity.
						\hypertarget{EMS-1842-JS-10-APR-1842-WW-e}%
					A
						\marginnote{f}%
					man is saved no faster than he gets knowledge for if he does not get knowledge he will be brought into Captivity by some evil power in the other world as evil spirits will have more knowled[g]e \& consequently more power than many men who are on the earth. hence it needs Revelation to assist us \& give us knowledge of the things of God.
						\hypertarget{EMS-1842-JS-10-APR-1842-WW-f}%
					What
						\marginnote{f}%
					is the reason that the Priest[s] of the day do not get Revelation They ask ownly to consume it upon their lust their hearts are corrupt \& they cloke their iniquity by saying that their is no more Revelations But if any revelations are given of God they are universally opposed by the priest \& christendon at large for it reveals their wickedmess \& abominations
					
		% -----
		\subsubsection{Letter to Nancy Rigdon, circa Mid-April 1842}
		% -----
			\label{EMS-1842-JS-CIR-MID-APR-1842}
			% \label{EMS-1842-JS-DAY-3LETTERMONTH-YEAR}
			
			\begin{description}
				\item[Print Sources]
			\end{description}

			\begin{tabular}{p{0.5in}p{1in}p{0.5in}p{1in}}
				JSP	 	& ---		& JSR 		& --- \\
				DC	 	& ---		& HC 		& --- \\
				HDDC 	& ---		& TCHC 		& --- \\
			\end{tabular}
			% HDDC = woodford's DC diss; JSR = Marquardt's JS Revelations; TCHC = Vogel HC
			
			\begin{description}
				\item[Online Access] \href{https://www.josephsmithpapers.org/paper-summary/appendix-letter-to-nancy-rigdon-circa-mid-april-1842/1}{Here}
				% \item[Analysis] \hyperref[XX]{Here}
			\end{description}
			
			\begin{description}
				\item[Background]
			\end{description}
			
				% It might also be useful to give a two sentence context of the period: e.g., "This speech was given in Nauvoo to a small congregation one month prior to the destruction of the Nauvoo Expositor. These and other tensions may have been on the prophet's mind when he gave the speech."
			
			\begin{description}
				\item[Original Source]
			\end{description}
			
				UNCERTAIN AUTHORSHIP---NOT KNOWN IF JS WROTE THIS. See the JSP historical introduction. See \href{https://drive.google.com/file/d/19dBq4Oej0R_q5ndI3bEY5R_1WRhYRbKI/view?usp=sharing}{Dirkmaat's 2016 article}, ``Searching for `Happiness': Joseph Smith's Alleged Authorship of the 1842 Letter to Nancy Rigdon,'' for more discussion. From 98--100 of that article:
				
				\begin{quote}
					Despite its popularity, however, it presents special problems of provenance and authenticity to historians. These problems compound the difficulty with which the document can be contextually understood. In effect, historians cannot demonstrate with certainty that Joseph Smith wrote the letter, as they can with other Joseph Smith documents. The only ``original'' text is found in the pages of the \textit{Sangamo Journal}. Even if Smith’s authorship is assumed, however, Bennett’s claim placing it in the context of a failed proposal to Nancy Rigdon also cannot be demonstrated with certainty, as there are competing accounts of the interactions between the two of them. It is not the intent of this article to make a definitive claim that Joseph Smith did not author the ``Happiness Letter'' nor that he did not make a polygamous proposal to Nancy Rigdon. Smith was already engaged in Nauvoo-era polygamy, and these marriages were kept secret from the general membership of the Church and the community at large. Marinda Hyde, who factored prominently in Bennett’s account of the proposal, eventually became a polyandrous wife of Smith, and she had already been approached by him with an explanation of ``the doctrine of Celestial marriage'' in the fall of 1841. Given these factors, a proposal by Smith to Nancy Rigdon would certainly not be extraordinary. Still, the purpose of this article is to illustrate the difficulties involved in any attempt to gain a clear understanding of this document and to outline the reasons to both doubt and to impute Smith’s authorship of the document. The ultimate determination of authorship hinges on the reasons why it was first excluded and then later included in the Manuscript History of the Church. As those reasons can only be met with speculation, a measure of caution should be employed when assigning authorship. Rather than attributing it casually and definitively to Joseph Smith, historians and theologians, Mormons and non-Mormons alike, should be aware of the questioned provenance when using the document and draw measured rather than expansive conclusions. Simply put, if this were any other document, historians would greatly question the claims of authorship for the reasons that will follow. 
				\end{quote}
				
				\noindent Dirkmaat's article also includes a discussion of how the letter ended up in the manuscript history, since its presence there is the reason why almost everyone has taken it as genuine.
				
				Then, again from Dirkmaat, page, 119:
				
				\begin{quote}
					The examination of the problematic provenance and questionable context of this document should in no way lead to the definitive conclusion that the letter was not authored or dictated by Joseph Smith. As stated before, the history of Joseph Smith's practice of polygamy in Nauvoo provides a context in which both the proposal to Nancy may have occurred and the subsequent letter indeed may have been produced. And the letter does seem to resemble Joseph Smith's language more readily than Bennett's other forgeries, but this has not been proven in a quantitative, academic way. Nevertheless, all users of this document should be aware of its questioned provenance, the inscrutable circumstances surrounding its inclusion and placement in the Manuscript History of the Church, and how it came to be regarded as unquestionably Joseph Smith's. It is simply not responsible to assert that the ``Happiness Letter'' was definitively authored by Smith when no original letter exists nor do any contemporary Mormons attribute it to him. Historical inertia has caused the document to be regarded as definitively Joseph Smith's rather than careful evaluation. Responsible historians should, after weighing the evidence, treat the letter, its contents, and its purported context very carefully. They should draw very measured and qualified conclusions when using the document either as a representation of Joseph Smith's doctrinal teachings or as context for Joseph Smith's practice of plural marriage in Nauvoo rather than relying on the presuppositions of an earlier age of writers, historians, apostates, or apologists. 
				\end{quote}
				
				If it's genuine, it's extremely important for many reasons; but it's doubtful. Best not to use this as a primary reason to establish any particular point of view.
			
				% brief description of original source material, with reference to JSP or somewhere else for more info
				% Background material: it might be helpful to have a note that says whether a given document was presented as a speech, was written in a journal, etc.
				% Also a comment here about how reliable the source might be, or what shape it is in, if there are large omissions or parts that are hard to understand, and so on.
			
			\begin{description}
				\item[Text]
			\end{description}
			
					\hypertarget{EMS-1842-JS-CIR-MID-APR-1842-a}%
				Happiness
					\marginnote{a}%
				is the object and design of our existence, and will be the end thereof if we pursue the path that leads to it; and this path is virtue, uprightness, faithfulness, holiness, and keeping \textsc{all} the commandments of God.
					\hypertarget{EMS-1842-JS-CIR-MID-APR-1842-b}%
				But
					\marginnote{b}%
				we cannot keep \textit{all} the commandments without first \textit{knowing} them, and we cannot expect to know \textit{all}, or more than we \textit{now} know, unless we \textit{comply} with or keep those we have \textit{already received}.
					\hypertarget{EMS-1842-JS-CIR-MID-APR-1842-c}%
				That
					\marginnote{c}%
				which is \textit{wrong} under one circumstance, may be, and often is, \textit{right} under another. God said \textit{thou shalt not kill},--- at another time he said \textit{thou shalt utterly destroy}. This is the principle on which the government of heaven is conducted---by \textit{revelation} adapted to the circumstances in which the children of the kingdom are placed. \textit{Whatever God requires is right}, \textsc{no matter what it is}, although we may not see the reason thereof till long after the events transpire.
					\hypertarget{EMS-1842-JS-CIR-MID-APR-1842-d}%
				If
					\marginnote{d}%
				we seek first the kingdom of God, all \textit{good things} will be added. So with Solomon---first he asked wisdom, and God gave it him, and with it every \textsc{desire} of his heart, even things which might be considered \textit{abominable} to all who do not\textit{understand} the order of heaven only in part, but which, in reality, were right, because God gave and sanctioned by \textit{special revelation}.
					\hypertarget{EMS-1842-JS-CIR-MID-APR-1842-e}%
				A
					\marginnote{e}%
				parent may whip a child, and justly too, because he stole an apple; whereas, if the child had asked for the apple, and the parent had given it, the child would have eaten it with a better appetite, there would have been no stripes---all the pleasures of the apple would have been received, all the misery of stealing lost. This principle will justly apply to all of God's dealings with his children.
					\hypertarget{EMS-1842-JS-CIR-MID-APR-1842-f}%
				Every
					\marginnote{f}%
				thing that God gives us is \textit{lawful and right}, and 'tis proper that we should \textit{enjoy} his gifts and blessings \textit{whenever and wherever} he is disposed to bestow; but if we should seize upon these same blessings and enjoyments without law, without \textit{revelation}, without commandment, those blessings and enjoyments would prove cursings and vexations in the end, and we should have to lie down in sorrow and wailings of everlasting regret. 
					\hypertarget{EMS-1842-JS-CIR-MID-APR-1842-g}%
				But
					\marginnote{g}%
				in obedience there is joy and peace unspotted, unalloyed, and as God has designed our happiness, the happiness of all his creatures, he never has, he never will, institute an ordinance, or give a commandment to his people that is not calculated in its nature to promote that happiness which he has designed, and which will not end in the greatest amount of good and glory to those who become the recipients of his laws and ordinances.
					\hypertarget{EMS-1842-JS-CIR-MID-APR-1842-h}%
				Blessings
					\marginnote{h}%
				offered, but rejected, are no longer blessings, but become like the talent hid in the earth by the wicked and slothful servant---the proffered good returns to the giver, the blessing is bestowed on those who will \textit{receive}, and \textit{occupy}; for unto him that hath shall be given, and he shall have \textit{abundantly}; but unto him that hath not, or will not receive, shall be taken away that which he hath, or might have had.

					\hypertarget{EMS-1842-JS-CIR-MID-APR-1842-i}%
				``\textit{Be wise}
					\marginnote{i}%
				\textsc{to-day}, \textit{'tis} \textsc{madness} \textit{to defer}.

				Next day the fatal precedent may plead:

				Thus on till wisdom is pushed out of time''

				Into eternity.

					\hypertarget{EMS-1842-JS-CIR-MID-APR-1842-j}%
				Our
					\marginnote{j}%
				heavenly father is more liberal in his views, and boundless in his mercies and blessings, than we are ready to believe or receive, and at the same time is more terrible to the workers of iniquity, more awful in the executions of his punishments, and more ready to detect every false \textit{way} than we are apt to suppose him to be.
					\hypertarget{EMS-1842-JS-CIR-MID-APR-1842-k}%
				He
					\marginnote{k}%
				will be enquired of by his children---he says \textit{ask} and \textsc{ye shall receive}, \textit{seek} and ye \textsc{shall} find, but if ye will take that which is not your own, or which I \textit{have not given you}, you shall be rewarded according to your deeds,
					\hypertarget{EMS-1842-JS-CIR-MID-APR-1842-l}%
				but
					\marginnote{l}%
				no good thing will I withhold from them who walk uprightly before me, and do my will in \textit{all things}, who will listen to \textit{my voice}, and to \textsc{to the voice of my servant whom I have sent}, for I delight in those who seek diligently to know my precepts, and abide by the laws of my kingdom, for all things shall be made known unto them in mine own due time, and in the end they shall have joy.
					
		% -----
		\subsubsection{Minutes and Discourse, 28 April 1842}
		% -----
			\label{EMS-1842-JS-28-APR-1842}
			% \label{EMS-1842-JS-DAY-3LETTERMONTH-YEAR}
			
			\begin{description}
				\item[Print Sources]
			\end{description}

			\begin{tabular}{p{0.5in}p{1in}p{0.5in}p{1in}}
				JSP	 	& ---		& JSR 		& --- \\
				DC	 	& ---		& HC 		& --- \\
				HDDC 	& ---		& TCHC 		& --- \\
			\end{tabular}
			% HDDC = woodford's DC diss; JSR = Marquardt's JS Revelations; TCHC = Vogel HC
			
			\begin{description}
				\item[Online Access] \href{https://www.josephsmithpapers.org/paper-summary/minutes-and-discourse-28-april-1842/1}{Here}
				% \item[Analysis] \hyperref[XX]{Here}
			\end{description}
			
			\begin{description}
				\item[Background]
			\end{description}
			
				% It might also be useful to give a two sentence context of the period: e.g., "This speech was given in Nauvoo to a small congregation one month prior to the destruction of the Nauvoo Expositor. These and other tensions may have been on the prophet's mind when he gave the speech."
			
			\begin{description}
				\item[Original Source]
			\end{description}
			
				See notes in the RS volume.
			
				% brief description of original source material, with reference to JSP or somewhere else for more info
				% Background material: it might be helpful to have a note that says whether a given document was presented as a speech, was written in a journal, etc.
				% Also a comment here about how reliable the source might be, or what shape it is in, if there are large omissions or parts that are hard to understand, and so on.
			
			\begin{description}
				\item[Text]
			\end{description}
			
					\hypertarget{EMS-1842-JS-28-APR-1842-a}%
				The
					\marginnote{a}%
				Meeting \sout{meet} met according to adjournment, present Pres\supers{t.} Joseph Smith and Elder W[illard] Richards---
				
				President Smith arose and said that the purport of his being present on the occasion was, to make observations respecting the Priesthood, and give instructions for the benefit of the Society That as his instructions were intended only for the Society; he requested that a vote should be taken on those present who were not members, to ascertain whether they should be admitted---
					\hypertarget{EMS-1842-JS-28-APR-1842-b}%
				he
					\marginnote{b}%
				exhorted the meeting to act honestly and uprightly in all their proceedings--- inasmuch as they would be call'd to give an account to Jehovah. All hearts must repent--- be pure and God will regard them and bless them in a manner that they could not be bless'd in any other way---
				
				A vote was then taken on the following names, to wit...
				
					\hypertarget{EMS-1842-JS-28-APR-1842-c}%
				Committee
					\marginnote{c}%
				retir'd--- and Prest. J. Smith arose and call'd the attention of the meeting to the 12\supers{th} Chap. of 1\supers{st} Cor. ``Now concerning spiritual gifts'' \&c.--- Said that the passage which reads ``no man can \uline{say} that Jesus is the the \sout{Christ} <Lord> but by the holy ghost,'' should be translated, no man can \uline{know} \&c

					\hypertarget{EMS-1842-JS-28-APR-1842-d}%
				He
					\marginnote{d}%
				continued to read the Chap. and give instructions respecting the different offices, and the necessity of every individual acting in the sphere allotted him or her; and filling the several offices to which they were appointed--- Spoke of the disposition of man, to consider the lower offices in the church dishonorable and to look with jealous eyes upon the standing of others--- that it was the nonsense of the human heart, for a person to be aspiring to other stations than appointed of God--- that it was better for individuals to magnify their respective callings, and wait patiently till God shall say to them come up higher.
					\hypertarget{EMS-1842-JS-28-APR-1842-e}%
				He
					\marginnote{e}%
				said the reason of these remarks being made, was, that some little thing was circulating in the Society, that some persons were not going right in laying hands on the sick \&c. Said if he had common sympathies, would rejoice that the sick could be heal'd: that the time had not been before, that these things could be in their proper order--- that the church is not now organiz'd in its proper order, and cannot be until the Temple is completed------
					\hypertarget{EMS-1842-JS-28-APR-1842-f}%
				Pres\supers{t.}
					\marginnote{f}%
				Smith continued the subject by adverting to the commission given to the ancient apostles ``Go ye into all the world'' \&c.--- no matter who believeth; these signs, such as healing the sick, casting out devils \&c. should follow all that believe whether male or female. He ask'd the Society if they could not see by this sweeping stroke, that wherein they are ordaind, it is the privilege of those set apart to administer in that authority which is confer'd on them--- and if the sisters should have faith to heal the sick, let all hold their tongues, and let every thing roll on.

					\hypertarget{EMS-1842-JS-28-APR-1842-g}%
				He
					\marginnote{g}%
				said, if God has appointed him, and chosen him as an instrument to lead the church, why not let him lead it through? Why stand in the way, when he is appointed to do a thing? Who knows the mind of God? Does he not reveal things differently from what we expect?--- He remark'd that he was continually rising--- altho' he had every thing bearing him down--- standing in his way and opposing--- after all he always comes out right in the end.

					\hypertarget{EMS-1842-JS-28-APR-1842-h}%
				Respecting
					\marginnote{h}%
				the female laying on hands, he further remark'd, there could be no devil in it if God gave his sanction by healing--- that there could be no more sin in any female laying hands on the sick than in wetting the face with water--- that it is no sin for any body to do it that has faith, or if the sick has faith to be heal'd by the administration.

					\hypertarget{EMS-1842-JS-28-APR-1842-i}%
				He
					\marginnote{i}%
				reprov'd those that were dispos'd to find fault with the management of concerns--- saying if he undertook to lead the church he would lead it right--- that he calculates to organize the church in proper order \&c.
				
					\hypertarget{EMS-1842-JS-28-APR-1842-j}%
				President
					\marginnote{j}%
				Smith continued by speaking of the difficulties he had to surmount ever since the commencement of the work in consequence of aspiring men, ``great big Elders'' as he call'd them, who had caused him much trouble, whom he had taught in the private counsel; and they would go forth into the world and ploclaim [proclaim] the things he had taught them; as their own revelations--- said the same aspiring disposition will be in this Society, and must be guarded against--- that every person should stand and act in the place appointed, and thus sanctify the Society and get it pure---

					\hypertarget{EMS-1842-JS-28-APR-1842-k}%
				He
					\marginnote{k}%
				said he had been trampled underfoot by aspiring Elders, for all were infected with that spirit, for instance P[arley P.] Pratt O[rson] Pratt, O[rson] Hyde and J[ohn E.] Page had been aspiring--- they could not be exalted but must run away as tho' the care and authority of the church were vested with them--- he said we had a subtle devil to deal with, and could only curb him by being humble.

					\hypertarget{EMS-1842-JS-28-APR-1842-l}%
				He
					\marginnote{l}%
				said as he had this opportunity, he was going to instruct the Society and point out the way for them to conduct, that they might act according to the will of God--- that he did not know as he should have many opportunities of teaching them--- that they were going to be left to themselves,--- they would not long have him to instruct them--- that the church would not have his instruction long, and the world would not be troubled with him a great while, and would not have his teachings--- He spoke of delivering the keys to this Society and to the church--- that according to his prayers God had appointed him elsewhere

					\hypertarget{EMS-1842-JS-28-APR-1842-m}%
				He
					\marginnote{m}%
				exhorted the sisters always to concentrate their faith and prayers for, and place confidence, in those whom God has appointed to honor, whom God has plac'd at the head to lead--- that we should arm them with our prayers.--- that the keys of the kingdom are about to be given to them, that they may be able to detect every thing false--- as well as to the Elders

					\hypertarget{EMS-1842-JS-28-APR-1842-n}%
				He
					\marginnote{n}%
				said if one member becomes corrupt and you know it; you must immediately put it away. The sympathies of the heads of the church have induc'd them to bear with those that were corrupt; in consequence of which all become contaminated--- you must put down iniquity and by your good example provoke the Elders to good works---
					\hypertarget{EMS-1842-JS-28-APR-1842-o}%
				if
					\marginnote{o}%
				you do right, no danger of going too fast: he said he did not care how fast we run in the path of virtue. Resist evil and there is no danger. God, men, angels, and devils can't condemn those that resist every thing that is evil--- as well might the devil seek to dethrone Jehovah, as that soul that resists every thing that is evil.

					\hypertarget{EMS-1842-JS-28-APR-1842-p}%
				The 
					\marginnote{p}%
				charitable Society--- this is according to your natures--- it is natural for females to have feelings of charity--- you are now plac'd in a situatio[n] where you can act according to those sympathies which God has planted in your bosoms. If you live up to these principles how great and glorious!--- if you live up to your privilege, the angels cannot be restrain'd from being your associates--- females, if they are pure and innocent can come into the presence of God; for what is more pleasing to God than innocence; you must be innocent or you cannot come up before God. If \sout{ye} <we> would come before God let us be pure ourselves.
					\hypertarget{EMS-1842-JS-28-APR-1842-q}%
				The
					\marginnote{q}%
				devil has great power--- he will so transform things as to make one gape at those who are doing the will of God--- You need not be tearing men for their deeds, but let the weight of innocence be felt; which is more mighty than a millstone hung about the neck. Not war, not jangle, not contradiction, but meekness, love, purity, these are the things that should magnify us.--- Achan must be brought to light--- iniquity must be purged out--- \uline{then} the vail will be rent and the blessings of heaven will flow down --- they will roll down like the Missisippi river.
					\hypertarget{EMS-1842-JS-28-APR-1842-r}%
				This
					\marginnote{r}%
				Society shall have power to command Queens in their midst--- I now deliver it as a prophecy that before ten years shall roll round, the queens of the earth shall come and pay their respects to this Society--- they shall come with their millions and shall contribute of their abundance for the relief of the poor--- If you will be pure, nothing can hinder.

					\hypertarget{EMS-1842-JS-28-APR-1842-s}%
				After
					\marginnote{s}%
				this instruction, you will be responsible for your own sins. It is an honor to save \sout{yourself} yourselves--- all are responsible to save themselves.

				Pres\supers{t.} Smith, after reading from the above mentioned Chapter, continued to give instruction respecting the order of God, as established in the church; saying every one should aspire only to magnify his own office \&c.------

					\hypertarget{EMS-1842-JS-28-APR-1842-t}%
				He
					\marginnote{t}%
				then commenc'd reading the 13\supers{th} chapter, ``Though I speak with the tongues of men'' \&c; and said don't be limited in your views with regard to your neighbors' virtues, but be limited towards your own virtues; and not think yourselves more righteous than others; you must enlarge your souls toward others if yould [you would?] do like Jesus, and carry your fellow creatures to Abram's bosom.

					\hypertarget{EMS-1842-JS-28-APR-1842-u}%
				He
					\marginnote{u}%
				said he had manifested long suffering and we must do so too------ Pres\supers{t.} Smith then read, ``Though I have the gift of prophecy'' \&c. He then said, though one should become mighty--- do great things--- overturn mountains \&c and should then turn to eat and drink with the drunken; all former deeds would not save him--- but he would go to destruction!

					\hypertarget{EMS-1842-JS-28-APR-1842-v}%
				As
					\marginnote{v}%
				you increase in innocence and virtue, as you increase in goodness, let your hearts expand--- let them be enlarged towards others--- you must be longsuff'ring and bear with the faults and errors of mankind. How precious are the souls of men!--- The female part of community are apt to be contracted in their views. You must not be contracted, but you must be liberal in your feelings.

					\hypertarget{EMS-1842-JS-28-APR-1842-w}%
				Let
					\marginnote{w}%
				this Society teach how to act towards husbands to treat them with mildness and affection. When a man is borne down with trouble--- when he is perplex'd; if he can meet a smile, an argument--- if he can meet with mildness, it will calm down his soul and soothe his feelings. When the mind is going to despair, it needs a solace.

					\hypertarget{EMS-1842-JS-28-APR-1842-x}%
				This
					\marginnote{x}%
				Society is to get instruction thro' the order which God has established--- thro' the medium of those appointed to lead--- and I now turn the key to you in the name of God and this Society shall rejoice and knowledge and intelligence shall flow down from this time--- this is the beginning of better days, to this Society

					\hypertarget{EMS-1842-JS-28-APR-1842-y}%
				When
					\marginnote{y}%
				you go home never give a cross word, but let kindness, charity and love, crown your works henceforward. Don't envy sinners--- have mercy on them, God will destroy them.--- Let your labors be confin'd mostly to those around you in your own circle; as far as knowledge is concerned, it may extend to all the world, but your administrations, should be confin'd to the circle of your immediate acquaintance, and more especially to the members of the Society.

				Those ordain'd to lead the Society, are authoriz'd to appoint to different offices as the circumstances shall require.

					\hypertarget{EMS-1842-JS-28-APR-1842-z}%
				If
					\marginnote{z}%
				any have a matter to reveal, let it be in your own tongue. Do not indulge too much in the gift of tongues, or the devil will take advantage of the innocent. You may speak in tongues for your comfort but I lay this down for a rule that if any thing is is taught by the gift of tongues, it is not to be received for doctrine.

					\hypertarget{EMS-1842-JS-28-APR-1842-aa}%
				Pres\supers{t.}
					\marginnote{aa}%
				S. then offered instruction respecting the propriety of females administering to the sick by the laying on of hands--- said it was according to revelation \&c. said he never was plac'd in similar circumstances, and never had given the same instruction.

				He clos'd his instructions by expressing his satisfaction in improving the opportunity.

				The spirit of the Lord was pour'd out in a very powerful manner, never to be forgotten by those present on that interesting occasion...
				
		% -----
		\subsubsection{Discourse, 1 May 1842}
		% -----
			\label{EMS-1842-JS-1-MAY-1842}
			% \label{EMS-1842-JS-DAY-3LETTERMONTH-YEAR}
			
			% --
			\paragraph{As Reported by Willard Richards}
			% --
				\label{EMS-1842-JS-1-MAY-1842-WR}
			
				\begin{description}
					\item[Print Sources]
				\end{description}

				\begin{tabular}{p{0.5in}p{1in}p{0.5in}p{1in}}
					JSP	 	& ---		& JSR 		& --- \\
					DC	 	& ---		& HC 		& --- \\
					HDDC 	& ---		& TCHC 		& --- \\
				\end{tabular}
				% HDDC = woodford's DC diss; JSR = Marquardt's JS Revelations; TCHC = Vogel HC
				
				\begin{description}
					\item[Online Access] \href{https://www.josephsmithpapers.org/paper-summary/discourse-1-may-1842-as-reported-by-willard-richards/1}{Here}
					% \item[Analysis] \hyperref[XX]{Here}
				\end{description}
				
				\begin{description}
					\item[Background]
				\end{description}
				
					% It might also be useful to give a two sentence context of the period: e.g., "This speech was given in Nauvoo to a small congregation one month prior to the destruction of the Nauvoo Expositor. These and other tensions may have been on the prophet's mind when he gave the speech."
				
				\begin{description}
					\item[Original Source]
				\end{description}
				
					% brief description of original source material, with reference to JSP or somewhere else for more info
					% Background material: it might be helpful to have a note that says whether a given document was presented as a speech, was written in a journal, etc.
					% Also a comment here about how reliable the source might be, or what shape it is in, if there are large omissions or parts that are hard to understand, and so on.
				
				\begin{description}
					\item[Text]
				\end{description}
				
						\hypertarget{EMS-1842-JS-1-MAY-1842-WR-a}%
					preached
						\marginnote{a}%
					in the grove on the keys of the kingdom charity \&\supers{c}.--- The keys are certain signs \& words by which false spirits \& personages may be detected from true.--- which cannot be revealed to the Elders till the Temple is completed.---
						\hypertarget{EMS-1842-JS-1-MAY-1842-WR-b}%
					The
						\marginnote{b}%
					rich can only get them in the Temple. The poor may get them on the Mountain top as did moses. The rich cannot be saved without cha[r]ity. giving to feed the poor. when \& how God requires as well as building. There are signs in heaven earth \& hell. the elders must know them all to be endued with power. to finish their work \& prevent imposition.
						\hypertarget{EMS-1842-JS-1-MAY-1842-WR-c}%
					The
						\marginnote{c}%
					devil knows many signs. but does not know the sign of the son of man. or Jesus. No one can truly say he knows God until he has handled something. \& \sout{these} this can only be in the holiest of Holies.
					
		% -----
		\subsubsection{Revelation, 19 May 1842}
		% -----
			\label{EMS-1842-JS-19-MAY-1842}
			% \label{EMS-1842-JS-DAY-3LETTERMONTH-YEAR}
			
			\begin{description}
				\item[Print Sources]
			\end{description}

			\begin{tabular}{p{0.5in}p{1in}p{0.5in}p{1in}}
				JSP	 	& ---		& JSR 		& --- \\
				DC	 	& ---		& HC 		& --- \\
				HDDC 	& ---		& TCHC 		& --- \\
			\end{tabular}
			% HDDC = woodford's DC diss; JSR = Marquardt's JS Revelations; TCHC = Vogel HC
			
			\begin{description}
				\item[Online Access] \href{https://www.josephsmithpapers.org/paper-summary/revelation-19-may-1842/1}{Here}
				% \item[Analysis] \hyperref[XX]{Here}
			\end{description}
			
			\begin{description}
				\item[Background]
			\end{description}
			
				% It might also be useful to give a two sentence context of the period: e.g., "This speech was given in Nauvoo to a small congregation one month prior to the destruction of the Nauvoo Expositor. These and other tensions may have been on the prophet's mind when he gave the speech."
			
			\begin{description}
				\item[Original Source]
			\end{description}
			
				% brief description of original source material, with reference to JSP or somewhere else for more info
				% Background material: it might be helpful to have a note that says whether a given document was presented as a speech, was written in a journal, etc.
				% Also a comment here about how reliable the source might be, or what shape it is in, if there are large omissions or parts that are hard to understand, and so on.
			
			\begin{description}
				\item[Text]
			\end{description}
			
				Verily thus saith the Lord unto you my servant Joseph by the voice of my Spirit, Hiram Kimball has been insinuating evil. \& forming evil opinions against you with. others. \& if he continue in them he \& they shall be accursed. for I am the Lord thy God \& will stand by thee \& bless thee. \uline{Amen}.
				
		% -----
		\subsubsection{Minutes and Discourse, 26 May 1842}
		% -----
			\label{EMS-1842-JS-26-MAY-1842}
			% \label{EMS-1842-JS-DAY-3LETTERMONTH-YEAR}
			
			\begin{description}
				\item[Print Sources]
			\end{description}

			\begin{tabular}{p{0.5in}p{1in}p{0.5in}p{1in}}
				JSP	 	& ---		& JSR 		& --- \\
				DC	 	& ---		& HC 		& --- \\
				HDDC 	& ---		& TCHC 		& --- \\
			\end{tabular}
			% HDDC = woodford's DC diss; JSR = Marquardt's JS Revelations; TCHC = Vogel HC
			
			\begin{description}
				\item[Online Access] \href{https://www.josephsmithpapers.org/paper-summary/minutes-and-discourse-26-may-1842/1}{Here}
				% \item[Analysis] \hyperref[XX]{Here}
			\end{description}
			
			\begin{description}
				\item[Background]
			\end{description}
			
				% It might also be useful to give a two sentence context of the period: e.g., "This speech was given in Nauvoo to a small congregation one month prior to the destruction of the Nauvoo Expositor. These and other tensions may have been on the prophet's mind when he gave the speech."
			
			\begin{description}
				\item[Original Source]
			\end{description}
			
				See notes in the RS volume.
			
				% brief description of original source material, with reference to JSP or somewhere else for more info
				% Background material: it might be helpful to have a note that says whether a given document was presented as a speech, was written in a journal, etc.
				% Also a comment here about how reliable the source might be, or what shape it is in, if there are large omissions or parts that are hard to understand, and so on.
			
			\begin{description}
				\item[Text]
			\end{description}
			
					\hypertarget{EMS-1842-JS-26-MAY-1842-a}%
				...Prest.
					\marginnote{a}%
				J. Smith \& wife then entered.

				Prest. J. Smith rose, read the 14\supers{th} Chap. of Ezekiel--- said the Lord had declar'd by the prophet that the people should each one stand for himself and depend on no man or men in that state of corruption of the Jewish church--- that righteous persons could only deliver their own souls---
					\hypertarget{EMS-1842-JS-26-MAY-1842-b}%
				app[l]ied
					\marginnote{b}%
				it to the present state of the church of Latter-Day Saints--- said if the people departed from the Lord, they must fall--- that they were depending on the prophet hence were darkened in their minds from neglect of themselves--- envious toward the innocent, while they afflict the virtuous with their shafts of envy.

					\hypertarget{EMS-1842-JS-26-MAY-1842-c}%
				There
					\marginnote{c}%
				is another error which opens a door for the adversary to enter. As females possess refin'd feelings and sensitivenes[s], they are also subject to an overmuch zeal which must ever prove dangerous, and cause them to be rigid in a religious capacity--- should be arm'd with mercy notwithstanding the iniquity among us. Said he had been instrumental in bringing it to light--- melancholy and awful that so many are under the condemnation of the devil \& going to perdition.

					\hypertarget{EMS-1842-JS-26-MAY-1842-d}%
				With
					\marginnote{d}%
				deep feeling said that they are our fellows--- we lov'd them once. Shall we not encourage them to reformation?

				We have not forgivn them seventy times--- perhaps we have not forgiven them once. There is now a day of salvation to such as repent and reform--- they should be cast out from this Society, yet we should woo them to return to God lest they escape not the damnation of hell!

					\hypertarget{EMS-1842-JS-26-MAY-1842-e}%
				When
					\marginnote{e}%
				there is a mountain top there also is a vally--- we should act in all things an a proper medium--- to every immortal spirit. Notwithstanding the unworthy are among us, the virtuous should not from self-importance grieve and oppress needlessly those unfortunate ones, even these should be encourag'd to hereafter live to be honored by this Society who are the best portions of community.
					\hypertarget{EMS-1842-JS-26-MAY-1842-f}%
				Said
					\marginnote{f}%
				he had two things to recommend to the Society, to put a double watch over the tongue. No organiz'd body can exist without this at all. All organiz'd bodies have their peculiar evils, weaknesses and difficulties--- the object is to make those not so good, equal with the good and ever hold the keys of pow'r which will influence to virtue and goodness.
					\hypertarget{EMS-1842-JS-26-MAY-1842-g}%
				Should
					\marginnote{g}%
				chasten and reprove and keep \sout{in} it all in silence, not even mention them again, then you will be established in power, virtue and holiness and the wrath of God will be turn'd away. One request to the Prest. and Society, that you search yourselves--- the tongue is an unruly member--- hold your tongues about things of no moment,--- a little tale will set the world on fire.
					\hypertarget{EMS-1842-JS-26-MAY-1842-h}%
				At
					\marginnote{h}%
				this time the truth on the guilty should not be told openly--- Strange as this may seem, yet this is policy. We must use precaution in bringing sinners to justice lest in exposing these heinous sins, we draw the indignation of a gentile world upon us (and to their imaginatio[n] justly too)

				It is necessary to hold an influence in the world and thus spare ourselves an extermination; and also accomplish our end in spreading the gospel or holiness in the earth.

					\hypertarget{EMS-1842-JS-26-MAY-1842-i}%
				If
					\marginnote{i}%
				we were brought to desolation, the disobedient would find no help. There are some who are obedient yet men cannot steady the ark--- my arm can not do it--- God must steady it. To the iniquitous show yourselves merciful. I am advis'd by some of the heads of the church to tell the Relief Society to be virtuous--- but to save the church from desolation and the sword beware, be still, be prudent.
					\hypertarget{EMS-1842-JS-26-MAY-1842-j}%
				Repent,
					\marginnote{j}%
				reform but do it in a way to not destroy all around you. I do not want to cloak iniquity--- all things contrary to the will of God, should be cast from us, but dont do more hurt than good with your tongues--- be pure in heart--- Jesus designs to save the people out of their sins. Said Jesus ye shall do the work which ye see me do--- These are the grand key words for the Society to act upon.

					\hypertarget{EMS-1842-JS-26-MAY-1842-k}%
				If
					\marginnote{k}%
				I were not in your midst to aid and council you, the devil would overcome you. I want the innocent to go free--- rather spare ten iniquitous among you than than condemn one innocent one. ``Fret not thyself because of evil doers.'' God will see to it.

				[\textit{1 line blank}]

					\hypertarget{EMS-1842-JS-26-MAY-1842-l}%
				Mrs.
					\marginnote{l}%
				Prest. rose and said all idle rumor and idle talk must be laid aside yet sin must not be covered, especially those sins which are against the law of God and the laws of the country--- all who walk disorderly must reform, and any knowing of heinous sins against the law of God, and refuse to expose them, becomes the offender--- said she wanted none in this Society who had violated the laws of virtue...
				
		% -----
		\subsubsection{Discourse, 5 June 1842}
		% -----
			\label{EMS-1842-JS-5-JUN-1842}
			% \label{EMS-1842-JS-DAY-3LETTERMONTH-YEAR}
			
			% --
			\paragraph{As Reported by \textit{The Wasp}}
			% --
				\label{EMS-1842-JS-5-JUN-1842-W}
			
				\begin{description}
					\item[Print Sources]
				\end{description}

				\begin{tabular}{p{0.5in}p{1in}p{0.5in}p{1in}}
					JSP	 	& ---		& JSR 		& --- \\
					DC	 	& ---		& HC 		& --- \\
					HDDC 	& ---		& TCHC 		& --- \\
				\end{tabular}
				% HDDC = woodford's DC diss; JSR = Marquardt's JS Revelations; TCHC = Vogel HC
				
				\begin{description}
					\item[Online Access] \href{https://www.josephsmithpapers.org/paper-summary/discourse-5-june-1842-as-reported-by-the-wasp/1}{Here}
					% \item[Analysis] \hyperref[XX]{Here}
				\end{description}
				
				\begin{description}
					\item[Background]
				\end{description}
				
					% It might also be useful to give a two sentence context of the period: e.g., "This speech was given in Nauvoo to a small congregation one month prior to the destruction of the Nauvoo Expositor. These and other tensions may have been on the prophet's mind when he gave the speech."
				
				\begin{description}
					\item[Original Source]
				\end{description}
				
					% brief description of original source material, with reference to JSP or somewhere else for more info
					% Background material: it might be helpful to have a note that says whether a given document was presented as a speech, was written in a journal, etc.
					% Also a comment here about how reliable the source might be, or what shape it is in, if there are large omissions or parts that are hard to understand, and so on.
				
				\begin{description}
					\item[Text]
				\end{description}
				
						\hypertarget{EMS-1842-JS-5-JUN-1842-W-a}%
					The
						\marginnote{a}%
					Prophet, Joseph Smith, delivered a powerful discourse, last Sabbath, to an attentive audience of about 8000, in this city. The subject matter was drawn from the 32nd and 33rd chapters of Ezekiel, wherein it was shown that Old Pharaoh was comforted, and greatly rejoiced, that he was honored as a kind of King Devil over those uncircumcised nations, that go down to hell for rejecting the word of the Lord, withstanding his mighty miracles, and fighting the saints:---%
						\hypertarget{EMS-1842-JS-5-JUN-1842-W-b}%
					The
						\marginnote{b}%
					whole exhibited as a pattern to this generation, and the nations now rolling in splendor over the globe, if they do not repent, that they shall go down to the pit also, and be rejoiced over, and ruled over by Old Pharaoh. King Devil of mobocrats, miracle rejecters, saint killers, hypocritical priests, and all other fit subjects to fester in their own infamy.--- \textit{Spring water tastes best right from the fountain.}
					
			% --
			\paragraph{As Reported by John D. Lee}
			% --
				\label{EMS-1842-JS-5-JUN-1842-JDL}
			
				\begin{description}
					\item[Print Sources]
				\end{description}

				\begin{tabular}{p{0.5in}p{1in}p{0.5in}p{1in}}
					JSP	 	& ---		& JSR 		& --- \\
					DC	 	& ---		& HC 		& --- \\
					HDDC 	& ---		& TCHC 		& --- \\
				\end{tabular}
				% HDDC = woodford's DC diss; JSR = Marquardt's JS Revelations; TCHC = Vogel HC
				
				\begin{description}
					\item[Online Access] \href{https://www.josephsmithpapers.org/paper-summary/discourse-5-june-1842-as-reported-by-john-d-lee/1}{Here}
					% \item[Analysis] \hyperref[XX]{Here}
				\end{description}
				
				\begin{description}
					\item[Background]
				\end{description}
				
					% It might also be useful to give a two sentence context of the period: e.g., "This speech was given in Nauvoo to a small congregation one month prior to the destruction of the Nauvoo Expositor. These and other tensions may have been on the prophet's mind when he gave the speech."
				
				\begin{description}
					\item[Original Source]
				\end{description}
				
					% brief description of original source material, with reference to JSP or somewhere else for more info
					% Background material: it might be helpful to have a note that says whether a given document was presented as a speech, was written in a journal, etc.
					% Also a comment here about how reliable the source might be, or what shape it is in, if there are large omissions or parts that are hard to understand, and so on.
				
				\begin{description}
					\item[Text]
				\end{description}
				
						\hypertarget{EMS-1842-JS-5-JUN-1842-JDL-a}%
					Revelation.
						\marginnote{a}%
					Given to J. Simth [Smith] Nauvoo.
					
					June 5th. A. D. 1842.
					
					As the word of the Lord. was unto pharoah--- King of Egypt--- by the mouth of Ezikiel. the prophet. E.Z. 32. <\&> 33.------ so--- is the word of the Lord: unto this--- generation. by the mouth of. Joseph Smith. if they Repent not. they shall be herld [hurled] down to Hell.
						\hypertarget{EMS-1842-JS-5-JUN-1842-JDL-b}%
					Also---
						\marginnote{b}%
					to you Latter days saints Repent \& forsake your sins--- or you shall likewise suffer--- verily thus saith the Lord--- to those that encourage Mobs. to opose \&--- distress the Mormons. (remember) the same. Mobs that you encourage--- shall return to your own houses--- \& Bosom \& shall distress you--- \& shall spread death \& distruction in your midst--- yea there be some here to day who shall witness the same \& acknowledge this to--- be the word of the Lord--- others to their shame \& everlasting condemnation---
						\hypertarget{EMS-1842-JS-5-JUN-1842-JDL-c}%
					Remember.
						\marginnote{c}%
					o. ye L.D. saints \& perish not. Repent of your <evil> doeings--- lest ye be herld down to Hell with pharaoh \& his hosts--- but if you will turn to the Lord \& remember to give him glory \& you shall live this shall be a sign unto you when it shall come to Pass------
					
		% -----
		\subsubsection{Minutes and Discourse, 9 June 1842}
		% -----
			\label{EMS-1842-JS-9-JUN-1842}
			% \label{EMS-1842-JS-DAY-3LETTERMONTH-YEAR}
			
			\begin{description}
				\item[Print Sources]
			\end{description}

			\begin{tabular}{p{0.5in}p{1in}p{0.5in}p{1in}}
				JSP	 	& ---		& JSR 		& --- \\
				DC	 	& ---		& HC 		& --- \\
				HDDC 	& ---		& TCHC 		& --- \\
			\end{tabular}
			% HDDC = woodford's DC diss; JSR = Marquardt's JS Revelations; TCHC = Vogel HC
			
			\begin{description}
				\item[Online Access] \href{https://www.josephsmithpapers.org/paper-summary/minutes-and-discourse-9-june-1842/1}{Here}
				% \item[Analysis] \hyperref[XX]{Here}
			\end{description}
			
			\begin{description}
				\item[Background]
			\end{description}
			
				% It might also be useful to give a two sentence context of the period: e.g., "This speech was given in Nauvoo to a small congregation one month prior to the destruction of the Nauvoo Expositor. These and other tensions may have been on the prophet's mind when he gave the speech."
			
			\begin{description}
				\item[Original Source]
			\end{description}
			
				See notes in the RS volume.
			
				% brief description of original source material, with reference to JSP or somewhere else for more info
				% Background material: it might be helpful to have a note that says whether a given document was presented as a speech, was written in a journal, etc.
				% Also a comment here about how reliable the source might be, or what shape it is in, if there are large omissions or parts that are hard to understand, and so on.
			
			\begin{description}
				\item[Text]
			\end{description}
			
				\begin{tightright}
						\hypertarget{EMS-1842-JS-9-JUN-1842-a}%
					Grove,
						\marginnote{a}%
					June 9\supers{th.}
				\end{tightright}

				Prest J. Smith opened the meeting by pray'r and proceeded to address the congregation on the design of the Institution--- said it is no matter how fast the Society increases if all are virtuous--- that we must be as particular with regard to the character of members, as when the Society first started--- that sometimes persons wish to put themselves into a Society of this kind, when they do not intend to pursue the ways of purity and righteousness, as if the Society would be a shelter to them in their iniquity.

					\hypertarget{EMS-1842-JS-9-JUN-1842-b}%
				Prest.
					\marginnote{b}%
				S. said that henceforth no person shall be admitted but by presenting regular petition signed by two or three members in good standing in the Society--- whoever comes in must be of good report.

				Harriet Luce and Mary Luce were receiv'd into the Society by recommend.

					\hypertarget{EMS-1842-JS-9-JUN-1842-c}%
				Objections
					\marginnote{c}%
				previously made against Mahala Overton were remov'd--- after which Prest Smith continued his address--- said he was going to preach mercy Supposing that Jesus Christ and angels should object to us on frivolous things, what would become of us? We must be merciful and overlook small things.

					\hypertarget{EMS-1842-JS-9-JUN-1842-d}%
				Respecting
					\marginnote{d}%
				the reception of Sis. Overton, Prest. Smith It grieves me that there is no fuller fellowship--- if one member suffer all feel it--- by union of feeling we obtain pow'r with God. Christ said he came to call sinners to repentance and save them. Christ was condemn'd by the righteous jews because he took sinners into his society--- he took them <up>on the principle that they repented of their sins.
					\hypertarget{EMS-1842-JS-9-JUN-1842-e}%
				It
					\marginnote{e}%
				is the object of this Society to reform persons, not to take those that are corrupt, but if they repent we are bound to take them and by kindness sanctify and cleanse from all unrighteousness, by our influence in watching over them--- nothing will have such influence over people, as the fear of being disfellowship'd by so goodly a Society as this. Then take Sis. O. as Jesus received sinners into his bosom.

				Sis. O. In the name of the Lord I now make you free, and from this hour if any thing should be found against you

					\hypertarget{EMS-1842-JS-9-JUN-1842-f}%
				Nothing
					\marginnote{f}%
				is so much calculated to lead people to forsake sin as to take them by the hand and watch over them with tenderness. When persons manifest the least kindness and love to me, O what pow'r it has over my mind, while the opposite course has a tendency to harrow up all the harsh feelings and depress the human mind.

					\hypertarget{EMS-1842-JS-9-JUN-1842-g}%
				It
					\marginnote{g}%
				is one evidence that men are unacquainted with the principle of godliness, to behold the contraction of feeling and lack of charity. The pow'r and glory of Godliness is spread out on a broad principle to throw out the mantle of charity. God does not look on sin with allowance, but when men have sin'd there must be allowance made for them.

					\hypertarget{EMS-1842-JS-9-JUN-1842-h}%
				All
					\marginnote{h}%
				the religious world is boasting of righteousness--- tis the doctrine of the devil to retard the human mind and retard our progress, by filling us with selfrighteousness--- The nearer we get to our heavenly Father, the more are we dispos'd to look with compassion on perishing souls--- to take them upon our shoulders and cast their sins behind our back. [\textit{blank}] I am going to talk to all this Society--- if you would have God have mercy on you, have mercy on one another.

					\hypertarget{EMS-1842-JS-9-JUN-1842-i}%
				Prest.
					\marginnote{i}%
				S. then refer'd them to the conduct of the Savior when he was taken and crucified \&c.

				He then made a promise in the name of the Lord saying, that soul that has righteousness enough to ask God in the secret place for life, every day of their lives shall live to three score years \& ten--- We must walk uprightly all day long--- How glorious are the principles of righteousness! We are full of selfishness--- the devil flatters us that we are very righteous, while we are feeding on the faults of others--- We can only live by worshipping our God--- all must do it for themselves--- none can do it for another.
					\hypertarget{EMS-1842-JS-9-JUN-1842-j}%
				How
					\marginnote{j}%
				mild the Savior dealt with Peter, saying ``when thou art converted, strengthen thy brethren''--- at an other time he said to him ``lovest thou me? ``Feed my sheep''.--- If the sisters love the Lord let them feed the sheep and not destroy them. How oft have wise men \& women sought to dictate br. Joseph by saying ``O if I were br. Joseph I would do this and that.'' But if they were in br. Joseph's shoes, they would find that men could not be compel'd into the kingdom of God, but must be dealt with in long suff'ring--- and at last we shall save them.
					\hypertarget{EMS-1842-JS-9-JUN-1842-k}%
				The
					\marginnote{k}%
				way to keep all the saints together and keep the work rolling, is to wait with all long suff'ring till God shall bring such character to justice. There should be no license for sin, but mercy should go hand in hand with reproof.

					\hypertarget{EMS-1842-JS-9-JUN-1842-l}%
				Sisters
					\marginnote{l}%
				of this Society, shall there be strife among you? I will not have it--- you must repent and get the love of God. Away with selfrighteousness. The best measure or principle to bring the poor to repentance is to administer to their wants--- the Society is not only to relieve the poor, but to save souls.

					\hypertarget{EMS-1842-JS-9-JUN-1842-m}%
				Prest.
					\marginnote{m}%
				S. then said that he would give a lot of land to the Society by deeding it to the Treasurer, that the Society may build houses for the poor. He also said he would give a house--- frame not finished--- said that br.
				
		% -----
		\subsubsection{Letter to Parley P. Pratt and Others, 12 June 1842}
		% -----
			\label{EMS-1842-JS-12-JUN-1842}
			% \label{EMS-1842-JS-DAY-3LETTERMONTH-YEAR}
			
			\begin{description}
				\item[Print Sources]
			\end{description}

			\begin{tabular}{p{0.5in}p{1in}p{0.5in}p{1in}}
				JSP	 	& ---		& JSR 		& --- \\
				DC	 	& ---		& HC 		& --- \\
				HDDC 	& ---		& TCHC 		& --- \\
			\end{tabular}
			% HDDC = woodford's DC diss; JSR = Marquardt's JS Revelations; TCHC = Vogel HC
			
			\begin{description}
				\item[Online Access] \href{https://www.josephsmithpapers.org/paper-summary/letter-to-parley-p-pratt-and-others-12-june-1842/1}{Here}
				% \item[Analysis] \hyperref[XX]{Here}
			\end{description}
			
			\begin{description}
				\item[Background]
			\end{description}
			
				% It might also be useful to give a two sentence context of the period: e.g., "This speech was given in Nauvoo to a small congregation one month prior to the destruction of the Nauvoo Expositor. These and other tensions may have been on the prophet's mind when he gave the speech."
			
			\begin{description}
				\item[Original Source]
			\end{description}
			
				See notes in the RS volume.
			
				% brief description of original source material, with reference to JSP or somewhere else for more info
				% Background material: it might be helpful to have a note that says whether a given document was presented as a speech, was written in a journal, etc.
				% Also a comment here about how reliable the source might be, or what shape it is in, if there are large omissions or parts that are hard to understand, and so on.
			
			\begin{description}
				\item[Text]
			\end{description}
			
				\begin{tightright}
						\hypertarget{EMS-1842-JS-12-JUN-1842-a}%
					Nauvoo
						\marginnote{a}%
					June 12\supersu{th}\supers{.} 1842.
				\end{tightright}

				To our well-beloved <brother> Parley P. Pratt, and to the Elders of the Church of Jesus Christ of Latter Day Saints, in England and scattered abroad throughout all Europe, and to the saints, Greeting, Whereas in times past, persons have been permitted to gather with the Saints \sout{with} at Nauvoo in North America; such as husbands leaving their wives and children behind; also, such as wives leaving their husbands and children behind; and such as women leaving their husbands, and such as husbands leaving their wives who have no children, and some because their companions are unbelievers, All this kind of proceedings we consider [to] be an error and for the want of proper information and the same should be taught to all the Saints; and not suffer families to be broken up on no account whatever, if it is possible to avoid it.
					\hypertarget{EMS-1842-JS-12-JUN-1842-b}%
				Suffer
					\marginnote{b}%
				no man to leave his wife because she is an unbeliever; nor no woman to leave her husband, because he is an unbeliever. These things are an evil, and must be forbidden by the authorities of the church, or they will come under condemnation; for the gathering is not in haste, nor by flight; but to prepare all things before you, and you know not but the unbeliever may be converted, and the Lord heal him: but let the believers exercise faith in God, and the unbelieving husband, shall be sanctified by the believing wife; and the unbelieving wife, by the believing husband; and families are preserved, and saved from a great evil; which we have seen verified before our eyes.
					\hypertarget{EMS-1842-JS-12-JUN-1842-c}%
				Behold
					\marginnote{c}%
				this is a wicked generation, full of lyings, and deceit, and craftiness. and the children of [t]he wicked are wiser than the children of light, <i.e. they are more crafty> and it seems that it has been the case in all ages of the world, and the man, when he leaves his wife and travels to a foreign nation; while on his way, darkness overpowers his mind, and Satan deceives him, and flatters him with the graces of the harlot; and before he is aware he is disgraced forever: and greater is the danger, for the woman that leaves her husband. And their are several instances where women have left their husbands, and come to this place, \& in a few weeks, or months, they have found themselves new husbands, and they are living in adultery; and we are obliged to cut them off from the church.
					\hypertarget{EMS-1842-JS-12-JUN-1842-d}%
				\phantom{ }
					\marginnote{d}%
				\sout{I presume.} There are men also that are guilty of the same crime, as we are credibly informed. We are \uline{knowing} to their having taken wives \uline{here} and are \uline{credibly} informed that they have wives in England. There is another evil which does exist There are poor men who come here and leave their families behind in a destitute situation, and beg for assistan[c]e to send back after the[i]r families. Every man should tarry with his family untill Providence provides for the whole, for there is no means here to be obtained to send back. Money is scarce and hard to be obtained. The people that gather to this place are generally poor, the gathering being attended with a great sacrifice; and money cannot be obtained by labor; but all kinds of produce is plentiful, and can be obtained by labor.
					\hypertarget{EMS-1842-JS-12-JUN-1842-e}%
				Therefore,
					\marginnote{e}%
				the poor man that leaves his family in England, cannot get means---which must be silver and Gold---to send for his family, but must remain under the painful sensation, that his family must be cast upon the mercy of the people and separated and put into the Poorhouse. Therefore, to remedy the evil, we forbid a man's leaving his family behind, because he has no means to bring them. If the Church is not able to bring them, and the Parrish will not send them let the man tarry with his family---live with them---and die with them, and not leave them untill providence [s]hall open a way for the[m] to all come together, and we also forbid that a woman shall leave her husband, because he is an unbeliever.
					\hypertarget{EMS-1842-JS-12-JUN-1842-f}%
				We
					\marginnote{f}%
				also forbid that a man shall leave his wife, because she is an unbeliever. If he is a bad man (i.e. the unbeliever) there is a law to remedy that evil. and if she is a bad woman, there is a law to remedy that evil: And if the law will divorce them, then they are at liberty. Otherwise they are bound, as long as they two shall live, and it is not our prerogative to go beyond this; if we do it, it will be at the expense of our reputation.

				These are the things in plainness, which we desire should be publickly known; and you can publish them in the Millenial Star, in full, or extract, as you please.

					\hypertarget{EMS-1842-JS-12-JUN-1842-g}%
				It
					\marginnote{g}%
				is a general time of health in Nauvoo. Every thing begins to flouri[s]h and look prosperous. Crops of grain have the appea[ra]nce of a rich harvest. Em[i]gration continues to increase, so does also the City. We expect to see brother P. P. Pratt probably as soon as next spring. Brother Amos Fielding will be the bearer of this, he will start from here in a few days.

				May the Lord bestow his blessings upon you richly, and hasten the gathering, and bring abou[t] the fulness of the everlasting Covenant are the prayers of your brethren.

				Written by Hyrum Smith, Patriarch, by the order of Joseph Smith, president over the whole Church of Jesus Christ of Latter Day Sai[n]ts.

				\begin{tightright}
					\uline{Hyrum Smith}
				\end{tightright}

					\hypertarget{EMS-1842-JS-12-JUN-1842-h}%
				N.
					\marginnote{h}%
				B. Brother P. P. Pratt will send over 3 families namely brother John Allabys, John Farrars, \& David Claytons, \sout{with} <by> the donation m[o]ney that shall be given in for the building of the Temple[.] They are now at work on [t]he Temple, under th[at] sp[e]cial contract, that their families shall be forwarded to this place by monies donated for the Temple. Brother John Allabys family lives in No 33 Brownlow Hill, Liverpool. John Farrar \& David Claytons families live at Messrs Bashale and Boardmans Mill Farrington [Farington], near Preston, Lancashire. Direct Ann Farrar care of Mr Thomas Beardwood, Shopkeeper Messors Bashale \& Boardmans Mill. Also Elizabeth Clayton care of Mr Thomas Beardwood, Shopkeeper Messrs Bashale \& Boardmans Mill, Farrington, near Preston, Lancashire. Brother Amos Fielding will understand the particulars. This is a precedent, that we cannot establish, therefore you will be particular and keep this to yourself.

				\begin{tightright}
					Joseph Smith

					<Trustee in Trust>

					Hyrum Smith
				\end{tightright}

				<We wish to have these families sent, this fall if possible, or they must, suffer.--->

				\begin{tightright}
					<J. S.>

					<H. S.>
				\end{tightright}

					\hypertarget{EMS-1842-JS-12-JUN-1842-i}%
				We
					\marginnote{i}%
				assure you that you have our best feelings, and our prayers, and have no fault to find. Believing every man has done the best he could, that is--- the Elders, such as have remained in England. And we desire your prayers, even all the Saints--- \&c. \&c.

				\begin{tightright}
					H. S.

					J. S.
				\end{tightright}

				\begin{tightcenter}
					M\supersu{r}\supers{.} Parley P. Pratt

					Liverpool

					England

					\uline{Care of bro Amos Fielding}
				\end{tightcenter}
				
		% -----
		\subsubsection{Account of Meeting and Discourse, 18 June 1842}
		% -----
			\label{EMS-1842-JS-18-JUN-1842}
			% \label{EMS-1842-JS-DAY-3LETTERMONTH-YEAR}
			
			% --
			\paragraph{As Reported by Wilford Woodruff}
			% --
				\label{EMS-1842-JS-18-JUN-1842-WW}
			
				\begin{description}
					\item[Print Sources]
				\end{description}

				\begin{tabular}{p{0.5in}p{1in}p{0.5in}p{1in}}
					JSP	 	& ---		& JSR 		& --- \\
					DC	 	& ---		& HC 		& --- \\
					HDDC 	& ---		& TCHC 		& --- \\
				\end{tabular}
				% HDDC = woodford's DC diss; JSR = Marquardt's JS Revelations; TCHC = Vogel HC
				
				\begin{description}
					\item[Online Access] \href{https://www.josephsmithpapers.org/paper-summary/account-of-meeting-and-discourse-18-june-1842-as-reported-by-wilford-woodruff/1}{Here}
					% \item[Analysis] \hyperref[XX]{Here}
				\end{description}
				
				\begin{description}
					\item[Background]
				\end{description}
				
					% It might also be useful to give a two sentence context of the period: e.g., "This speech was given in Nauvoo to a small congregation one month prior to the destruction of the Nauvoo Expositor. These and other tensions may have been on the prophet's mind when he gave the speech."
				
				\begin{description}
					\item[Original Source]
				\end{description}
				
					% brief description of original source material, with reference to JSP or somewhere else for more info
					% Background material: it might be helpful to have a note that says whether a given document was presented as a speech, was written in a journal, etc.
					% Also a comment here about how reliable the source might be, or what shape it is in, if there are large omissions or parts that are hard to understand, and so on.
				
				\begin{description}
					\item[Text]
				\end{description}
				
						\hypertarget{EMS-1842-JS-18-JUN-1842-WW-a}%
					June
						\marginnote{a}%
					18\supersu{th} Saterday The Citizens of Nauvoo Both Male \& female assembled near the Temple for a general meeting many thousands were assembled Joseph the seer arose \& spoke upon several subjects.
					
					Among other subjects he spoke his mind in great plainness concerning the iniquity \& wickedness of Gen John Cook Bennet[t], \& exposed him before the public. He also prophesied in the Name of the Lord Concerning the merchants in the City that if they with the rich did not open their hearts \& contribute to the poor that they would be cursed by the hand of God \& be cut off from the land of the living,
						\hypertarget{EMS-1842-JS-18-JUN-1842-WW-b}%
					The
						\marginnote{b}%
					main part of the day was taken up upon the business of the agricultural \& manufacturing society \uline{ie} we have a Charter granted us by the Legislator of the state for that purpose \& the time has come for us to make use of that Charter it is divided into stock of \$50. dollars each share, any person owning one share became a member of the society a stockholder, each share is entitled to one vote this is esstablished with a view of helping the poor arangments were entered into to commence operations immediately
					
						\hypertarget{EMS-1842-JS-18-JUN-1842-WW-c}%
					Also
						\marginnote{c}%
					Joseph commmanded the Twelve to organize the Church more according to the Law of \uline{God} \uline{that} \uline{is} to require of those that come in to be setteled according to their council \& also to appoint a committee to wait upon all who arive \& make them welcome \& council them what to do B[righam] Young. H[eber] C. Kimball G[eorge] A. Smith \& Hyrum Smith was the committee appointed to wait upon emigrants \& settle them.
					
		% -----
		\subsubsection{Letter to Jennetta Richards Richards, 23 June 1842}
		% -----
			\label{EMS-1842-JS-23-JUN-1842}
			% \label{EMS-1842-JS-DAY-3LETTERMONTH-YEAR}
			
			\begin{description}
				\item[Print Sources]
			\end{description}

			\begin{tabular}{p{0.5in}p{1in}p{0.5in}p{1in}}
				JSP	 	& ---		& JSR 		& --- \\
				DC	 	& ---		& HC 		& --- \\
				HDDC 	& ---		& TCHC 		& --- \\
			\end{tabular}
			% HDDC = woodford's DC diss; JSR = Marquardt's JS Revelations; TCHC = Vogel HC
			
			\begin{description}
				\item[Online Access] \href{https://www.josephsmithpapers.org/paper-summary/letter-to-jennetta-richards-richards-23-june-1842/1}{Here}
				% \item[Analysis] \hyperref[XX]{Here}
			\end{description}
			
			\begin{description}
				\item[Background]
			\end{description}
			
				% It might also be useful to give a two sentence context of the period: e.g., "This speech was given in Nauvoo to a small congregation one month prior to the destruction of the Nauvoo Expositor. These and other tensions may have been on the prophet's mind when he gave the speech."
			
			\begin{description}
				\item[Original Source]
			\end{description}
			
				% brief description of original source material, with reference to JSP or somewhere else for more info
				% Background material: it might be helpful to have a note that says whether a given document was presented as a speech, was written in a journal, etc.
				% Also a comment here about how reliable the source might be, or what shape it is in, if there are large omissions or parts that are hard to understand, and so on.
			
			\begin{description}
				\item[Text]
			\end{description}
			
				\begin{tightright}
						\hypertarget{EMS-1842-JS-23-JUN-1842-a}%
					Nauvoo
						\marginnote{a}%
					June 23\supers{rd.} 1842
				\end{tightright}

				Sister Jennetta [Richards] Richards;

				Agreabley to your request, in the Midst of all the bustle, and buisness of the day, and the care of all the Churches boath at home and abroad. I now imbrace a moment to adress a few words to you thinking peradventure it may be a consolation to you to know that you too are remembered by me as well as all the saints. my hearts desire and prayr to God is all the day long for all the saints and in an especial and poticular [particular] manner for those whom he hath chosen and anointed to bear the heaviest burthens in the heat of the day among which number is your husband \sout{raised} received a man in whom I have the most implicit confidence and trust you say I have got him so I have in the which I rejoice, for he has done me great good and taken a great \sout{burthen} burden off my shoulders since his arrival in Nauvoo never did I have greater intimacy with any man than with him may the blessings of Elijah crown his head forever and ever.
					\hypertarget{EMS-1842-JS-23-JUN-1842-b}%
				we
					\marginnote{b}%
				are about to send him in a few days after his dear family he shall have our pray'rs fervently for his safe arrival to their imbraces and may God speed his Journey and return him quickly to our society. and I want you beloved Sister, to be a Genral in this matter, in helping him along. which I know you will he will be able to teach you many things which you never have heard you may have implicit confidence in the same. I have heard much about you by the twelve and in consequence of the great friendship that exists between your husband and me and the information they all have given me of your virtue and strong attachment to the truth of the work of God in the Last Days I have formed a very strong Brotherly friendship and attachment for you in the bonds of the Gosple,
					\hypertarget{EMS-1842-JS-23-JUN-1842-c}%
				Although
					\marginnote{c}%
				I never saw you I shall be exceedingly glad to see you face to face and be able to administer in the name of the Lord some of the words of Life to your consolation and I hope that you may be kept steadfast in the faith even unto the end, I want you should give my love and tender reguard to B\supersu{r} Richards [Willard Richards's] familey and those who are friendly enough to me to enquire after me: in that region of Country, not having but little time to apportion to anyone \& \sout{have} <having> stolen this oppertunity I therefore subscribe myself in haste your most Obedient Brother in the fulness of the Gosple Joseph Smith

				P.S. B\supers{ro} Richards having been with me for a long time can give you any information which you need and will tell you all about me. I shall be very anxious for his return he is a grate prop to me in my Labours.---

				Mrs. Jennetta Richards

				Richmond

				Massachusetts
				
		% -----
		\subsubsection{Letter to the Church, 23 June 1842}
		% -----
			\label{EMS-1842-JS-23-JUN-1842-C}
			% \label{EMS-1842-JS-DAY-3LETTERMONTH-YEAR}
			
			\begin{description}
				\item[Print Sources]
			\end{description}

			\begin{tabular}{p{0.5in}p{1in}p{0.5in}p{1in}}
				JSP	 	& ---		& JSR 		& --- \\
				DC	 	& ---		& HC 		& --- \\
				HDDC 	& ---		& TCHC 		& --- \\
			\end{tabular}
			% HDDC = woodford's DC diss; JSR = Marquardt's JS Revelations; TCHC = Vogel HC
			
			\begin{description}
				\item[Online Access] \href{https://www.josephsmithpapers.org/paper-summary/letter-to-the-church-23-june-1842/1}{Here}
				% \item[Analysis] \hyperref[XX]{Here}
			\end{description}
			
			\begin{description}
				\item[Background]
			\end{description}
			
				% It might also be useful to give a two sentence context of the period: e.g., "This speech was given in Nauvoo to a small congregation one month prior to the destruction of the Nauvoo Expositor. These and other tensions may have been on the prophet's mind when he gave the speech."
			
			\begin{description}
				\item[Original Source]
			\end{description}
			
				% brief description of original source material, with reference to JSP or somewhere else for more info
				% Background material: it might be helpful to have a note that says whether a given document was presented as a speech, was written in a journal, etc.
				% Also a comment here about how reliable the source might be, or what shape it is in, if there are large omissions or parts that are hard to understand, and so on.
			
			\begin{description}
				\item[Text]
			\end{description}
			
					\hypertarget{EMS-1842-JS-23-JUN-1842-C-a}%
				TO
					\marginnote{a}%
				THE CHURCH OF JESUS CHRIST OF LATTER DAY SAINTS, AND TO ALL THE HONORABLE PART OF COMMUNITY.

				It becomes my duty to lay before the Church of Jesus Christ of Latter Day Saints, and the public generally, some important facts relative to the conduct and character of Dr. \textsc{John C. Bennett}, who has lately been expelled from the aforesaid church; that the honorable part of community may be aware of his proceedings, and be ready to treat him and regard him as he ought to be regarded, viz: as an imposter and base adulterer.

					\hypertarget{EMS-1842-JS-23-JUN-1842-C-b}%
				It
					\marginnote{b}%
				is a matter of notoriety that said Dr. J. C. Bennett, became favorable to the doctrines taught by the elders of the church of Jesus Christ of Latter Day Saints, and located himself in the city of Nauvoo, about the month of August 1840, and soon after joined the church. Soon after it was known that he had become a member of said church, a communication was received at Nauvoo, from a person of respectable character, and residing in the vicinity where Bennett had lived.
					\hypertarget{EMS-1842-JS-23-JUN-1842-C-c}%
				This
					\marginnote{c}%
				letter cautioned us against him, setting forth that he was a very mean man, and had a wife, and two or three children in McConnelsville, Morgan county, Ohio; but knowing that it is no uncommon thing for good men to be evil spoken against, the above letter was kept quiet, but held in reserve.

					\hypertarget{EMS-1842-JS-23-JUN-1842-C-d}%
				He	
					\marginnote{d}%
				had not been long in Nauvoo before he began to keep company with a young lady, one of our citizens; and she being ignorant of his having a wife living, gave way to his addresses, and became confident, from his behavior towards her, that he intended to marry her; and this he gave her to understand he would do. I, seeing the folly of such an acquaintance, persuaded him to desist; and, on account of his continuing his course, finally threatened to expose him if he did not desist. This, to outward appearance, had the desired effect, and the acquaintance beneonate them was broken off.

					\hypertarget{EMS-1842-JS-23-JUN-1842-C-e}%
				But,
					\marginnote{e}%
				like one of the most abominable and depraved beings which could possibly exist, he only broke off his publicly wicked actions, to sink deeper into iniquity and hypocrisy. When he saw that I would not submit to any such conduct, he went to some of the females in the city, who knew nothing of him but as an honorable man, \& began to teach them that promiscous intercourse between the sexes, was a doctrine believed in by the Latter-Day Saints, and that there was no harm in it;
					\hypertarget{EMS-1842-JS-23-JUN-1842-C-f}%
				but
					\marginnote{f}%
				this failing, he had recourse to a more influential and desperately wicked course; and that was, to persuade them that myself and others of the authorities of the church not only sanctioned, but practiced the same wicked acts; and when asked why I publicly preached so much against it, said that it was because of the prejudice of the public, and that it would cause trouble in my own house. He was well aware of the consequence of such wilful and base falsehoods, if they should come to my knowledge;
					\hypertarget{EMS-1842-JS-23-JUN-1842-C-g}%
				and
					\marginnote{g}%
				consequently endeavored to persuade his dupes to keep it a matter of secresy, persuading them there would be no harm if they should not make it known. This proceeding on his part, answered the desired end; he accomplished his wicked purposes; he seduced an innocent female by his lying, and subjected her character to public disgrace, should it ever be known.

					\hypertarget{EMS-1842-JS-23-JUN-1842-C-h}%
				But
					\marginnote{h}%
				his depraved heart would not suffer him to stop here. Not being contented with having disgraced one female, he made an attempt upon others; and, by the same plausible tale, overcame them also; evidently not caring whose character was ruined, so that his wicked, lustful appetities might be gratified.

					\hypertarget{EMS-1842-JS-23-JUN-1842-C-i}%
				Sometime
					\marginnote{i}%
				about the early part of July 1841, I received a letter from Elder H[yrum] Smith and Wm. Law, who were then at Pittsburgh, Penn. This letter was dated June 15th, and contained the particulars of a conversation betwixt them and a respectable gentleman from the neighborhood where Bennett's wife and children resided. He stated to them that it was a fact that Bennett had a wife and children living, and that she had left him because of his ill-treatment towards her. This letter was read to Bennett, which he did not attempt to deny; but candidly acknowledged the fact.

					\hypertarget{EMS-1842-JS-23-JUN-1842-C-j}%
				Soon
					\marginnote{j}%
				after this information reached our ears, Dr. Bennett made an attempt at suicide, by taking poison; but he being discovered before it had taken effect, and the proper antidotes being administered, he again recovered; but he very much resisted when an attempt was made to save him. The public impression was, that he was so much ashamed of his base and wicked conduct, that he had recourse to the above deed to escape the censures of an indignant community.

					\hypertarget{EMS-1842-JS-23-JUN-1842-C-k}%
				It
					\marginnote{k}%
				might have been supposed that these circumstances transpiring in the manner they did, would have produced a thorough reformation in his conduct; but, alas! like a being totally destitute of common decency, and without any government over his passions, he was soon busily engaged in the same wicked career, and continued until a knowledge of the same reached my ears.
					\hypertarget{EMS-1842-JS-23-JUN-1842-C-l}%
				I
					\marginnote{L}%
				immediately charged him with it, and he admitted that it was true; but in order to put a stop to all such proceedings for the future, I publicly proclaimed against it, and had those females notified to appear before the proper officers that the whole subject might be investigated and thoroughly exposed.

					\hypertarget{EMS-1842-JS-23-JUN-1842-C-m}%
				During
					\marginnote{m}%
				the course of investigation, the foregoing facts were proved by credible witnesses, and were sworn and subscribed to before an alderman of the city, on the 15 ult. The documents containing the evidence are now in my possession.

					\hypertarget{EMS-1842-JS-23-JUN-1842-C-n}%
				We
					\marginnote{n}%
				also ascertained by the above investigation, that others had been led by his conduct to persue the same adulterous practice, and in order to accomplish their detestable designs made use of the same language insinuated by Bennett, with this difference, that they did not hear me say any thing of the kind, but Bennett was one of the heads of the church, and he had informed them that such was the fact, and they credited his testimony.

					\hypertarget{EMS-1842-JS-23-JUN-1842-C-o}%
				The
					\marginnote{o}%
				public will perceive the aggravating nature of this case; and will see the propriety of this exposure. Had he only been guilty of adultry, that was sufficient to stamp disgrace upon him because he is a man of better information, and has been held high in the estimation of many. But when it is considered that his mind was so intent upon his cruel, and abominable deeds, and his own reputation not being sufficient to enable him to do it, he must make use of my name in order to effect his purposes, an enlightened public will not be astonished at the course I have pursued.

					\hypertarget{EMS-1842-JS-23-JUN-1842-C-p}%
				In
					\marginnote{p}%
				order that it may be distinctly understood that he wilfully and knowingly lied, in the above insinuations, I will lay before my readers an affidavit taken before an alderman of the city, after I had charged him with these things:
				
				\begin{quote}
					\textsc{State of Illinois},)

					City of Nauvoo.)

					Personally appeared before me, Daniel H. Wells, an Alderman of said city of Nauvoo, John C. Bennett, who being duly sworn according to law, deposeth and saith: that he never was taught anything in the least contrary to the strictest principles of the Gospel, or of virtue, or of the laws of God, or man, under any circumstances, or upon any occasion either directly or indirectly, in word or deed, by Joseph Smith;
						\hypertarget{EMS-1842-JS-23-JUN-1842-C-q}%
					and
						\marginnote{q}%
					that he never knew the said Smith to countenance any improper conduct whatever, either in public or private; and that he never did teach to me in private that an illegal illicit intercourse with females was, under any circumstances, justifiable; and that I never knew him so to teach others.

					\begin{tightright}
						JOHN C. BENNETT.
					\end{tightright}

					Sworn to, and subscribed, before me, this 17th day of May, A. D. 1842.

					\begin{tightright}
						DANIEL H. WELLS, Alderman.
					\end{tightright}
				\end{quote}

					\hypertarget{EMS-1842-JS-23-JUN-1842-C-r}%
				The
					\marginnote{r}%
				following conversation took place in the City Council, and was elicited in consequence of its being reported that the Doctor had stated that I had acted in an indecorous manner, and given countenance to vices practised by the Doctor, and others:
				
				\begin{quote}

					\hypertarget{EMS-1842-JS-23-JUN-1842-C-s}%
				Dr.
					\marginnote{s}%
				John C. Bennett, ex-Mayor, was then called upon by the Mayor to state if he knew aught aganst him; when Mr. Bennett replied: ``I know what I am about, and the heads of the Church know what they are about. I expect I have no difficulty with the heads of the church. I publicly avow that any one who has said that I have stated that General Joseph Smith has given me authority to hold illicit intercourse with women is a liar in the face of God, those who have said it are damned liars; they are infernal liars.
					\hypertarget{EMS-1842-JS-23-JUN-1842-C-t}%
				He
					\marginnote{t}%
				never, either in public or private, gave me any such authority or license, and any person who states it is a scoundrel and a liar. I have heard it said that I should become a second [Sampson] Avard by withdrawing from the church, and that I was at variance with the heads and should use an influence against them because I resigned the office of Mayor; this is false.
					\hypertarget{EMS-1842-JS-23-JUN-1842-C-u}%
				I
					\marginnote{u}%
				have no difficulty with the heads of the church, and I intend to continue with you, and hope the time may come when I may be restored to full confidence, and fellowship, and my former standing in the church; and that my conduct may be such as to warrant my restoration---and should the time ever come that I may have the opportunity to test my faith it will then be known whether I am a traitor or a true man.''

					\hypertarget{EMS-1842-JS-23-JUN-1842-C-v}%
				Joseph
					\marginnote{v}%
				Smith then asked: ``Will you please state definitely whether you know any thing against my character either in public or private?''

				Gen. Bennett answered: ``I do not; in all my intercourse with Gen. Smith, in public and in private, he has been strictly virtuous.
				
				\begin{tabular}{p{6cm}p{6cm}}
				\textit{Aldermen.} & GEO. A. SMITH,				\\
				N[ewel] K. WHITNEY, & WILSON LAW,               \\
				HIRAM KIMBALL, & B[righam] YOUNG,               \\
				ORSON SPENCER, & JOHN TAYLOR,                   \\
				GUST[avus] HILLS, & H[eber] C. KIMBALL,         \\
				G[eorge] W. HARRIS, & W[ilford] WOODRUFF,       \\
				Counsellors. & JOHN P. GREEN[E],                \\
				WILLARD RICHARDS,	                            
				\end{tabular}                                   

				\begin{tightright}
					JAMES SLOAN, City Recorder.
				\end{tightright}

				May 19, 1842.

				\end{quote}

					\hypertarget{EMS-1842-JS-23-JUN-1842-C-w}%
				After
					\marginnote{w}%
				I had done all in my power to persuade him to amend his conduct, and these facts were fully established, (not only by testimony, but by his own concessions,) he having acknowledged that they were true, and seeing no prospects of any satisfaction from his future life, the hand of fellowship was withdrawn from him as a member of the church, by the officers; but on account of his earnestly requesting that we would not publish him to the world, we concluded not to do so at that time, but would let the matter rest until we saw the effect of what we had already done.

					\hypertarget{EMS-1842-JS-23-JUN-1842-C-x}%
				It
					\marginnote{x}%
				appears evident, that as soon as he perceived that he could no longer maintain his standing as a member of the church, nor his respectability as a citizen, he came to the conclusion to leave the place; which he has done; and that very abruptly; and had he done so quietly, and not attempted to deceive the people around him, his case would not have excited the indignation of the citizens, so much as his real conduct has done.
					\hypertarget{EMS-1842-JS-23-JUN-1842-C-y}%
				In
					\marginnote{y}%
				order to make his case look plausible, he has reported, ``that he had withdrawn from the church because we were not worthy of his society;'' thus instead of manifesting a spirit of repentance, he has to the last, proved himself to be unworthy the confidence or regard of any upright person, by lying, to deceive the innocent, and committing adultery in the most abominable and degraded manner.

					\hypertarget{EMS-1842-JS-23-JUN-1842-C-z}%
				We
					\marginnote{z}%
				are credibly informed that he has colleagued with some of our former wicked persecutors, the Missourians, and has threatened destruction upon us; but we should naturally suppose, that he would be so much ashamed of himself at the injury he has already done to those who never injured, but befriended him in every possible manner, that he could never dare to lift up his head before an enlightened public, with the design either to misrepresent or persecute;
					\hypertarget{EMS-1842-JS-23-JUN-1842-C-aa}%
				but
					\marginnote{aa}%
				be that as it may, we neither dread him nor his influence; but this much we believe, that unless he is determined to fill up the measure of his iniquity, and bring sudden destruction upon himself from the hand of the Almighty; he will be silent, and never more attempt to injure those concerning whom he has testified upon oath he knows nothing but that which is good and virtuous.

					\hypertarget{EMS-1842-JS-23-JUN-1842-C-bb}%
				Thus
					\marginnote{bb}%
				I have laid before the Church of Latter Day Saints, and before the public, the character and conduct of a man who has stood high in the estimation of many; but from the foregoing facts it will be seen that he is not entitled to any credit, but rather to be stamped with indignity and disgrace so far as he may be known.
					\hypertarget{EMS-1842-JS-23-JUN-1842-C-cc}%
				What
					\marginnote{cc}%
				I have stated I am prepared to prove, having all the documents concerning the matter in my possession, but I think that to say further is unnecessary, as the subject is so plain that no one can mistake the true nature of the case.

				\begin{tightcenter}
					I remain yours, respectfully,
				\end{tightcenter}

				\begin{tightright}
					JOSEPH SMITH.
				\end{tightright}

				Nauvoo, June 23, 1842.
				
		% -----
		\subsubsection{Letter to Thomas Carlin, 24 June 1842}
		% -----
			\label{EMS-1842-JS-24-JUN-1842}
			% \label{EMS-1842-JS-DAY-3LETTERMONTH-YEAR}
			
			\begin{description}
				\item[Print Sources]
			\end{description}

			\begin{tabular}{p{0.5in}p{1in}p{0.5in}p{1in}}
				JSP	 	& ---		& JSR 		& --- \\
				DC	 	& ---		& HC 		& --- \\
				HDDC 	& ---		& TCHC 		& --- \\
			\end{tabular}
			% HDDC = woodford's DC diss; JSR = Marquardt's JS Revelations; TCHC = Vogel HC
			
			\begin{description}
				\item[Online Access] \href{https://www.josephsmithpapers.org/paper-summary/letter-to-thomas-carlin-24-june-1842/1}{Here}
				% \item[Analysis] \hyperref[XX]{Here}
			\end{description}
			
			\begin{description}
				\item[Background]
			\end{description}
			
				% It might also be useful to give a two sentence context of the period: e.g., "This speech was given in Nauvoo to a small congregation one month prior to the destruction of the Nauvoo Expositor. These and other tensions may have been on the prophet's mind when he gave the speech."
			
			\begin{description}
				\item[Original Source]
			\end{description}
			
				% brief description of original source material, with reference to JSP or somewhere else for more info
				% Background material: it might be helpful to have a note that says whether a given document was presented as a speech, was written in a journal, etc.
				% Also a comment here about how reliable the source might be, or what shape it is in, if there are large omissions or parts that are hard to understand, and so on.
			
			\begin{description}
				\item[Text]
			\end{description}
			
				\begin{tightright}
						\hypertarget{EMS-1842-JS-24-JUN-1842-a}%
					Nauvoo
						\marginnote{a}%
					June 24\supersu{th}\supers{.} 1842
				\end{tightright}

				Thomas Carlin Governor of the State of Ill,

				D\supersu{r}\supers{.} Sir,

				It becomes my duty to lay before you some facts relative to the conduct of our Major, General John C. Bennett; which have been proven beyond the possibility of dispute, and which he himself has admitted to be true, in my presence.

					\hypertarget{EMS-1842-JS-24-JUN-1842-b}%
				It
					\marginnote{b}%
				is evident that his general character is that of an adulterer of the worst kind, and although he has a wife and Children living, circumstances which have transpired in Nauvoo, have proven to a demonstration that he cares not whose character is disgraced whose honor is destroyed nor who suffers so that his lustful appetite may be gratified and further he cares not how many, nor how abominable the falsehood he has to make use of to accomplish his wicked purposes, even should it be that he brings disgrace upon a whole community.

					\hypertarget{EMS-1842-JS-24-JUN-1842-c}%
				Some
					\marginnote{c}%
				time ago, it having been reported to me that some of the most aggravating cases of adultery had been committed upon some previously respectable females in our City, I took proper measures to ascertain the truth of the report, and was soon enabled to bring sufficient witnesses before proper Authority to establish the following facts,
					\hypertarget{EMS-1842-JS-24-JUN-1842-d}%
				More
					\marginnote{d}%
				than twelve months ago Bennett went to a Lady in the City and began to teach her that promiscuous intercourse between the sexes was lawful and no harm in it, and requested the privilege of gratifying his passions but she refused in the strongest terms saying that it was very wrong to do so, and it would bring a disgrace on the church Finding this argument ineffectual he told her that men in higher standing in the church than himself not only sanctioned but practised the same deeds,
					\hypertarget{EMS-1842-JS-24-JUN-1842-e}%
				and
					\marginnote{e}%
				in order to finish the controversy said and affirmed that I both taught and acted in the same manner, but publicly proclaimed against it in consequence of the prejudace of the people and fear of trouble in my own house. By this means he accomplished his designs, he seduced a respectable female with lying and subjected her to public infamy and disgrace.

					\hypertarget{EMS-1842-JS-24-JUN-1842-f}%
				Not
					\marginnote{f}%
				contented with what he had already done he made the attempt on others and by using the same language seduced them also.

				about the early part of July 1841 I received a letter from Pittsburgh Pa In it was contained information setting forth that said Bennett had a wife and two or three children then living. This I re[a]d to him and he acknowledged it was true
		
					\hypertarget{EMS-1842-JS-24-JUN-1842-g}%
				A
					\marginnote{g}%
				very short time after this he attempted to destroy himself by taking poison but being discovered before it had taken sufficient affect, and proper antidotes administered he again recovered.

				The impression made on the minds of the public by this event was; that he was so ashamed of his base conduct that he took this coursse to escape the censures of a justly indignant community.
					\hypertarget{EMS-1842-JS-24-JUN-1842-h}%
				It
					\marginnote{h}%
				might have been supposed that after this he would have broke off his adulterous proceedings but to the contrary the public consternation had scarcly ceased before he was again deeply involved in the same wicked proceedings, and continued untill a knowledge of the fact reached my ears. I immediately charged him with the whole circumstance and he candidly acknowledged the truth of the whole.
	
					\hypertarget{EMS-1842-JS-24-JUN-1842-i}%
				The
					\marginnote{i}%
				foregoing facts were established on oath before an alderman of the City.--- the affidavits are now in my possession.

				In order that the truth might be fully established I asked Bennett to testify before an alderman wether I had given him any cause for such aggravating conduct He testified that I never taught to him that illicit intercourses with females was under any circumstances justifiable neither did he ever hear me teach any thing but the strictest principles of righteousness and virtue. This affidavit is also in my possession.

					\hypertarget{EMS-1842-JS-24-JUN-1842-j}%
				I
					\marginnote{j}%
				have also a similar affidavit taken before the city council and signed by the members of the council.

				after these things transpired, and finding that I should resist all such wicked conduct, and knowing that he could no longer maintain himself as a respectable citizen he has seen fit to leave Nauvoo, and that very abruptly

				I have been credibly informed that he is colleaguing with some of our former cruel persecuters the missourians and that he is threatening destruction upon us; and under these circumstances I consider it my duty to give you information on the subject that a knowledge of his proceedings may be before you, in due season

					\hypertarget{EMS-1842-JS-24-JUN-1842-k}%
				It
					\marginnote{k}%
				can be proven by hundreds of witnesses that he is one of the basest of liars that his whole routine of proceedings whilst amongst us has been of the basest kind.

				He also stated here that he had resigned his commission as Major General to the Governor weither this be true or not I have no knowledge, I wish to be informed on the subject that we may know how to act in relasion to the Legion.

					\hypertarget{EMS-1842-JS-24-JUN-1842-l}%
				A.
					\marginnote{l}%
				short time ago I was told by a friend of mine (not a member of the Church) that some of the missourians were conspiring to come up to Nauvoo and Kidnap me, and not doubting but that it might be true I consulted with General Bennett upon the most proper course to be pursued, we concluded to write to you on the subject, and I requested him to do so. I understand he has wrote to you but I know not in what manner, and I should be very much pleased if you would write to me on reciept of this giving me the contents of his communication.

					\hypertarget{EMS-1842-JS-24-JUN-1842-m}%
				I
					\marginnote{m}%
				have also heard that yourself has entertained of late very unfavourable feelings towards us as a People, and espicially so with reguard to myself and that you have said I ought to be shot \&c. If this be true I should be pleased to know from your self the reason of such hostile feelings, for I know of no cause which can possibly exist that might produce such feelings in your breast,

					\hypertarget{EMS-1842-JS-24-JUN-1842-n}%
				It
					\marginnote{n}%
				is rumoured and strong evidence exists that Bennett and David [Kilbourne] and Edward Kilbourn[e] have posted Bills in Galena calling upon the people to hold meetings and have themselves in readiness at a moments warning to assemble and come here and mob us out of the place and try to Kidnap me we know not as to the truth of this report but we have conversed with some transient persons who had the report from a Gentleman who lately came from there and had seen those hand Bills posted in Galena---

					\hypertarget{EMS-1842-JS-24-JUN-1842-o}%
				In
					\marginnote{o}%
				case of a mob coming upon us I wish to be informed by the Governor what will be the best course for us to pursue, and how he wishes us to act in reguard to this matter

				\begin{tightright}
					Joseph Smith Le\supersu{u}\supers{.} General

					Nauvoo. Legion.
				\end{tightright}
				
		% -----
		\subsubsection{Letter to James Arlington Bennet, 30 June 1842}
		% -----
			\label{EMS-1842-JS-30-JUN-1842}
			% \label{EMS-1842-JS-DAY-3LETTERMONTH-YEAR}
			
			\begin{description}
				\item[Print Sources]
			\end{description}

			\begin{tabular}{p{0.5in}p{1in}p{0.5in}p{1in}}
				JSP	 	& ---		& JSR 		& --- \\
				DC	 	& ---		& HC 		& --- \\
				HDDC 	& ---		& TCHC 		& --- \\
			\end{tabular}
			% HDDC = woodford's DC diss; JSR = Marquardt's JS Revelations; TCHC = Vogel HC
			
			\begin{description}
				\item[Online Access] \href{https://www.josephsmithpapers.org/paper-summary/letter-to-james-arlington-bennet-30-june-1842/1}{Here}
				% \item[Analysis] \hyperref[XX]{Here}
			\end{description}
			
			\begin{description}
				\item[Background]
			\end{description}
			
				% It might also be useful to give a two sentence context of the period: e.g., "This speech was given in Nauvoo to a small congregation one month prior to the destruction of the Nauvoo Expositor. These and other tensions may have been on the prophet's mind when he gave the speech."
			
			\begin{description}
				\item[Original Source]
			\end{description}
			
				% brief description of original source material, with reference to JSP or somewhere else for more info
				% Background material: it might be helpful to have a note that says whether a given document was presented as a speech, was written in a journal, etc.
				% Also a comment here about how reliable the source might be, or what shape it is in, if there are large omissions or parts that are hard to understand, and so on.
			
			\begin{description}
				\item[Text]
			\end{description}
			
				\begin{tightright}
						\hypertarget{EMS-1842-JS-30-JUN-1842-a}%
					Nauvoo
						\marginnote{a}%
					June 30\supersu{th}\supers{.} 1842
				\end{tightright}

				Gen\supersu{l}\supers{.} James Arlington Bennett [Bennet] \sout{Esqr}

				Dr Sir / Permit me to introduce to your kindness and attention, Dr Willard Richards, my private secretary, and General Business Agent. Any attention or favors conferred upon him, I shall esteem in the same degree as though conferred upon myself.

					\hypertarget{EMS-1842-JS-30-JUN-1842-b}%
				We
					\marginnote{b}%
				have been under the necessity of publishing the character and conduct; of our major \uline{G}eneral Dr John Cook Bennett, a circumstance which has been to us truly unpleasant; but his conduct whilst amongst us at Nauvoo, and up to the present time, has been such as to make it a matter of duty, which I consider we owe to the public generally, and could not be dispensed with.
					\hypertarget{EMS-1842-JS-30-JUN-1842-c}%
				As
					\marginnote{c}%
				Dr Richards is conversant with all the facts relative to his case, it would be superfluous to state them in this brief communication. Any information you may desire on the subject; or concerning the Legion, the prog[r]ess of the City, prospects of business, or any other matter he will cheerfully give, and to greater satisfaction than could be done with pen.

					\hypertarget{EMS-1842-JS-30-JUN-1842-d}%
				I
					\marginnote{d}%
				should be highly pleased if you should see proper to open a correspondence with me as soon as convenient; and also to receive a visit from you whenever circumstances shall render it consistent.

				You will please accept of my kind thanks for the several instances of remembrance from your pen, both as poetic effusions, and in defence of the glorious principles of truth, and I pray my heavenly Father that his spirit may continue to rest upon you forever.

				\begin{tightcenter}
						\hypertarget{EMS-1842-JS-30-JUN-1842-e}%
					I
						\marginnote{e}%
					am Sir.--- Yours resp[ectfull]\supers{y}
				\end{tightcenter}

				\begin{tightright}
					Joseph Smith

					Mayor
				\end{tightright}

				Gen\supersu{l}\supers{.} Ja\supers{s} Arlington Bennett \sout{Esqr}

				[\textit{blank}]

				[\textit{blank}] 

				\begin{tightcenter}
					Gen\supersu{l}\supers{.} Ja\supersu{s}\supers{.} Arlington Bennett

					Arlington House

					L. I. [Long Island]

					New York
				\end{tightcenter}

				Per Dr Richards

				<From the Prophet at Nauvoo Illinois.>

				<Aug. 6th. left the 8 \uline{1842}>
				
		% -----
		\subsubsection{Letter to the Citizens of Hancock County, circa 2 July 1842}
		% -----
			\label{EMS-1842-JS-CIR-2-JUL-1842}
			% \label{EMS-1842-JS-DAY-3LETTERMONTH-YEAR}
			
			\begin{description}
				\item[Print Sources]
			\end{description}

			\begin{tabular}{p{0.5in}p{1in}p{0.5in}p{1in}}
				JSP	 	& ---		& JSR 		& --- \\
				DC	 	& ---		& HC 		& --- \\
				HDDC 	& ---		& TCHC 		& --- \\
			\end{tabular}
			% HDDC = woodford's DC diss; JSR = Marquardt's JS Revelations; TCHC = Vogel HC
			
			\begin{description}
				\item[Online Access] \href{https://www.josephsmithpapers.org/paper-summary/letter-to-the-citizens-of-hancock-county-circa-2-july-1842/1}{Here}
				% \item[Analysis] \hyperref[XX]{Here}
			\end{description}
			
			\begin{description}
				\item[Background]
			\end{description}
			
				% It might also be useful to give a two sentence context of the period: e.g., "This speech was given in Nauvoo to a small congregation one month prior to the destruction of the Nauvoo Expositor. These and other tensions may have been on the prophet's mind when he gave the speech."
			
			\begin{description}
				\item[Original Source]
			\end{description}
			
				% brief description of original source material, with reference to JSP or somewhere else for more info
				% Background material: it might be helpful to have a note that says whether a given document was presented as a speech, was written in a journal, etc.
				% Also a comment here about how reliable the source might be, or what shape it is in, if there are large omissions or parts that are hard to understand, and so on.
			
			\begin{description}
				\item[Text]
			\end{description}
			
				\begin{tightcenter}
						\hypertarget{EMS-1842-JS-CIR-2-JUL-1842-a}%
					For
						\marginnote{a}%
					the Wasp.

					TO THE CITIZENS OF HANCOCK CO.
				\end{tightcenter}

				As a people, the church of Jesus Christ of Latter-Day Saints are found ``more sinned against, than sinning.''---In political affairs we are ever ready to yield to our fellow citizens of [t]he county, equal participation in the selection of candidates for office---we have been disappointed in our hopes of being met with the same disposition on the part of some of the old citizens of the county---they indeed seem to manifest a spirit of intolerance and exclusion, incompatible with the liberal doctrines of true republicanism. 
					\hypertarget{EMS-1842-JS-CIR-2-JUL-1842-b}%
				At
					\marginnote{b}%
				the late Anti-Mormon convention, a complete set of candidates, pledged to a man to receive no support from, and to yield no quarters to, Mormons, are commended to \textit{all} the citizens of this county for their suffrages! As a portion of said citizens of Hancock we embrace the occasion to decline \textit{this ticket} for the want of \textit{reciprocity} in its term, [a]nd honesty and intelligence in the character of some of its candidates.
				
					\hypertarget{EMS-1842-JS-CIR-2-JUL-1842-c}%
				If
					\marginnote{c}%
				the old citizens of the county are still desirous of equal participations with us in the choice of candidates, we are ready to co-operate with them---If independent gentlemen will announce themselves, and possess the requisite qualities, \textit{capacity} and integrity, they \textit{will} receive the united support of our people in the county---The time for holding a convention seems to have already gone by---%
					\hypertarget{EMS-1842-JS-CIR-2-JUL-1842-d}%
				there
					\marginnote{d}%
				is time enough for the friends of justice and fair play to \textit{elect} a ticket, to be announced in the independent manner we have suggested. Let the gentlemen who have the \textit{courage} to oppose the spirit of dictation which governed the Anti-Mormon convention candidates, show themselves, and we will exercise enough, on the terms proposed in this article, to ensure complete success.
				
				\begin{tightright}
					JOSEPH SMITH.
				\end{tightright}
				
		% -----
		\subsubsection{Account of Meeting, 15 July 1842}
		% -----
			\label{EMS-1842-JS-15-JUL-1842}
			% \label{EMS-1842-JS-DAY-3LETTERMONTH-YEAR}
			
			\begin{description}
				\item[Print Sources]
			\end{description}

			\begin{tabular}{p{0.5in}p{1in}p{0.5in}p{1in}}
				JSP	 	& ---		& JSR 		& --- \\
				DC	 	& ---		& HC 		& --- \\
				HDDC 	& ---		& TCHC 		& --- \\
			\end{tabular}
			% HDDC = woodford's DC diss; JSR = Marquardt's JS Revelations; TCHC = Vogel HC
			
			\begin{description}
				\item[Online Access] \href{https://www.josephsmithpapers.org/paper-summary/account-of-meeting-15-july-1842/1}{Here}
				% \item[Analysis] \hyperref[XX]{Here}
			\end{description}
			
			\begin{description}
				\item[Background]
			\end{description}
			
				% It might also be useful to give a two sentence context of the period: e.g., "This speech was given in Nauvoo to a small congregation one month prior to the destruction of the Nauvoo Expositor. These and other tensions may have been on the prophet's mind when he gave the speech."
			
			\begin{description}
				\item[Original Source]
			\end{description}
			
				% brief description of original source material, with reference to JSP or somewhere else for more info
				% Background material: it might be helpful to have a note that says whether a given document was presented as a speech, was written in a journal, etc.
				% Also a comment here about how reliable the source might be, or what shape it is in, if there are large omissions or parts that are hard to understand, and so on.
			
			\begin{description}
				\item[Text]
			\end{description}
			
					\hypertarget{EMS-1842-JS-15-JUL-1842-a}%
				After
					\marginnote{a}%
				considerable search had been made but to no effect a meeting was called at the Grove where Joseph stated before the public a general outline of J[ohn] C. Bennetts conduct and especially with regard to Sis P [Sarah Marinda Bates Pratt] Met again in the P.M. when Hyrum [Smith] \& H[eber] C. Kimball spake on the same subject
					\hypertarget{EMS-1842-JS-15-JUL-1842-b}%
				after
					\marginnote{b}%
				which Joseph arose and said that he would state to those present some things which he had heard respecting Edward \& D Kilbourn [David Kilbourne] being conspiring with J. C. Bennett in endeavoring to bring a Mob upon us, and as Mr E. Kilbourn was then present he would have the privilege of either admitting or denying it.
					\hypertarget{EMS-1842-JS-15-JUL-1842-c}%
				Question
					\marginnote{c}%
				by E. Kilbourn ``Who did Bennett tell that I and my brother were conspiring to bring a mob upon you'' Answer by Joseph ``He told me and he told [\textit{blank}] Allred and Orson Pratts wife \& others''. Q by E Kilboun ``Where did he say we were going to bring a mob from''. Ans. by Joseph. ``From Galena''. 
					\hypertarget{EMS-1842-JS-15-JUL-1842-d}%
				Mr.
					\marginnote{d}%
				Kilbourn then arose and said, ``I was conversing with my brother this morning and he said he had never seen Bennett since he had us before him last year for conspiracy. I have only seen him twice since last fall, I saw <him> once then. I was going to Galena about 2 weeks ago. The Boat I was on stopped at the upper Landing place and I came ashore a little while. The first person I saw was Bennett; we entered into conversation, but there was no mention made of mobs. I have not seen him since.
					\hypertarget{EMS-1842-JS-15-JUL-1842-e}%
				I
					\marginnote{e}%
				always regarded Bennett the same as I regard you (Joseph) and thought you were pretty well matched. If any one says that I have conspired to bring a Mob upon you it is false''. The meeting was then peaceably dismissed.
				
		% -----
		\subsubsection{Editorial, 15 July 1842--A}
		% -----
			\label{EMS-1842-JS-15-JUL-1842-A}
			% \label{EMS-1842-JS-DAY-3LETTERMONTH-YEAR}
			
			\begin{description}
				\item[Print Sources]
			\end{description}

			\begin{tabular}{p{0.5in}p{1in}p{0.5in}p{1in}}
				JSP	 	& ---		& JSR 		& --- \\
				DC	 	& ---		& HC 		& --- \\
				HDDC 	& ---		& TCHC 		& --- \\
			\end{tabular}
			% HDDC = woodford's DC diss; JSR = Marquardt's JS Revelations; TCHC = Vogel HC
			
			\begin{description}
				\item[Online Access] \href{https://www.josephsmithpapers.org/paper-summary/editorial-15-july-1842-a/1}{Here}
				% \item[Analysis] \hyperref[XX]{Here}
			\end{description}
			
			\begin{description}
				\item[Background]
			\end{description}
			
				% It might also be useful to give a two sentence context of the period: e.g., "This speech was given in Nauvoo to a small congregation one month prior to the destruction of the Nauvoo Expositor. These and other tensions may have been on the prophet's mind when he gave the speech."
			
			\begin{description}
				\item[Original Source]
			\end{description}
			
				% brief description of original source material, with reference to JSP or somewhere else for more info
				% Background material: it might be helpful to have a note that says whether a given document was presented as a speech, was written in a journal, etc.
				% Also a comment here about how reliable the source might be, or what shape it is in, if there are large omissions or parts that are hard to understand, and so on.
			
			\begin{description}
				\item[Text]
			\end{description}
			
				\begin{tightcenter}
						\hypertarget{EMS-1842-JS-15-JUL-1842-A-a}%
					THE
						\marginnote{a}%
					GOVERNMENT OF GOD.
				\end{tightcenter}

				The government of the Almighty, has always been very dissimilar to the government of men; whether we refer to his religious government, or to the government of nations. The government of God has always tended to promote peace, unity, harmony, strength and happiness; while that of man has been productive of confusion, disorder, weakness and misery.
					\hypertarget{EMS-1842-JS-15-JUL-1842-A-b}%
				The
					\marginnote{b}%
				greatest acts of the mighty men have been to depopulate nations, and to overthrow kingdoms; and whilst they have exalted themselves and become glorious, it has been at the expense of the lives of the innocent---the blood of the oppressed---the moans of the widow, and the tears of the orphan. Egypt, Babylon, Greece, Persia, Carthage, Rome---each were raised to dignity amid the clash of arms, and the din of war;
					\hypertarget{EMS-1842-JS-15-JUL-1842-A-c}%
				and
					\marginnote{c}%
				whilst their triumphant leaders led forth their victorious armies to glory and victory, their ears were saluted with the groans of the dying, and the misery and distress of the human family;---before them the earth was a paradise, and behind them a desolate wilderness; their kingdoms were founded in carnage and bloodshed, and sustained by oppression, tyranny, and despotism.
					\hypertarget{EMS-1842-JS-15-JUL-1842-A-d}%
				The
					\marginnote{d}%
				designs of God, on the other hand, have been to promote the universal good, of the universal world;---to establish peace and good will among men;---to promote the principles of eternal truth;---to bring about a state of things that shall unite man to his fellow man---cause the world to ``beat their swords into plow-shares, and their spears into pruning-hooks''---make the nations of the earth dwell in peace; and to bring about the millenial glory---when ``the earth shall yield its increase, resume its paradisean glory, and become as the garden of the Lord.''

					\hypertarget{EMS-1842-JS-15-JUL-1842-A-e}%
				The
					\marginnote{e}%
				great and wise of ancient days have failed in all their attempts to promote eternal power, peace, and happiness. Their nations have crumbled to pieces; their thrones have been cast down in their turn; and their cities, and their mightiest works of art have been annihilated; or their dilapidated towers, or time worn monuments have left us but feeble traits of their former magnificence, and ancient grandeur. They proclaim as with a voice of thunder, those imperishable truths---that man's strength is weakness, his wisdom is folly, his glory is his shame.

					\hypertarget{EMS-1842-JS-15-JUL-1842-A-f}%
				Monarchical,
					\marginnote{f}%
				aristocratic, and republican forms of government, of their various kinds and grades, have in their turn been raised to dignity and prostrated in the dust. The plans of the greatest politicians, the wisest senators, and most profound statesmen have been exploded; and the proceedings of the greatest chieftains, the bravest generals, and the wisest kings have fallen to the ground. Nation has succeeded nation, and we have inherited nothing but their folly. History records their puerile plans, their short lived glory, their feeble intellect, and their ignoble deeds.

					\hypertarget{EMS-1842-JS-15-JUL-1842-A-g}%
				Have
					\marginnote{g}%
				we increased in knowledge or intelligence? where is there a man that can step forth and alter the destiny of nations, and promote the happiness of the world? Or where is there a kingdom or nation, that can promote the universal happiness of its own subjects, or even their general well being? Our nation, which possesses greater resources than any other, is rent from center to circumference, with party strife, political intrigue, and sectional interest; our counsellors are panic struck, our legislators are astonished, and our senators are confounded; our merchants are paralized, our tradesmen are disheartened, our mechanics out of employ, our farmers distressed, and our poor crying for bread.
					\hypertarget{EMS-1842-JS-15-JUL-1842-A-h}%
				Our
					\marginnote{h}%
				banks are broken, our credit ruined, and our states overwhelmed in debt;---yet we are, and have been in peace.--- What is the matter? Are we alone in this thing? Verily, no. With all our evils we are better situated than any other nations. Let Egypt, Turkey, Spain, France, Italy, Portugal, Germany, England, China, or any other nation speak, and tell the tale of their trouble---their perplexity, and distress, and we should find that their cup was full, and that they were preparing to drink the dregs of sorrow.
					\hypertarget{EMS-1842-JS-15-JUL-1842-A-i}%
				England,
					\marginnote{i}%
				that boasts of her literature, her science, commerce, \&c., has her hands reeking with the blood of the innocent, abroad; and she is saluted with the cries of the oppressed, at home.---Chartism, O'Connelism, and Radicalism are gnawing her vitals at home; and Ireland, Scotland, Canada, and the East, are threatening her destruction abroad. France is rent to the core---intrigue, treachery, and treason lurk in the dark; and murder, and assassination stalk forth at noonday.
					\hypertarget{EMS-1842-JS-15-JUL-1842-A-j}%
				Turkey,
					\marginnote{j}%
				once the glory of European nations, has been shorn of her strength---has dwindled into her dotage, and has been obliged to ask her allies to propose to her tributary terms of peace: and Russia, and Egypt are each of [t]hem opening their jaws to devour her. Spain has been the theatre of bloodshed, of misery and woe, for years past. Syria is now convulsed with war and bloodshed.
					\hypertarget{EMS-1842-JS-15-JUL-1842-A-k}%
				The
					\marginnote{k}%
				great and powerful empire of China, which has for centuries resisted the attacks of barbarians, has become tributary to a foreign foe; her batteries thrown down; many of her cities destroyed, and her villages deserted. We might mention the Eastern rajahs; the miseries and oppressions of the Irish; the convulsed state of Central America; the situation of Texas and Mexico; the state of Greece, Switzerland, and Poland---nay, the world itself presents one great theatre of misery, woe, and ``distress of nations with perplexity.'' All, all speak with a voice of thunder, that man is not able to govern himself---to legislate for himself---to protect himself---to promote his own good, nor the good of the world.

					\hypertarget{EMS-1842-JS-15-JUL-1842-A-l}%
				It
					\marginnote{l}%
				has been the design of Jehovah, from the commencement of the world, and is his purpose now, to regulate the affairs of the world in his own time; to stand as head of the universe, and take the reigns of government into his own hand. When that is done judgment will be administered in righteousness; anarchy and confusion will be destroyed, and ``nations will learn war no more.'' It is for want of this great governing principle that all this confusion has existed; ``for it is not in man that walketh to direct his steps;'' this we have fully shewn.

					\hypertarget{EMS-1842-JS-15-JUL-1842-A-m}%
				If
					\marginnote{m}%
				there was any thing great or good in the world it came from God. The construction of the first vessel was given to Noah, by revelation. The design of the ark was given by God ``a pattern of heavenly things.'' The learning of the Egyptians, and their knowledge of astronomony was no doubt taught them by Abraham and Joseph, as their records testify, who received it from the Lord. The art of working in brass, silver, gold, and precious stones, was taught by revelation, in the wilderness.
					\hypertarget{EMS-1842-JS-15-JUL-1842-A-n}%
				The
					\marginnote{n}%
				architectural designs of the Temple at Jerusalem, together with its ornament and beauty was given of God. Wisdom to govern the house of Israel was given to Solomon, and to the judges of Israel; and if he had always been their king, and they subject to his mandate, and obedient to his laws, they would still have been a great and mighty people; the rulers of the universe---and the wonder of the world. If Nebuchadnezzar, or Darius, or Cyrus, or any other king possessed knowledge or power it was from the same source, as the scriptures abundantly testify.
					\hypertarget{EMS-1842-JS-15-JUL-1842-A-o}%
				If
					\marginnote{o}%
				then, God puts up one, and sets down another, at his pleasure---and made instruments of kings, unknown to themselves, to fulfill his prophesies, how much more was he able, if man would have been subject to his mandate, to regulate the affairs of this world, and promote peace and happiness among the human family.

					\hypertarget{EMS-1842-JS-15-JUL-1842-A-p}%
				The
					\marginnote{p}%
				Lord has at various times commenced this kind of government, and tendered his services to the human family. He selected Enoch, whom he directed, and gave his law unto, and to the people who were with him; and when the world in general would not obey the commands of God, after walking with God, he translated Enoch and his church, and the priesthood or government of heaven, was taken away.

					\hypertarget{EMS-1842-JS-15-JUL-1842-A-q}%
				Abraham
					\marginnote{q}%
				was guided in all his family affairs by the Lord; was told where to go, and when to stop; was conversed with by angels, and by the Lord; and prospered exceedingly in all that he put his hand unto; it was because he and his family obeyed the counsel of the Lord.--- When Egypt was under the superintendence of Joseph, it prospered, because he was taught of God; when they oppressed the Israelites destruction came upon them.
					\hypertarget{EMS-1842-JS-15-JUL-1842-A-r}%
				When
					\marginnote{r}%
				the children of Israel were chosen with Moses at their head, they were to be a peculiar people, among whom God should place his name: their motto was ``The Lord is our lawgiver; the Lord is our judge; the Lord is our king, and he shall reign over us.'' While in this state they might truly say ``happy is that people whose God is the Lord.'' Their government was a theocracy; they had God to make their laws, and men chosen by him to administer them; he was their God, and they were his people.
					\hypertarget{EMS-1842-JS-15-JUL-1842-A-s}%
				Moses
					\marginnote{s}%
				received the word of the Lord from God himself; he was the mouth of God to Aaron, and Aaron taught the people in both civil and ecclasiastical affairs; they were both one; there was no distinction; so will it be when the purposes of God shall be accomplished; when ``the Lord shall be king over the whole earth'' and ``Jerusalem his throne.'' ``The law shall go forth from Zion and the word of the Lord from Jerusalem.''

					\hypertarget{EMS-1842-JS-15-JUL-1842-A-t}%
				This
					\marginnote{t}%
				is the only thing that can bring about the ``restitution of all things, spoken of by all the holy prophets since the world was''---``the dispensation of the fulness of times, when GOD shall gather together all things in one.'' Other attempts to promote universal peace and happiness in the human family have proven abortive; every effort has failed; every plan and design has fallen to the ground; it needs the wisdom of God, the intelligence of God, and the power of God to accomplish this.
					\hypertarget{EMS-1842-JS-15-JUL-1842-A-u}%
				The
					\marginnote{u}%
				world has had a fair trial for six thousand years; the Lord will try the seventh thousand himself; ``he whose right it is will possess the kingdom, and reign until he has put all things under his feet;'' iniquity will hide its hoary head, Satan will be bound, and the works of darkness destroyed; righteousness will be put to the line, and judgment to the plummet, and ``he that fears the Lord will alone be exalted in that day.'' To bring about this state of things there must of necessity be great confusion among the nations of the earth; distress of nations with perplexity.''--- Am I asked what is the cause of the present distress? I would answer: ``Shall there be evil in a city and the Lord hath not done it.''
					\hypertarget{EMS-1842-JS-15-JUL-1842-A-v}%
				The
					\marginnote{v}%
				earth is groaning under corruption, oppression, tyranny, and bloodshed; and God is coming out of his hiding place, as he said that he would do, to vex the nations of the earth. Daniel, in his vision, saw convulsion upon convulsion; he ``saw till thrones were cast down, and the ancient of days did sit; and one was brought before him like unto the Son of man; and all nations, kindreds, tongues, and people, did serve and obey him.'' It is for us to be righteous that we may be wise and understand, for ``none of the wicked shall understand; but the wise shall understand, and they that turn many to righteousness, as the stars for ever and ever.''
					\hypertarget{EMS-1842-JS-15-JUL-1842-A-w}%
				As
					\marginnote{w}%
				a church, and a people it beho[o]ves us to be wise, and to seek to know the will of God, and then be willing to \textit{do it}; for ``blessed is he that heareth the word of the Lord and keepeth it,'' says the scriptures. ``Watch and pray always,'' says our Savior, ``that ye may be accounted worthy to escape the things that are coming on the earth, and to stand before the Son of man.'' If Enoch, Abraham, Moses, the children of Israel, and all God's people were saved by keeping the commandments of God, we, if saved at all, shall be saved upon the same principle.
					\hypertarget{EMS-1842-JS-15-JUL-1842-A-x}%
				As
					\marginnote{x}%
				God governed Abraham, Isaac and Jacob, as families, and the children of Israel as a nation, so we, as a church, must be under his guidance if we are prospered, preserved, and sustained. Our only confidence can be in God; our only wisdom obtained from him; and he alone must be our protector and safeguard, spiritually and temporally, or we fall.

					\hypertarget{EMS-1842-JS-15-JUL-1842-A-y}%
				We
					\marginnote{y}%
				have been chastened by the hand of God heretofore for not obeying his commands, although we never violated any human law, or transgressed any human precept: yet we have treated lightly his commands, and departed from his ordinances, and the Lord has chastened us sore, and we have felt his arm, and kissed the rod: let us we [be] wise in time to come, and ever remember that ``to obey is better than sacrifice; and to hearken than the fat of rams.''
					\hypertarget{EMS-1842-JS-15-JUL-1842-A-z}%
				The
					\marginnote{z}%
				Lord has told us to build the temple, and the Nauvoo House, and that command is as binding upon us as any other; and that man who engages not in these things is as much a transgressor as though he broke any other command---he is not a doer of God's will, nor a fulfiller of his laws.

					\hypertarget{EMS-1842-JS-15-JUL-1842-A-aa}%
				In
					\marginnote{aa}%
				regard to the building up of Zion it has to be done by the counsel of Jehovah; by the revelations of heaven, and we should feel to say ``if the Lord go not with us, carry us not up hence.'' We would say to the saints that come here, we have laid the foundation for the gathering of God's people to this place, and expect that when the saints do come they will be under the counsel of those that God has appointed.
					\hypertarget{EMS-1842-JS-15-JUL-1842-A-bb}%
				The
					\marginnote{bb}%
				Twelve are set apart to counsel the saints pertaining to this matter: and we expect that those who come here will send before them their wise men according to revelation; or if not practicable, be subject to the counsel that God has given or they cannot receive an inheritance among the saints, or be considered as God's people; and they will be dealt with as transgressors of the laws of God; we are trying here to gird up our loins, and purge from our midst the workers of iniquity;
					\hypertarget{EMS-1842-JS-15-JUL-1842-A-cc}%
				and
					\marginnote{cc}%
				we hope that when our brethren arrive from abroad, they will assist us to roll forth this good work, and to accomplish this great design; that ``Zion may be built up in righteousness; and all nations flock to her stan[d]ard;'' that as God's people, under his direction, and obedient to his law, we may grow up in righteousness, and truth; that when his purposes shall be accomplished, we may receive an inheritance among those that are sanctified.---\textsc{Ed}.
				
		% -----
		\subsubsection{Editorial, 15 July 1842--B}
		% -----
			\label{EMS-1842-JS-15-JUL-1842-B}
			% \label{EMS-1842-JS-DAY-3LETTERMONTH-YEAR}
			
			\begin{description}
				\item[Print Sources]
			\end{description}

			\begin{tabular}{p{0.5in}p{1in}p{0.5in}p{1in}}
				JSP	 	& ---		& JSR 		& --- \\
				DC	 	& ---		& HC 		& --- \\
				HDDC 	& ---		& TCHC 		& --- \\
			\end{tabular}
			% HDDC = woodford's DC diss; JSR = Marquardt's JS Revelations; TCHC = Vogel HC
			
			\begin{description}
				\item[Online Access] \href{https://www.josephsmithpapers.org/paper-summary/editorial-15-july-1842-b/1}{Here}
				% \item[Analysis] \hyperref[XX]{Here}
			\end{description}
			
			\begin{description}
				\item[Background]
			\end{description}
			
				% It might also be useful to give a two sentence context of the period: e.g., "This speech was given in Nauvoo to a small congregation one month prior to the destruction of the Nauvoo Expositor. These and other tensions may have been on the prophet's mind when he gave the speech."
			
			\begin{description}
				\item[Original Source]
			\end{description}
			
				% brief description of original source material, with reference to JSP or somewhere else for more info
				% Background material: it might be helpful to have a note that says whether a given document was presented as a speech, was written in a journal, etc.
				% Also a comment here about how reliable the source might be, or what shape it is in, if there are large omissions or parts that are hard to understand, and so on.
			
			\begin{description}
				\item[Text]
			\end{description}
			
				\begin{tightcenter}
						\hypertarget{EMS-1842-JS-15-JUL-1842-B-a}%
					AMERICAN
						\marginnote{a}%
					ANTIQUITIES.
				\end{tightcenter}

				Some have supposed that all the great works of the west, of which we have been treating, belong to our present race of Indians; but from continued wars with each other, have driven themselves from agricultural pursuits, and thinned away their numbers, to that degree, that the wild animals and fishes of the rivers, and wild fruit of the forests, were found sufficient to give them abundant support: on which account, they were reduced to savagism.

					\hypertarget{EMS-1842-JS-15-JUL-1842-B-b}%
				But
					\marginnote{b}%
				this is answered by the Antiquarian Society, as follows: ``Have our present race of Indians ever buried their dead in mounds by thousands? Were they acquainted with the uses of silver or copper? These metals curiously wrought have been found. Did the ancients of our Indians burn the bodies of distinguished chiefs, on funeral piles, and then raise a lofty tumulus over the urn containing their ashes? Did the Indians erect any thing like the ``walled towns,'' on Paint Creek? Did they ever dig such wells as are found at Marietta, Portsmouth, and above all, such as those in Paint Creek? Did they manufacture vessels from calcareous breccia, equal to any now made in Italy?

					\hypertarget{EMS-1842-JS-15-JUL-1842-B-c}%
				To
					\marginnote{c}%
				this we respond, they never have: no, not even their traditions afford a glimpse of the existence of such things, as forts, tumuli, roads, wells, mounds, walls enclosing between one and two hundred, and even five hundred acres of land; some of them of stone, and others of earth, twenty feet in thickness, and exceeding high, are works requiring too much labor for Indians ever to have performed.

					\hypertarget{EMS-1842-JS-15-JUL-1842-B-d}%
				An
					\marginnote{d}%
				idol found in a tumulus near Nashville, Tennessee, and now in the Museum of Mr. Clifford, of Lexington, is made of clay, peculiar for its fineness. With this clay was mixed a small portion of gypsum or plaster of Paris. This Idol was made to represent a man, in a state of nudity or nakedness, whose arms had been cut off close to the body, and whose nose and chin have been mutilated, with a fillet and cake upon its head.

					\hypertarget{EMS-1842-JS-15-JUL-1842-B-e}%
				Some
					\marginnote{e}%
				years since a clay vessel was discovered, about twenty feet below the surface, in alluvial earth, in digging a well near Nashville, Tennessee, and was found standing on a rock, from whence a spring of water issued. This vessel was taken to Peale's Museum, at Philadelphia. It contains about one gallon; was circular in its shape, with a flat bottom, from which it rises in a somewhat globose form, terminating at the summit with the figure of a female head; the place where the water was introduced, or poured out, was on the one side of it, nearly at the top of the globose part.

					\hypertarget{EMS-1842-JS-15-JUL-1842-B-f}%
				Another
					\marginnote{f}%
				idol was, a few years since, dug up in Natchez, on the Mississippi, on a piece of ground where, according to tradition, long before Europeans visited this country, stood an Indian temple.--- This idol is of stone, and is nineteen inches in height, nine inches in width, and seven inches thich at the extremities.--- On its breast, as represented on the plate of the idol, were five marks, which were evidently characters of some kind, resembling as supposed, the Persian; probably expressing, in the language of its authors, the name and supposed attributes of the senseless god of stone.

					\hypertarget{EMS-1842-JS-15-JUL-1842-B-g}%
				One
					\marginnote{g}%
				of the arts known to the builders of Babel, was that of brick making; this art was also known to the people who built the works in the west. The knowledge of copper was known to the people of the plains of Shinar, for Noah must have communicated it, as he lived an hundred and fifty years among them after the flood; also, copper was known to the antediluvians. Copper was also known to the authors of the western monuments.
					\hypertarget{EMS-1842-JS-15-JUL-1842-B-h}%
				Iron
					\marginnote{h}%
				was known to the antediluvians; it was also known to the ancients of the west; however, it is evident that very little iron was among them, as very few instances of its discovery in their works have occurred; and for this very reason we draw a conclusion that they came to this country very soon after the dispersion, and brought with them such few articles of iron as have been found in their works in an oxydized state.

					\hypertarget{EMS-1842-JS-15-JUL-1842-B-i}%
				Copper
					\marginnote{i}%
				ore is very abundant in many places of the west; and therefore, as they had a knowledge of it, when they first came here they knew how to work it, and form it into tools and ornaments. This is the reason why so many articles of this metal are found in their works; and even if they had a knowledge of iron ore, and knew how to work it, all articles made of it must have become oxydized as appears from what few specimens have been found, while those of copper are more imperishable. Gold ornaments are said to have been found in several tumuli. Silver very well plated on copper, has been found in several mounds, besides those at Circleville and Marietta. An ornament of copper was found in a stone mound near Chilicothe; it was a bracelet for the ancle or wrist.

					\hypertarget{EMS-1842-JS-15-JUL-1842-B-j}%
				The
					\marginnote{j}%
				ancients of Asia, immediately after the dispersion, were acquainted with ornaments made of the various metals; for the family of \textit{Terah}, who was the father of \textit{Abraham} and \textit{Nahor}, we find these ornaments in use for the beautifying of females. See the servant of Abraham, at the well of Bethuel in the country of ``Ur of the Chaldeans,'' or Mesopotamia, which is not very far from the place where Babel stood---putting a jewell of gold upon the face or forehead of Rebecca, weighing half a shekel, and two bracelets for her wrists, or arms.
					\hypertarget{EMS-1842-JS-15-JUL-1842-B-k}%
				Bracelets
					\marginnote{k}%
				for the same use have been found in the west; all of which circumstances go to establish the acquaintance of those who made those ornaments of silver and copper found in the mounds of the west, equal with those of Ur in Chaldea. The families of Peleg, Reu, Serug, and Nahor, who were the immediate progenitors of Abraham, lived at an era but little after the flood; and yet we find them in the possession of ornaments of this kind; from which we conclude a knowledge both of the metals, and how to make ornaments, as above described, was brought by Noah and his family from beyond the flood.

					\hypertarget{EMS-1842-JS-15-JUL-1842-B-l}%
				On
					\marginnote{l}%
				the shores of the Mississippi, some miles below Lake Pepin, on a fine plain, exists an artificial elevation of about four feet high, extending a full mile, in somewhat of a circular form. It is sufficiently capacious to have covered 5000 men. Every angle of the breast work is yet traceable, though much defaced by time. Here, it is likely, conflicting realms as great as those of the ancient Greeks and Persians, decided the fate of ambitious Monarchs, of the Chinese, Mongol descent.

					\hypertarget{EMS-1842-JS-15-JUL-1842-B-m}%
				Weapons
					\marginnote{m}%
				of brass have been found in many parts of America, as in the Canadas, Florida, \&c., with curiously sculptured stones, all of which go to prove that this country was once peopled with civilized, industrious nations,---now traversed the greater part by savage hunters.--- Priests American Antiquities.

				The Book of Mormon speaks of ores, swords, cities, armies, \&c., and we extract the following:---

				\begin{quote}
						\hypertarget{EMS-1842-JS-15-JUL-1842-B-n}%
					And
						\marginnote{o}%
					it came to pass that we did find upon the land of promise, as we journeyed in the wilderness, that there were beasts in the forests of every kind, both the cow and the ox, and the ass, and the horse, and the goat, and the wild goat, and all manner of wild animals, which were for the use of men. And we did find all manner of ore, both of gold, and of silver, and of copper.

					And it came to pass that the Lord commanded me, wherefore I did make plates of ore, that I might engraven upon them the record of my people. * * *

						\hypertarget{EMS-1842-JS-15-JUL-1842-B-o}%
					And
						\marginnote{o}%
					it came to pass that we began to prosper exceedingly, and to multiply in the land. And I, Nephi, did take the sword of Laban, and after the manner of it did make many swords, lest by any means the people who were now called Lamanites, should come upon us and destroy us: for I knew their hatred towards me and my children, and those who were called my people. And I did teach my people to build buildings; and to work in all manner of wood, and of iron, and of copper, and of brass, and of steel, and of gold, and of silver, and of precious ores, which were in great abundance. And I, Nephi, did build a temple; and I did construct it after the manner of the temple of Solomon, save it were not built of so many precious things: for they were not to be found upon the land; wherefore it could not be built like unto Solomon's temple. But the manner of the construction was like unto the temple of Solomon; and the workmenship thereof was exceeding fine.
				\end{quote}

					\hypertarget{EMS-1842-JS-15-JUL-1842-B-p}%
				In
					\marginnote{p}%
				regard to there being great wars, the following will shew:---

				\begin{quote}
					And it came to pass when Coriantumr had recovered of his wounds, be began to remember the words which Ether had spoken unto him...he saw that there had been slain by the sword already nearly two millions of his people, and he began to sorrow in his heart; yea, there had been slain two millions of mighty men, and also their wives and their children. He began to repent of the evil which he had done; he began to remember the words which had been spoken by the mouth of all the prophets, and he saw them that they were fulfilled, thus far, every whit; and his soul mourned, and refused to be comforted.....

						\hypertarget{EMS-1842-JS-15-JUL-1842-B-q}%
					And
						\marginnote{q}%
					it came to pass that they did gather together all the people, upon all the face of the land, who had not been slain, save it was Ether. And it came to pass that Ether did behold all the doings of the people; and he beheld that the people who were for Coriantumr, were gathered together for the army of Coriantumr; and the people who were for Shiz, were gathered together to the army of Shiz; wherefore they were for the space of four years gathering together the people, that they might get all who were upon the face of the land, and that they might receive all the strength which it was profitable that they could receive. And it came to pass that when they were all gathering together, every one to the army which he would with their wives and their children; both men, women, and children being armed with weapons of war, having shields and breast plates, and head plates, and being clothed after the manner of war, they did march forth one against another, to battle; and they fought all that day, and conquered not. And it came to pass that when it was night they were weary, and retired to their camps; and after they had retired to their camps, they took up a howling and a lamentation for the loss of the slain of their people; and so great were their cries, their howlings and lamentations, that it did rend the air exceedingly.
				\end{quote}

					\hypertarget{EMS-1842-JS-15-JUL-1842-B-r}%
				If
					\marginnote{r}%
				men, in their researches into the history of this country, in noticing the mounds, fortifications, statues, architecture, implements of war, of husbandry, and ornaments of silver, brass, \&c.---were to examine the Book of Mormon, their conjectures would be removed, and their opinions altered; uncertainty and doubt would be changed into certainty and facts; and they would find that those things that they are anxiously prying into were matters of history, unfolded in that book. They would find their conjectures were more than realized---that a great and a mighty people had inhabited this continent---that the arts sciences and religion, had prevailed to a very great extent, and that there was as great and mighty cities on this continent as on the continent of Asia.
					\hypertarget{EMS-1842-JS-15-JUL-1842-B-s}%
				Babylon,
					\marginnote{s}%
				Ninevah, nor any of the ruins of the Levant could boast of more perfect sculpture, better architectural designs, and more imperishable ruins, than what are found on this continent. Stephens and Catherwood's researches in Central America abundantly testify of this thing. The stupendous ruins, the elegant sculpture, and the magnificence of the ruins of Guatamala, and other cities, corroborate this statement, and show that a great and mighty people---men of great minds, clear intellect, bright genius, and comprehensive designs inhabited this continent. Their ruins speak of their greatness; the Book of Mormon unfolds their history.---\textsc{Ed}.
				
		% -----
		\subsubsection{Minutes, 22 July 1842, as Published by Times and Seasons}
		% -----
			\label{EMS-1842-JS-22-JUL-1842}
			% \label{EMS-1842-JS-DAY-3LETTERMONTH-YEAR}
			
			\begin{description}
				\item[Print Sources]
			\end{description}

			\begin{tabular}{p{0.5in}p{1in}p{0.5in}p{1in}}
				JSP	 	& ---		& JSR 		& --- \\
				DC	 	& ---		& HC 		& --- \\
				HDDC 	& ---		& TCHC 		& --- \\
			\end{tabular}
			% HDDC = woodford's DC diss; JSR = Marquardt's JS Revelations; TCHC = Vogel HC
			
			\begin{description}
				\item[Online Access] \href{https://www.josephsmithpapers.org/paper-summary/minutes-22-july-1842-as-published-by-times-and-seasons/1}{Here}
				% \item[Analysis] \hyperref[XX]{Here}
			\end{description}
			
			\begin{description}
				\item[Background]
			\end{description}
			
				% It might also be useful to give a two sentence context of the period: e.g., "This speech was given in Nauvoo to a small congregation one month prior to the destruction of the Nauvoo Expositor. These and other tensions may have been on the prophet's mind when he gave the speech."
			
			\begin{description}
				\item[Original Source]
			\end{description}
			
				% brief description of original source material, with reference to JSP or somewhere else for more info
				% Background material: it might be helpful to have a note that says whether a given document was presented as a speech, was written in a journal, etc.
				% Also a comment here about how reliable the source might be, or what shape it is in, if there are large omissions or parts that are hard to understand, and so on.
			
			\begin{description}
				\item[Text]
			\end{description}
			
					\hypertarget{EMS-1842-JS-22-JUL-1842-a}%
				At
					\marginnote{a}%
				a meeting of the citizens of the city of Nauvoo held in said city at the meeting ground, July 22d 1842.

				Orson Spencer Esq. was called to the chair, and Gustavus Hills was appointed clerk.

				The meeting was called to order by the chairman, who stated the object of the meeting to be to obtain an expression of the public mind in reference to the reports gone abroad, calumniating the character of Pres. Joseph Smith. Gen. Wilson Law then rose and presented the following resolution.

					\hypertarget{EMS-1842-JS-22-JUL-1842-b}%
				\phantom{ }
					\marginnote{b}%
				\textit{Resolved}---That, having heard that John C. Bennett was circulating many base falsehoods respecting a number of the citizens of Nauvoo, and especially against our worthy and respected Mayor, Joseph Smith, we do hereby manifest to the world that so far as we are acquainted with Joseph Smith we know him to be a good, moral, virtuous, peaceable and patriotic man, and a firm supporter of law, justice and equal rights; that he at all times upholds and keeps inviolate the constitution of this State and of the United States.

				A vote was then called and the resolution adopted by a large concourse of citizens, numbering somewhere about a thousand men. Two or three, voted in the negative.

					\hypertarget{EMS-1842-JS-22-JUL-1842-c}%
				Elder
					\marginnote{c}%
				Orson Pratt then rose and spoke at some length in explanation of his negative vote. Pres. Joseph Smith spoke in reply---

				Question to Elder Pratt, `Have you personally a knowledge of any immoral act in me toward the female sex, or in any other way?' Answer, by Elder O. Pratt, `Personally, toward the female sex, I have not.'

				Elder O. Pratt responded at some length. Elder B[righam] Young then spoke in reply, and was followed by Elders Wm. Law H[eber] C. Kimball and Pres. H[yrum] Smith. Several others spoke bearing testimony of the iniquity of those who had calumniated Pres. J. Smith's character.

				Meeting adjourned for one hour...
				
		% -----
		\subsubsection{Revelation, 27 July 1842}
		% -----
			\label{EMS-1842-JS-27-JUL-1842}
			% \label{EMS-1842-JS-DAY-3LETTERMONTH-YEAR}
			
			\begin{description}
				\item[Print Sources]
			\end{description}

			\begin{tabular}{p{0.5in}p{1in}p{0.5in}p{1in}}
				JSP	 	& ---		& JSR 		& --- \\
				DC	 	& ---		& HC 		& --- \\
				HDDC 	& ---		& TCHC 		& --- \\
			\end{tabular}
			% HDDC = woodford's DC diss; JSR = Marquardt's JS Revelations; TCHC = Vogel HC
			
			\begin{description}
				\item[Online Access] \href{https://www.josephsmithpapers.org/paper-summary/revelation-27-july-1842/1}{Here}
				% \item[Analysis] \hyperref[XX]{Here}
			\end{description}
			
			\begin{description}
				\item[Background]
			\end{description}
			
				% It might also be useful to give a two sentence context of the period: e.g., "This speech was given in Nauvoo to a small congregation one month prior to the destruction of the Nauvoo Expositor. These and other tensions may have been on the prophet's mind when he gave the speech."
			
			\begin{description}
				\item[Original Source]
			\end{description}
			
				There are at least two other old sources for this revelation; see the JSP historical introduction (there are links to them).
			
				% brief description of original source material, with reference to JSP or somewhere else for more info
				% Background material: it might be helpful to have a note that says whether a given document was presented as a speech, was written in a journal, etc.
				% Also a comment here about how reliable the source might be, or what shape it is in, if there are large omissions or parts that are hard to understand, and so on.
			
			\begin{description}
				\item[Text]
			\end{description}
			
				\begin{tightright}
						\hypertarget{EMS-1842-JS-27-JUL-1842-a}%
					Wednesday
						\marginnote{a}%
					27th July 1842.
				\end{tightright}
					
				Verily, thus saith the Lord unto my servant N[ewel] K. Whitney, the thing that my servant Joseph Smith has made known unto you and your family and which you have agreed upon, is right in mine eyes, and shall be crowned upon your heads with honor and immortality and eternal life to all your house both old and young.
					\hypertarget{EMS-1842-JS-27-JUL-1842-b}%
				Because
					\marginnote{b}%
				of the lineage of my Priesthood, saith the Lord, it shall be upon you and upon your children after you from generation to generation, by virtue of the holy promise which I now make unto you, saith the Lord. These are the words which you shall pronounce upon my servant Joseph and your daughter S. A. [Sarah Ann] Whitney:
					\hypertarget{EMS-1842-JS-27-JUL-1842-c}%
				They
					\marginnote{c}%
				shall take each other by the hand, and you shall say, You both mutually agree (calling them by name) to be each other's companion so long as you both shall live, preserving yourselves for each other and from all others, and also throughout eternity, reserving only those rights which have been given to my servant Joseph by revelation and commandment and by legal authority in times past.
					\hypertarget{EMS-1842-JS-27-JUL-1842-d}%
				If
					\marginnote{d}%
				you both agree to covenant and \sout{to} <do> this, I then give you, S. A. Whitney, my daughter, to Joseph Smith, to be his wife, to observe all the rights between you both that belong to that condition.
					\hypertarget{EMS-1842-JS-27-JUL-1842-e}%
				I
					\marginnote{e}%
				do it in my own name and in the name of my wife, your mother, and in the name of my holy progenitors by the right of birth, which is of Priesthood vested in me by revelation, and commandment, and promise of the living God, obtained by the Holy Mechisedek, Jethro, and others of the holy fathers, commanding, in the name of the Lord, all those powers to concentrate in you, and through you to your posterity forever.
					\hypertarget{EMS-1842-JS-27-JUL-1842-f}%
				All
					\marginnote{f}%
				these things I do in the name of the Lord Jesus Christ, that through this order He may be glorified, and that through the power of anointing, David may reign King over Israel, which shall hereafter be revealed. Let immortality and eternal life henceforth be sealed upon your heads forever and ever. [\textit{1/2 page blank}]
				
		% -----
		\subsubsection{Letter to Wilson Law, 14 August 1842}
		% -----
			\label{EMS-1842-JS-14-AUG-1842}
			% \label{EMS-1842-JS-DAY-3LETTERMONTH-YEAR}
			
			\begin{description}
				\item[Print Sources]
			\end{description}

			\begin{tabular}{p{0.5in}p{1in}p{0.5in}p{1in}}
				JSP	 	& ---		& JSR 		& --- \\
				DC	 	& ---		& HC 		& --- \\
				HDDC 	& ---		& TCHC 		& --- \\
			\end{tabular}
			% HDDC = woodford's DC diss; JSR = Marquardt's JS Revelations; TCHC = Vogel HC
			
			\begin{description}
				\item[Online Access] \href{https://www.josephsmithpapers.org/paper-summary/letter-to-wilson-law-14-august-1842/1}{Here}
				% \item[Analysis] \hyperref[XX]{Here}
			\end{description}
			
			\begin{description}
				\item[Background]
			\end{description}
			
				% It might also be useful to give a two sentence context of the period: e.g., "This speech was given in Nauvoo to a small congregation one month prior to the destruction of the Nauvoo Expositor. These and other tensions may have been on the prophet's mind when he gave the speech."
			
			\begin{description}
				\item[Original Source]
			\end{description}
			
				% brief description of original source material, with reference to JSP or somewhere else for more info
				% Background material: it might be helpful to have a note that says whether a given document was presented as a speech, was written in a journal, etc.
				% Also a comment here about how reliable the source might be, or what shape it is in, if there are large omissions or parts that are hard to understand, and so on.
			
			\begin{description}
				\item[Text]
			\end{description}
			
				\begin{tightright}
						\hypertarget{EMS-1842-JS-14-AUG-1842-a}%
					Head
						\marginnote{a}%
					Quarters of Nauvoo Legion

					Aug\supers{t.} 15 [14]--- 1842
				\end{tightright}

				Major Gen. Law

				Dr General---

				I take this opportunity to give you some instructions how I wish you to act in case our persecutors should carry their pursuits so far as to tread upon our rights as free-born American Citizens. The orders which I am about to give you is the result of a long series of contemplation since I saw you.--- I have come fully to the conclusion both since this last difficulty commenced, as before, that I never would suffer myself to go into the hands of the Missourians alive; and to go into the hands of the Officers of this state is nothing more nor less, than to go into the hands of the Missourians; for the whole farce has been gotten up, unlawfully and unconstitutionally, as well on the part of the Governor as others; by a mob spirit for the purpose of carrying out mob violence, to carry on mob tolerance in a religious persecution.
					\hypertarget{EMS-1842-JS-14-AUG-1842-b}%
				I
					\marginnote{b}%
				am determined therefore, to keep out of their hands, and thwart their designs if possible, that perhaps they may not urge the necessity of force and blood against their own fellow-citizens and loyal subjects; and become ashamed and withdraw their pursuits. But if they should not do this and shall urge the necessity of force; and \sout{I} if I by any means should be taken, these are therefore to command you forthwith, without delay, regardless of life or death to rescue me out of their hands. And further, to treat any pretensions to the contrary, unlawful and unconstitutional and as a mob got up for the purpose as religious persecution to take away the rights of men.
					\hypertarget{EMS-1842-JS-14-AUG-1842-c}%
				And
					\marginnote{c}%
				further, that our chartered rights and privileges shall be considered by us as holding the supremacy in the premises and shall be maintained; nothing short of the supreme court of this State, having authority to dis-annul them; and the Municipal Court having jurisdiction in my case. You will see therefore that the peace of the City of Nauvoo is kept, let who will endeavor to disturb it. You will see also that whenever any mob force or violence is used, on any citizen thereof, or that belongeth thereunto, you will see that that force or violence is immediately dispersed and brought to punishment; or meet it, and contest it at the point of the sword with firm and undaunted and unyeilding valor; and let them know that the spirit of \sout{of} old seventysix; and of George Washington yet lives, and is contained in the bosoms and blood of the children of the fathers thereof.
					\hypertarget{EMS-1842-JS-14-AUG-1842-d}%
				If
					\marginnote{d}%
				there are any threats in the city let legal steps be taken on the part of those that make the threats: and let no man, woman or child be intimidated nor suffer it to be done. Nevertheless as I said in the first place we will take every measure that lays in our power and make every sacrifice that God or man could require at our hands to preserve the peace and safety of the people without colition [collision?]. And if sacrificing my own liberty for months and years without stooping to the disgrace of Missouri persecution and violence, and [Thomas] Carlins mis-rule and corruption I bow to my fate with cheerfulness and all due deference in the consideration of the lives, safety and welfare of others.
					\hypertarget{EMS-1842-JS-14-AUG-1842-e}%
				But
					\marginnote{e}%
				if this policy cannot accomplish the desired object; let our charter, and our Municipality; free trade and Sailors rights be our motto, and go a-head David Crockett like, and lay down our lives like men, and defend ourselves to the best advantage we can to the very last. You are therefore, hereby authorised and commanded, by virtue of the authority which I hold, and commission granted me by the Executive of this State, to maintain the very letter and spirit of the above contents of this letter to the very best of your ability; to the extent of our lives, and our fortunes; and to the lives and the fortunes of the [Nauvoo] Legion; as also all those who may volunteer their lives and their fortunes with ours; for the defence of our wives, our children, our fathers and our mothers; our homes; our grave yards, and our tombs; and our dead and their tombstones, and our dear bought American liberties with the blood of our fathers, and all that is dear and sacred to man. Shall we shrink at the onset?
					\hypertarget{EMS-1842-JS-14-AUG-1842-f}%
				No,
					\marginnote{f}%
				let every mans brow be as the face of a Lion; let his heart be unshaken as the mighty oak, and his knee confirmed as the sappline [sapling] of the forrest; and by the voice and loud roar of the cannon; and the loud peals and thundering of Artillery; and by the voice of the thunderings of heaven as upon mount Sinai; and by the voice of the Heavenly Hosts; and by the voice of the Eternal God; and by the voice of innocent blood; and by the voice of innocence; and by the voice of all that is sacred and dear to man, let us plead the justice of our cause; trusting in the arm of Jehovah the Eloheem who sits enthroned in the heavens:
					\hypertarget{EMS-1842-JS-14-AUG-1842-g}%
				that
					\marginnote{g}%
				peradventure he may give us the victory; and if we bleed we shall bleed in a good cause--- in the cause of innocence and truth: and from henceforth will their not be a crown of glory for us? And will not those who come after us, hold our names in sacred remembrance? and will our enemies dare to brand us with cowardly reproach? With these considerations, I subscribe myself Yours most faithfully and respectfully with acknowledgements of your high and honored trusts as Major Gen. of the Nauvoo Legion

				\begin{tightright}
					Joseph Smith--- Mayor of the City of Nauvoo

					and Lieu\supersu{t} Gen. of the Nauvoo Legion of Illinois. Militia
				\end{tightright}

				Wilson Law.

					\hypertarget{EMS-1842-JS-14-AUG-1842-h}%
				\phantom{ }
					\marginnote{h}%
				\uline{Major Gen. of the Nauvoo Legion}.---

				P.S. I want you to communicate all the information to me of all the transactions, as they are going on daily, in writing by the hand of my aidecamps. As I am not willing that any thing that goes from my hands to you should be made a public matter, I enjoin upon you to keep all things in your own bosom; and I want every thing that comes from you to come through my Aids. The bearer of this will be able to pilot them in a way that will not be prejudicial to my safety---

				\begin{tightright}
					Joseph Smith.---
				\end{tightright}
				
		% -----
		\subsubsection{Letter to Emma Smith, 16 August 1842}
		% -----
			\label{EMS-1842-JS-16-AUG-1842}
			% \label{EMS-1842-JS-DAY-3LETTERMONTH-YEAR}
			
			\begin{description}
				\item[Print Sources]
			\end{description}

			\begin{tabular}{p{0.5in}p{1in}p{0.5in}p{1in}}
				JSP	 	& ---		& JSR 		& --- \\
				DC	 	& ---		& HC 		& --- \\
				HDDC 	& ---		& TCHC 		& --- \\
			\end{tabular}
			% HDDC = woodford's DC diss; JSR = Marquardt's JS Revelations; TCHC = Vogel HC
			
			\begin{description}
				\item[Online Access] \href{https://www.josephsmithpapers.org/paper-summary/letter-to-emma-smith-16-august-1842/1}{Here}
				% \item[Analysis] \hyperref[XX]{Here}
			\end{description}
			
			\begin{description}
				\item[Background]
			\end{description}
			
				% It might also be useful to give a two sentence context of the period: e.g., "This speech was given in Nauvoo to a small congregation one month prior to the destruction of the Nauvoo Expositor. These and other tensions may have been on the prophet's mind when he gave the speech."
			
			\begin{description}
				\item[Original Source]
			\end{description}
			
				% brief description of original source material, with reference to JSP or somewhere else for more info
				% Background material: it might be helpful to have a note that says whether a given document was presented as a speech, was written in a journal, etc.
				% Also a comment here about how reliable the source might be, or what shape it is in, if there are large omissions or parts that are hard to understand, and so on.
			
			\begin{description}
				\item[Text]
			\end{description}
			
				\begin{tightright}
						\hypertarget{EMS-1842-JS-16-AUG-1842-a}%
					Nauvoo
						\marginnote{a}%
					August 16\supersu{th}\supers{.} 1842
				\end{tightright}

				My Dear Emma

				I embrace this opportunity to express to you some of my feelings this morning. First of all, I take the liberty to tender you my sincere thanks for the two interesting and consoling visits that you have made me during my almost exiled situation. Tongue can not express the gratitude of my heart, for the warm and true-hearted friendship you have manifested in these things toward me. The time has passed away since you left me, very agreeably; thus far, my mind being perfectly reconciled to my fate, let it be what it may. I have been kept from melancholy and dumps, by the kind-heartedness of brother [Erastus] Derby, and his interesting chit-chat from time to time, which has called my mind from the more strong contemplations of things, and subjects that would have preyed more earnestly upon my feelings. Last night---in the night---brother Hyrum [Smith], [George] Miller, [William] Law \& others came to see us.
					\hypertarget{EMS-1842-JS-16-AUG-1842-b}%
				They
					\marginnote{b}%
				seemed much agitated, and expressed some fears in consequence of some manouverings and some flying reports which they had heard in relation to our safety; but after relating what it was, I was able to comprehend the whole matter to my entire satisfaction, and did not feel at all alarmed or uneasy. They think, however, that the Militia will be called out to search the city, and if this should be the case I would be much safer for the time being at a little distance off, untill Governor [Thomas] Carlin could get weary and be made ashamed of his corrupt and unhallowed pro-ceedings.
					\hypertarget{EMS-1842-JS-16-AUG-1842-c}%
				I
					\marginnote{c}%
				had supposed, however, that if there were any serious operations taking by the Governor; that Judge [James] Ralston or brother [David S.] Hollister would have notified us; and cannot believe that any thing very serious is to be apprehended, untill we obtain information from a source that can be relied on. I have consulted wether it is best for you to go to Quincy, and see the Governor; but on the whole, he is a fool; and the impressions that are suggested to my mind, are, that it will be of no use; and the more we notice him, and flatter him, the more eager he will be for our destruction. You may write to him, whatever you see proper, but to go and see him, I do not give my consent at present.
					\hypertarget{EMS-1842-JS-16-AUG-1842-d}%
				Brother
					\marginnote{d}%
				Miller again suggested to me the propriety of my accompanying him to the Pine woods, and then he return, and bring you and the children. My mind will eternally revolt at every suggestion of that kind. More especially since the dream and vision that was manifested to me on the last night. My safety is with you, if you want to have it so. Any thing more or less than this cometh of evil. My feelings and council I think ought to be abided. If I go to the Pine country, you shall go along with me, and the children; and if you and the children go not with me, I dont go.
					\hypertarget{EMS-1842-JS-16-AUG-1842-e}%
				I
					\marginnote{e}%
				do not wish to exile myself for the sake of my own life, I would rather fight it out. It is for your sakes, therefore, that I would do such a thing. I will go with you then, in the same carriage and on Horse back, from time to time, as occasion may require; for I am not willing to trust you, in the hands of those who cannot feel the same interest for you, that I feel; to be subject to the caprice, temptations, or notions of any-body whatever.
					\hypertarget{EMS-1842-JS-16-AUG-1842-f}%
				And
					\marginnote{f}%
				I must say that I am pre-possessed somewhat, with the notion of going to the Pine Country any how; for I am tired of the mean, low, and unhallowed vulgarity, of some portions of the society in which we live; and I think if I could have a respite of about six months with my family, it would be a savor of life unto life, with my house. Nevertheless if it were possible I would like to live here in peace and wind up my business; but if it should be ascertained to a dead certainty that there is no other remedy, then we will round up our shoulders and cheerfully endure it; and this will be the plan. Let my horse, saddle, saddle-bags, and valice to put some shirts and clothing in, be sent to me.
					\hypertarget{EMS-1842-JS-16-AUG-1842-g}%
				Let
					\marginnote{g}%
				brother Derby and Miller take a horse and put it into my Buggy with a trunk containing my heavier cloths, shoes and Boots \&c and let brother [John] Taylor accompany us to his fathers, and there we will tarry, taking every precaution to keep out of the hands of the enemy, untill you can arrive with the children. Let brother Hyrum bring you. Let Lorain [Lorin Walker] and brother [William] Clayton come along and bring all the writings and papers, books and histories, for we shall want a scribe in order that we may pour upon the world the truth like the Lava of Mount Vesuvius.
					\hypertarget{EMS-1842-JS-16-AUG-1842-h}%
				Then,
					\marginnote{h}%
				let all the goods, household furniture, cloths and Store Goods that can be procured be put on to the Boat, and let 20 or 30 of the best men that we can find be put on board to man it, and let them meet us at Prairie Du Chien; and from thence, we will wend our way like larks up the Mississippi untill the touring [towering?] mountains and rocks, shall reminds us of the places of our nativity, and shall look like safety and home; and then we will bid defiance to the world, to Carlin, [Lilburn W.] Boggs, [John C.] Bennett, and all their whorish whores, and motly clan, that follow in their wake, Missouri not excepted; and until the damnation of hell rolls upon them, by the voice, and dread thunders, and trump of the eternal God; then, in that day will we not shout in the victory and be crowned with eternal joys, for the battles we have fought, having kept the faith and overcome the world.
					\hypertarget{EMS-1842-JS-16-AUG-1842-i}%
				Tell
					\marginnote{i}%
				the children that it is well with their father, as yet; and that he remains in fervent prayer to Almighty God for the safety of himself, and for you, and for them. Tell Mother [Lucy Mack] Smith that it shall be well with her son, wether in life or in death; for thus saith the Lord God. Tell her that I remember her all the while, as well as Lucy [Smith Millikin] and all the rest; they all must be of good cheer. Tell Hyrum to be sure and not fail to carry out my instructions, but at the same time if the Militia does not come, and we should get any favorable information all may be well yet. Yours in haste, Your affectionate husband untill death, through all eternity forevermore

				\begin{tightright}
					Joseph Smith
				\end{tightright}

					\hypertarget{EMS-1842-JS-16-AUG-1842-j}%
				P.S.
					\marginnote{j}%
				I want you to write to Lorenzo [D.] Wasson, and get him to make affidavit to all he knows about Bennett and forward it. I also want you to ascertain from Hyrum wether he will conform to what I have requested. And you must write me an answer per bearer, giving me all the news you have, and what is the appearance of things this morning

				\begin{tightright}
					J.S.---
				\end{tightright}
				
		% -----
		\subsubsection{Letter to Wilson Law, 16 August 1842}
		% -----
			\label{EMS-1842-JS-16-AUG-1842-W}
			% \label{EMS-1842-JS-DAY-3LETTERMONTH-YEAR}
			
			\begin{description}
				\item[Print Sources]
			\end{description}

			\begin{tabular}{p{0.5in}p{1in}p{0.5in}p{1in}}
				JSP	 	& ---		& JSR 		& --- \\
				DC	 	& ---		& HC 		& --- \\
				HDDC 	& ---		& TCHC 		& --- \\
			\end{tabular}
			% HDDC = woodford's DC diss; JSR = Marquardt's JS Revelations; TCHC = Vogel HC
			
			\begin{description}
				\item[Online Access] \href{https://www.josephsmithpapers.org/paper-summary/letter-to-wilson-law-16-august-1842/1}{Here}
				% \item[Analysis] \hyperref[XX]{Here}
			\end{description}
			
			\begin{description}
				\item[Background]
			\end{description}
			
				% It might also be useful to give a two sentence context of the period: e.g., "This speech was given in Nauvoo to a small congregation one month prior to the destruction of the Nauvoo Expositor. These and other tensions may have been on the prophet's mind when he gave the speech."
			
			\begin{description}
				\item[Original Source]
			\end{description}
			
				% brief description of original source material, with reference to JSP or somewhere else for more info
				% Background material: it might be helpful to have a note that says whether a given document was presented as a speech, was written in a journal, etc.
				% Also a comment here about how reliable the source might be, or what shape it is in, if there are large omissions or parts that are hard to understand, and so on.
			
			\begin{description}
				\item[Text]
			\end{description}
			
					\hypertarget{EMS-1842-JS-16-AUG-1842-W-a}%
				Head
					\marginnote{a}%
				Quarters of Nauvoo Legion

				August 16 1842

				Major Gen. [Wilson] Law.

				Beloved brother and friend

				Those few lines which I received from you written on the 15\supersu{th}\supers{.} was to me like apples of Gold in pictures of Silver. I rejoice with exceeding great joy to be associated in the high and responsible stations which we hold, whose mind and feelings and heart is so congenial with my own.
					\hypertarget{EMS-1842-JS-16-AUG-1842-W-b}%
				I
					\marginnote{b}%
				love that soul that is so nobly entabernacled in that clay of yours, may God Almighty grant, that it may be satiated with seeing a fuffilment of every virtuous \sout{desire} and manly desire that you possess. May we be able to triumph gloriously over those who seek our destruction and overthrow, which I believe we shall.
					\hypertarget{EMS-1842-JS-16-AUG-1842-W-c}%
				The
					\marginnote{c}%
				news you wrote me was more favorable than that which was communicated by the brethren. They seemed a little agitated for my safety and advise me for the Pine Woods. But I succeeded admirably in calming all their fears. But nevertheless as I said in my former letter, I was willing to exile myself for months and years, if it would be for the safety and welfare of the people;
					\hypertarget{EMS-1842-JS-16-AUG-1842-W-d}%
				and
					\marginnote{d}%
				I do not know but it would be as well for me to take a trip to the Pine Countries and remain untill arrangements can be made for my most perfect safety when I return. These are therefore to confer with you on this subject as I want to have a concert of action in every thing that I do.
					\hypertarget{EMS-1842-JS-16-AUG-1842-W-e}%
				If
					\marginnote{e}%
				I [k]new that they would oppress me alone, and let the rest of you dwell peaceably and quietly, I think I<t> be the wisest plan to absent myself for a little season if by that means we can prevent the profusion of blood Please write and give me your mind on that subject and all other information that has come to hand today and what are the signs of the times.

					\hypertarget{EMS-1842-JS-16-AUG-1842-W-f}%
				I
					\marginnote{f}%
				have no news for I am where I cannot get much all is quiet and peaceable around. I therefore wait with earnest expectation for your advices. I am anxious to know your opinions on any course that I may see proper to take, for in the multitude of council there is safety.

				I add no more, but subscribe myself your faithful and most obedient servant friend and brother

				\begin{tightright}
					Joseph Smith

					Lieu\supersu{t} Gen. of Nauvoo Legion of Illinois Millitia [\textit{5 lines blank}]
				\end{tightright}

				<Copy>

				<Maj Gen Law Aug\supersu{t}\supers{.} 16\supersu{th}\supers{.} 1842>
				
		% -----
		\subsubsection{Reflections and Blessings, 16 and 23 August 1842}
		% -----
			\label{EMS-1842-JS-16-23-AUG-1842}
			% \label{EMS-1842-JS-DAY-3LETTERMONTH-YEAR}
			
			\begin{description}
				\item[Print Sources]
			\end{description}

			\begin{tabular}{p{0.5in}p{1in}p{0.5in}p{1in}}
				JSP	 	& ---		& JSR 		& --- \\
				DC	 	& ---		& HC 		& --- \\
				HDDC 	& ---		& TCHC 		& --- \\
			\end{tabular}
			% HDDC = woodford's DC diss; JSR = Marquardt's JS Revelations; TCHC = Vogel HC
			
			\begin{description}
				\item[Online Access] \href{https://www.josephsmithpapers.org/paper-summary/reflections-and-blessings-16-and-23-august-1842/1}{Here}
				% \item[Analysis] \hyperref[XX]{Here}
			\end{description}
			
			\begin{description}
				\item[Background]
			\end{description}
			
				% It might also be useful to give a two sentence context of the period: e.g., "This speech was given in Nauvoo to a small congregation one month prior to the destruction of the Nauvoo Expositor. These and other tensions may have been on the prophet's mind when he gave the speech."
			
			\begin{description}
				\item[Original Source]
			\end{description}
			
				% brief description of original source material, with reference to JSP or somewhere else for more info
				% Background material: it might be helpful to have a note that says whether a given document was presented as a speech, was written in a journal, etc.
				% Also a comment here about how reliable the source might be, or what shape it is in, if there are large omissions or parts that are hard to understand, and so on.
			
			\begin{description}
				\item[Text]
			\end{description}
			
					\hypertarget{EMS-1842-JS-16-23-AUG-1842-a}%
				``Blessed
					\marginnote{a}%
				is Brother Erastus H. Derby, and he shall be blessed of the Lord; he possesses a sober mind, and a faithful heart; the snares therefore that are subsequent to befall other men, who are treacherous and rotten-hearted, shall not come nigh unto his doors, but shall be far from the path of his feet. He loveth wisdom, and shall be found possessed of her. Let their be a crown of glory, and a diadem upon his head.
					\hypertarget{EMS-1842-JS-16-23-AUG-1842-b}%
				Let
					\marginnote{b}%
				the light of eternal Truth shine forth upon his understanding; let his name be had in everlasting remembrance; let the blessings of Jehovah be crowned upon his posterity after him, for he rendered me consolation, in the lonely places of my retreat: How good, and glorious, it has seemed unto me, to find pure and holy friends, who are faithful, just and true, and whose hearts fail not; and whose knees are confirmed and do not faulter; while they wait upon the Lord, in administering to my necessities;
					\hypertarget{EMS-1842-JS-16-23-AUG-1842-c}%
				\phantom{ }
					\marginnote{c}%
				<See Page 164> [p. 135] <From Page 135> in the day when the wrath of mine enemies was poured out upon me. In the name of the Lord, I feel in my heart to bless them, and to say in the name of Jesus Christ of Nazareth that these are the ones that shall inherit eternal life. I say it by virtue of the Holy Priesthood, and by the ministering of Holy Angels, and by the gift and power of the Holy Ghost.
					\hypertarget{EMS-1842-JS-16-23-AUG-1842-d}%
				How
					\marginnote{d}%
				glorious were my feelings when I met that faithful and friendly band, on the night of the eleventh on thursday, on the Island, at the mouth of the slough, between Zarahemla and Nauvoo. With what unspeakable delight, and what transports of joy swelled my bosom, when I took by the hand on that night, my beloved Emma, she that was my wife, even the wife of my youth; and the choice of my heart.
					\hypertarget{EMS-1842-JS-16-23-AUG-1842-e}%
				Many
					\marginnote{e}%
				were the re-vibrations of my mind when I contemplated for a moment the many \sout{passt} scenes we had been called to pass through. The fatigues, and the toils, the sorrows, and sufferings, and the joys and consolations from time to time had strewed our paths and crowned our board. Oh! what a co-mingling of thought filled my mind for the moment, Again she <is> here, even in the seventh trouble, undaunted, firm and unwavering, unchangeable, affectionate Emma.
					\hypertarget{EMS-1842-JS-16-23-AUG-1842-f}%
				There
					\marginnote{f}%
				was brother Hyrum [Smith] who next took me by the hand, a natural brother; thought I to myself, brother Hyrum, what a faithful heart you have got. Oh, may the eternal Jehovah crown eternal blessings upon your head, as a reward for the care you have had for my soul. O how many are the sorrows have we shared together, and again we find ourselves shackled with the unrelenting hand of oppression. Hyrum, thy name shall be written in the Book of the Law of the Lord, for those who come after thee to look upon, that they may pattern after thy works.
					\hypertarget{EMS-1842-JS-16-23-AUG-1842-g}%
				Said
					\marginnote{g}%
				I to myself here is brother Newel K. Whitney also, how many scenes of sorrow, have strewed our paths together; and yet we meet once more to share again. Thou art a faithful friend in whom the afflicted sons of men can confide, with the most perfect safety. Let the blessings of the eternal be crowned also upon his head; how warm that heart! how anxious that soul! for the welfare of one who has been cast out, and hated of almost all men.
					\hypertarget{EMS-1842-JS-16-23-AUG-1842-h}%
				Brother
					\marginnote{h}%
				Whitney, thou knowest not how strong those ties are, that bind my soul and heart to thee. My heart was overjoyed, as I took the faithful band by hand, that stood upon the shore one by one. W\supersu{m}\supers{.} Law, W\supersu{m}\supers{.} Clayton, Dimick B. Huntington, George Miller were there. The above names constituted the little group. I do not think to mention the particulars of the history of that sacred night, which shall forever be remembered by me.
					\hypertarget{EMS-1842-JS-16-23-AUG-1842-i}%
				But
					\marginnote{i}%
				the names of the faithful are what I wish to record in this place. These I have met in prosperity and they were my friends, I now meet them in adversity, and they are still my warmer friends. These love the God that I serve; they love the truths that I promulge; they love those virtuous, and those holy doctrines that I cherish in my bosom with the warmest feelings of my heart; and with that zeal which cannot be denied.
					\hypertarget{EMS-1842-JS-16-23-AUG-1842-j}%
				I
					\marginnote{j}%
				love friendship and truth; I love virtue [p. 164] and Law; I love the God of Abraham and of Isaac and of Jacob, and they are my brethren, and I shall live; and because I live, they shall live also. These are not the only ones, who have administered to my necessity; whom the Lord will bless. There is brother John D Parker, and brother Amasa Lyman, and brother Wilson Law, and brother Henry G. Sherwood, my heart feels to reciprocate the unweried kindnesses that have been bestowed upon me by these men.
					\hypertarget{EMS-1842-JS-16-23-AUG-1842-k}%
				They
					\marginnote{k}%
				are men of noble stature, of noble hands, and of noble deeds; possessing noble and daring, and giant hearts and souls. There is brother Joseph B. Nobles [Noble] also, I would call up in remembrance before the Lord. There is brother Samuel Smith, a natural brother; he is, even as Hyrum. There is brother Arthur Millikin also, who married my youngest sister, Lucy. He is a faithful, an honest, and an upright man.
					\hypertarget{EMS-1842-JS-16-23-AUG-1842-l}%
				While
					\marginnote{l}%
				I call up in remembrance before the Lord these men, I would be doing injustice to those who rowed me in the skiff up the river that night, after I parted with the lovely group; who brought me to this my safe and lonely and private retreat; brother Jonathon [Jonathan] Dunham and the other whose name I do not know. Many were the thoughts that swelled my aching heart, while they were toiling faithfully with their oars.
					\hypertarget{EMS-1842-JS-16-23-AUG-1842-m}%
				The<y>
					\marginnote{m}%
				complained not at hardship and fatigue to secure my safety. My heart would have been harder than an adamantine stone, if I had not have prayed for them, with anxious and fervent desire. I did so, and the still small voice whispered to my soul, these that share your toils with such faithful hearts, shall reigne with you in the kingdom of their God; but I parted with them in silence and came to my retreat.
					\hypertarget{EMS-1842-JS-16-23-AUG-1842-n}%
				I
					\marginnote{n}%
				hope I shall see them again that I may toil for them and administer to their comfort also. They shall not want a friend while I live. My heart shall love those; and my hands shall toil for those, who love and toil for me, and shall ever be found faithful to my friends. Shall I be ungrateful? verily no! God forbid!['']
				
				[. . .] [p. 165] [. . . ]
				
					\hypertarget{EMS-1842-JS-16-23-AUG-1842-o}%
				``While
					\marginnote{o}%
				I contemplate the virtues and the good qualifications and characterestics of the faithful few, which I am now recording in the Book of the Law of the Lord, of such as have stood by me in every hour of peril, for these fifteen long years past; say for instance; my aged and beloved brother Joseph Knights [Knight] Sen\supers{r}, who was among the number of the first to administer to my necessities, while I was laboring, in the commencement of the bringing forth of the work of the Lord, and of laying the foundation of the Church of Jesus Christ of Latter Day Saints:
					\hypertarget{EMS-1842-JS-16-23-AUG-1842-p}%
				for
					\marginnote{p}%
				fifteen years has he been faithful and true, and even handed, and exemplary and virtuous, and kind; never deviating to the right hand nor to the left. Behold he is a righteous man. May God Almighty lengthen out the old mans days; and may his trembling, tortured and broken body be renewed, and the vigor of health turn upon him; if it can be thy will, consistently, O God; and it shall be said of him by the sons of Zion, while there is one of them remaining; that this man, was a faithful man in Israel; therefore his name shall never be forgotten.
					\hypertarget{EMS-1842-JS-16-23-AUG-1842-q}%
				There
					\marginnote{q}%
				is his son Newel Knight and Joseph Knight [Jr.] whose names I record in the Book of the Law of the Lord, with unspeakable delight, for they are my friends. There are a numerous host of faithful souls, whose names I could wish to record in the Book of the Law of the Lord; but time and chance would fail. I will mention therefore only a few of them as emblematical of those who are to numerous to be written.
					\hypertarget{EMS-1842-JS-16-23-AUG-1842-r}%
				But
					\marginnote{r}%
				there is one man I would mention namely [Orrin] Porter Rockwell, who is now a fellow-wanderer with myself--- an exile from his home because of the murderous deeds and infernal fiendish disposition of the indefatigable and unrelenting hand of the Missourians. He is an innocent and a noble boy; may God Almighty deliver him from the hands of his pursuers. He was an innocent and a noble child, and my soul loves him; Let this be recorded for ever and ever. Let the blessings of Salvation and honor be his portion.
					\hypertarget{EMS-1842-JS-16-23-AUG-1842-s}%
				But
					\marginnote{s}%
				as I said before, so say I again while I remember the faithful few who are now living, I would remember also the faithful of my friends who are dead, for they are many; and many are the acts of kindness, and paternal, and brotherly kindnesses which they have bestowed upon me. And since I have been hunted by the Missourians many are the scenes which have been called to my mind. Many thoughts have rolled through my head, and across my breast.
					\hypertarget{EMS-1842-JS-16-23-AUG-1842-t}%
				I
					\marginnote{t}%
				have remembered the scenes of my child-hood I have thought of my father who is dead; who died by disease which was brought upon him through suffering by the hands of ruthless mobs. He was a great and a good man. The envy of knaves and fools was heaped upon him, and this was his lot and portion all the days of his life. He was of noble stature, and possessed a high, and holy, and exalted, and a virtuous mind.
					\hypertarget{EMS-1842-JS-16-23-AUG-1842-u}%
				His
					\marginnote{u}%
				soul soared above all those mean [p. 179] and grovelling principles that are so subsequent to the human heart. I now say, that he never did a mean act that might be said was ungenerous, in his life, to my knowledge. I loved my father and his memory; and the memory of his noble deeds, rest with ponderous weight upon my mind; and many of his kind and parental words to me, are written on the tablet of my heart.
					\hypertarget{EMS-1842-JS-16-23-AUG-1842-v}%
				Sacred
					\marginnote{v}%
				to me, are the thoughts which I cherish of the history of his life, that have rolled through my mind and has been implanted there, by my own observation since I was born. Sacred to me is his dust, and the spot where he is laid. Sacred to me is the tomb I have made to encircle o'er his head. \sout{that} <Let> the memory of my father eternally live. \sout{Let the faults, and the follies} Let his soul, or the Spirit my follies forgive.
					\hypertarget{EMS-1842-JS-16-23-AUG-1842-w}%
				With
					\marginnote{w}%
				him may I reign one day, in the mansions above; and tune up the Lyre of Anthems, of the eternal Jove. May the God that I love look down from above, and save me from my enemies here, and take me by the hand; that on Mount Zion I may stand and with my father crown me eternally there. Words and language, is inadequate to express the gratitude that I owe to God for having given me so honorable a parentage.
					\hypertarget{EMS-1842-JS-16-23-AUG-1842-x}%
				My
					\marginnote{x}%
				mother also is one of the noblest, and the best of all women. May God grant to prolong her days, and mine; that we may live to enjoy each others society long yet in the enjoyment of liberty, and to breath the free air. Alvin my oldest brother, I remember well the pangs of sorrow that swelled my youthful bosom and almost burst my \sout{aching} <tender> heart, when he died. He was the oldest, and the noblest of my fathers family.
					\hypertarget{EMS-1842-JS-16-23-AUG-1842-y}%
				He
					\marginnote{y}%
				was one of the noblest of the sons of men: Shall his name not be recorded in this Book? Yes, Alvin; let it be had here, and be handed down upon these sacred pages, forever and ever. In him there was no guile. He lived without spot from the time he was a child. From the time of his birth, he never knew mirth. He was candid and sober and never would play; and minded his father, and mother, in toiling all day. He was one of the soberest of men and when he died the Angel of the Lord visited him in his last moments. These childish lines I record in remembrance of my child-hood scenes.
					\hypertarget{EMS-1842-JS-16-23-AUG-1842-z}%
				My
					\marginnote{z}%
				Brother Don Carloss [Carlos] Smith, whose name I desire to record also, was a noble boy. I never knew any fault in him. I never saw the first immoral act; or the first irreligious, or ignoble disposition in the child. From the time that he was born, till the time of his death; he was a lovely, a good natured, and a kind-hearted, and a virtuous and a faithful upright child. And where his soul goes let mine go also. He lays by the side of my father.
					\hypertarget{EMS-1842-JS-16-23-AUG-1842-aa}%
				Let
					\marginnote{aa}%
				my father, Don Carlos, and Alvin, and children that I have buried be brought and laid in the Tomb I have built. Let my mother, and my brethren, and my sisters be laid there also; and let it be called the Tomb of Joseph, a descendant of Jacob; and when I die, let me be gathered to the Tomb of my father. 
					\hypertarget{EMS-1842-JS-16-23-AUG-1842-bb}%
				here
					\marginnote{bb}%
				are many souls, whom I have loved stronger than death; to them I have proved faithful; to them I [p. 180] am determined to prove faithful, untill God calls me to resign up my breath. O, thou who seeeth, and knoweth the hearts of all men; thou eternal, omnipotent, omnicient, and omnipresent Jehovah, God; thou Eloheem, that sitteth, as sayeth the psalmist, enthroned in heaven; look down upon thy servant Joseph, at this time; and let faith on the name of thy Son Jesus Christ, to a greater degree than thy servant ever yet has enjoyed, be conferred upon him; even the faith of Elijah; and let the Lamp of eternal life, be lit up in his heart, never to be taken away; and let the words of eternal life, be poured upon the soul of thy servant; that he may know thy will, thy statutes, and thy commandments, and thy judgements to do them.
					\hypertarget{EMS-1842-JS-16-23-AUG-1842-cc}%
				As
					\marginnote{cc}%
				the dews upon Mount Hermon may the distillations of thy divine grace, glory and honor in the plenitude of thy mercy, and power and goodness be poured down upon the head of thy servant.
					\hypertarget{EMS-1842-JS-16-23-AUG-1842-dd}%
				O
					\marginnote{dd}%
				Lord God, my heavenly Father, shall it be in vain, that thy servant must needs be exiled from the midst of his friends; or be dragged from their bosoms, to clank in cold and iron chains; to be thrust within the dreary prison walls; to spend days of sorrow, and of grief and misery their [there], by the hand of an infuriated, insensed and infatuated foe; to glut their infernal and insatiable desire upon innocent blood; and for no other cause on the part of thy servant, than for the defence of innocence, and thou a just God will not hear his cry?
					\hypertarget{EMS-1842-JS-16-23-AUG-1842-ee}%
				O,
					\marginnote{ee}%
				no, thou wilt hear me; a child of woe, pertaining to this mortal life; because of sufferings here, but not for condemnation that shall come upon him in eternity; for thou knowest O God, the integrity of his heart. Thou hearest me, and I knew that thou wouldst hear me, and mine enemies shall not prevail; they all shall melt like wax before thy face; and as the mighty floods, and waters roar;
					\hypertarget{EMS-1842-JS-16-23-AUG-1842-ff}%
				\phantom{ }
					\marginnote{ff}%
				\sout{so} \sout{shall} or as the billowing earth-quake's, devouring gulf; or rolling Thunders loudest peal; or vivid, forked lightnings flash; or sound of the Arch-Angels trump; or voice of the Eternal God, shall the souls of my enemies be made to feel in an instant, suddenly; and shall be taken, and ensnared; and fall back-wards, and stumble in the ditch they have dug for my feet, and the feet of my friends; and perish in their own infamy and shame,--- be thrust down to an eternal hell, for their murderous and hellish deeds.'' [. . .] [p. 181]
				
		% -----
		\subsubsection{Letter to Newel K., Elizabeth Ann Smith, and Sarah Ann Whitney, 18 August 1842}
		% -----
			\label{EMS-1842-JS-18-AUG-1842}
			% \label{EMS-1842-JS-DAY-3LETTERMONTH-YEAR}
			
			\begin{description}
				\item[Print Sources]
			\end{description}

			\begin{tabular}{p{0.5in}p{1in}p{0.5in}p{1in}}
				JSP	 	& ---		& JSR 		& --- \\
				DC	 	& ---		& HC 		& --- \\
				HDDC 	& ---		& TCHC 		& --- \\
			\end{tabular}
			% HDDC = woodford's DC diss; JSR = Marquardt's JS Revelations; TCHC = Vogel HC
			
			\begin{description}
				\item[Online Access] \href{https://www.josephsmithpapers.org/paper-summary/letter-to-newel-k-elizabeth-ann-smith-and-sarah-ann-whitney-18-august-1842/1}{Here}
				% \item[Analysis] \hyperref[XX]{Here}
			\end{description}
			
			\begin{description}
				\item[Background]
			\end{description}
			
				% It might also be useful to give a two sentence context of the period: e.g., "This speech was given in Nauvoo to a small congregation one month prior to the destruction of the Nauvoo Expositor. These and other tensions may have been on the prophet's mind when he gave the speech."
			
			\begin{description}
				\item[Original Source]
			\end{description}
			
				% brief description of original source material, with reference to JSP or somewhere else for more info
				% Background material: it might be helpful to have a note that says whether a given document was presented as a speech, was written in a journal, etc.
				% Also a comment here about how reliable the source might be, or what shape it is in, if there are large omissions or parts that are hard to understand, and so on.
			
			\begin{description}
				\item[Text]
			\end{description}
			
					\hypertarget{EMS-1842-JS-18-AUG-1842-a}%
				<Nauvoo
					\marginnote{a}%
				August 18\supersu{th}\supers{.} 1842>
				
				\textbf{Dear, and Beloved, Brother [Newel K. Whitney] and Sister, [Elizabeth Ann Smith] Whitney, and \&c. [Sarah Ann Whitney]---}
				
				\textbf{I take this oppertunity to communi[c]ate, some of my feelings, privetely at this time, which I want you three Eternaly to keep in your own bosams; for my feelings are so strong for you since what has pased lately between us, that the time of my abscence from you seems so long, and dreary, that it seems, as if I could not live long in this way: and <if you> three would come and see me in this my lonely retreat, it would afford me great relief, of mind, if those with whom I am alied, do love me, now is the time to afford me succour, in the days of exile, for you know I foretold you of these things.}
					\hypertarget{EMS-1842-JS-18-AUG-1842-b}%
				\phantom{ }
					\marginnote{b}%
				\textbf{I am now at Carlos Graingers [Granger's], Just back of Brother Hyrams [Hyrum Smith's] farm, it is only one mile from town, the nights are very pleasant indeed, all three of <you> \sout{come} <can> come and see me in the fore part of the night, let Brother Whitney come a little a head, and nock at the south East corner of the house at <the> window; it <is> next to the cornfield; I have a room intirely by myself, the whole matter can be attended to with most perfect saf[e]ty, I <know> it is the will of God that you should comfort <me> now in this time of affliction, or not \sout{all} at all now is the time or never, but I hav[e] no kneed of saying any such thing, to you, for I know the goodness of your hearts, and that you will do the will of the Lord, when it is made known to you;}
					\hypertarget{EMS-1842-JS-18-AUG-1842-c}%
				\phantom{ }
					\marginnote{c}%
				\textbf{the only thing to be careful of; is to find out when Emma comes then you cannot be safe, but when she is not here, there is the most perfect \uline{safty}: only be careful to escape observation, as much as possible, I know it is a heroick undertakeing; but so much the greater frendship, and the more Joy, when I see you I <will> tell you all my plans, I cannot write them on paper, burn this letter as soon as you read it, keep all locked up in your breasts, my life depends upon it, one thing I want to see you for is <to> get the fulness of my blessings sealed upon our heads, \&c.}
					\hypertarget{EMS-1842-JS-18-AUG-1842-d}%
				\phantom{ }
					\marginnote{d}%
				\textbf{you will pardon me for my earnestness on <this subject> when you consider how lonesome I must be, your good feelings know how to <make> every allowance for me, I close my letter, I think Emma wont come to night if she dont dont fail to come to night, I subscribe myself your most obedient, <and> affectionate, companion, and friend.}
				
				\begin{tightright}
					\textbf{Joseph Smith}
				\end{tightright}
				
		% -----
		\subsubsection{Discourse, 29 August 1842}
		% -----
			\label{EMS-1842-JS-29-AUG-1842}
			% \label{EMS-1842-JS-DAY-3LETTERMONTH-YEAR}
			
			% --
			\paragraph{As Reported by William Clayton}
			% --
				\label{EMS-1842-JS-29-AUG-1842-WC}
			
				\begin{description}
					\item[Print Sources]
				\end{description}

				\begin{tabular}{p{0.5in}p{1in}p{0.5in}p{1in}}
					JSP	 	& ---		& JSR 		& --- \\
					DC	 	& ---		& HC 		& --- \\
					HDDC 	& ---		& TCHC 		& --- \\
				\end{tabular}
				% HDDC = woodford's DC diss; JSR = Marquardt's JS Revelations; TCHC = Vogel HC
				
				\begin{description}
					\item[Online Access] \href{https://www.josephsmithpapers.org/paper-summary/discourse-29-august-1842-as-reported-by-william-clayton/1}{Here}
					% \item[Analysis] \hyperref[XX]{Here}
				\end{description}
				
				\begin{description}
					\item[Background]
				\end{description}
				
					% It might also be useful to give a two sentence context of the period: e.g., "This speech was given in Nauvoo to a small congregation one month prior to the destruction of the Nauvoo Expositor. These and other tensions may have been on the prophet's mind when he gave the speech."
				
				\begin{description}
					\item[Original Source]
				\end{description}
				
					% brief description of original source material, with reference to JSP or somewhere else for more info
					% Background material: it might be helpful to have a note that says whether a given document was presented as a speech, was written in a journal, etc.
					% Also a comment here about how reliable the source might be, or what shape it is in, if there are large omissions or parts that are hard to understand, and so on.
				
				\begin{description}
					\item[Text]
				\end{description}
				
						\hypertarget{EMS-1842-JS-29-AUG-1842-WC-a}%
					``He
						\marginnote{a}
					had told them formerly about fighting the Missourians, and about fighting alone. He had not fought them with the sword nor by carnal weapons; he had done it by stratagem or by outwitting them, and there had been no lives lost, and there would be no lives lost if they would hearken to his council.
						\hypertarget{EMS-1842-JS-29-AUG-1842-WC-b}%
					Up
					\marginnote{b}%
					to this day God had given him wisdom to save the people who took council. None had ever been killed who abode by his council. At Hauns Mill the brethren went contrary to his council, if they had not there lives would have been spared. He has been in Nauvoo all the while, and outwitted Bennetts associates and attended to his own business in the City all the time.
						\hypertarget{EMS-1842-JS-29-AUG-1842-WC-c}%
					We
						\marginnote{c}%
					want to whip the world mentally and they will whip themselves physically. The brethren cant have the tricks played upon them that were done at Kirtland and Far-west, they have seen enough of the tricks of their enemies and know better''. Orson Pratt has attempted to destroy himself,--- caused all the City almost to go in search of him.
						\hypertarget{EMS-1842-JS-29-AUG-1842-WC-d}%
					Is
						\marginnote{d}%
					it not enough to put down all the infernal influence of the Devil what we have felt and seen, handled and evidenced of this work of God? But the Devil had influence among the Jews to cause the death of Jesus Christ by hanging between heaven and earth. O. P [Orson Pratt] and others of the same class caused trouble by telling stories to people who would betray me and `they must believe these stories because his wife told him so'! I will live to trample on their ashes with the soles of my feet.
						\hypertarget{EMS-1842-JS-29-AUG-1842-WC-e}%
					I
						\marginnote{e}%
					prophecy in the name of Jesus Christ that such shall not prosper, they shall be cut down in their own plans. They would deliver me up Judas like, but a small band of us shall overcome. We dont want or mean to fight with the sword of the flesh but we will fight with the broad Sword of the spirit. Our enemies say our Charter and Writs of Habeus Corpus are worth nothing. We say they came from the highest authority in the States, and we will hold to them. They cannot be disannulled or taken away.''

						\hypertarget{EMS-1842-JS-29-AUG-1842-WC-f}%
					He
						\marginnote{f}%
					then told the brethren what he was going to do, viz; to send all the Elders away and when the mob came there would only be women and children to fight and they would be ashamed, He said

					``I dont want you to fight but to go and gather tens, hundreds and thousands to fight for you. If oppression comes I will then shew them that there is a Moses and a Joshua amongst us; and I will fight them if they dont take off oppression from me, I will do as I have done this time, I will run into the woods.
						\hypertarget{EMS-1842-JS-29-AUG-1842-WC-g}%
					I
						\marginnote{g}%
					will fight them in my own way. I will send bro. Hyrum [Smith] to call conferences every where through-out the States, and let documents be taken along and show to the world the corrupt and oppressive <conduct> of [Lilburn W.] Boggs. [Thomas] Carlin and others, that the public may have the truth laid before them.
						\hypertarget{EMS-1842-JS-29-AUG-1842-WC-h}%
					Let
						\marginnote{h}%
					the Twelve send all who will support the character of the Prophet--- the Lords anointed. And if all who go will support my character I prophecy in the name of the Lord Jesus whose servant I am, that you will prosper in your missions. I have the whole plan of the kingdom before me, and no other person has.
						\hypertarget{EMS-1842-JS-29-AUG-1842-WC-i}%
					And
						\marginnote{i}%
					as to all that Orson Pratt, Sidney Rigdon or George W. Robinson can do to prevent me I can kick them off my heels, as many as you can name, I know what will become of them''. He concluded his remarks by saying ``I have the best of feelings towards my brethren since this last trouble began, but to the apostates and enemies I will give a lashing every oppertunity and I will curse them.''
					
		% -----
		\subsubsection{Discourse, 31 August 1842}
		% -----
			\label{EMS-1842-JS-31-AUG-1842}
			% \label{EMS-1842-JS-DAY-3LETTERMONTH-YEAR}
		
			\begin{description}
				\item[Print Sources]
			\end{description}

			\begin{tabular}{p{0.5in}p{1in}p{0.5in}p{1in}}
				JSP	 	& ---		& JSR 		& --- \\
				DC	 	& ---		& HC 		& --- \\
				HDDC 	& ---		& TCHC 		& --- \\
			\end{tabular}
			% HDDC = woodford's DC diss; JSR = Marquardt's JS Revelations; TCHC = Vogel HC
			
			\begin{description}
				\item[Online Access] \href{https://www.josephsmithpapers.org/paper-summary/minutes-and-discourse-31-august-1842/1}{Here}
				% \item[Analysis] \hyperref[XX]{Here}
			\end{description}
			
			\begin{description}
				\item[Background]
			\end{description}
			
				% It might also be useful to give a two sentence context of the period: e.g., "This speech was given in Nauvoo to a small congregation one month prior to the destruction of the Nauvoo Expositor. These and other tensions may have been on the prophet's mind when he gave the speech."
			
			\begin{description}
				\item[Original Source]
			\end{description}
			
				% brief description of original source material, with reference to JSP or somewhere else for more info
				% Background material: it might be helpful to have a note that says whether a given document was presented as a speech, was written in a journal, etc.
				% Also a comment here about how reliable the source might be, or what shape it is in, if there are large omissions or parts that are hard to understand, and so on.
			
			\begin{description}
				\item[Text]
			\end{description}
				
				\begin{tightright}
						\hypertarget{EMS-1842-JS-31-AUG-1842-a}%
					Grove,
						\marginnote{a}%
					August 31\supers{st.}
				\end{tightright}

				Pres\supers{t.} Joseph Smith opened the meeting by addressing the Society. He commenced by expressing his happiness and thankfulness for the privilege of being present on the occasion. He said that great exertions had been made on the part of our enemies, but they had not accomplished their purpose--- God had enabled him to keep out of their hands--- he had war'd a good warfare inasmuch as he had whip'd out all of [John C.] Bennett's host---
					\hypertarget{EMS-1842-JS-31-AUG-1842-b}%
				his
					\marginnote{b}%
				feelings at the present time were, that inasmuch as the Lord Almighty had preserv'd him today. He said it reminded him of the Savior, when he said to the pharisees, ``Go ye and tell that fox, Behold I cast out devils, and I do cures today and tomorrow, and the third day I shall be perfected.'' \&c.

				He said he expected the heavenly Father had decreed that the Missourians shall not get him--- if they do, it will be because he does not keep out of the way.

					\hypertarget{EMS-1842-JS-31-AUG-1842-c}%
				Pres\supers{t.}
					\marginnote{c}%
				S. continued by saying, I shall triumph over my enemies--- I have begun to triumph over them at home and I shall do it abroad--- all those that rise up against me will feel the weight of their iniquity upon their own heads--- those that speak evil are abominable characters--- and full of iniquity--- All the fuss and all the stir against me, is like the jack in the lantern, it cannot be found.
					\hypertarget{EMS-1842-JS-31-AUG-1842-d}%
				Altho'
					\marginnote{d}%
				I do wrong, I do not the wrongs that I am charg'd with doing--- the wrong that I do is thro' the frailty of human nature like other men. No man lives without fault. Do you think that even Jesus, if he were here would be without fault in your eyes? Th[e]y said all manner of evil against him--- they all watch'd for iniquity.

					\hypertarget{EMS-1842-JS-31-AUG-1842-e}%
				How
					\marginnote{e}%
				easy it was for Jesus to call out all the iniquity of the hearts of those whom he was among? The servants of the Lord are required to guard against those thing[s] that are calculated to do the most evil--- the little foxes spoil the vines--- little evils do the most injury to the church. If you have evil feelings and speak of them to one an other, it has a tendency to do mischief--- these things result in those evils which are calculated to cut the throats of the heads of the church.

					\hypertarget{EMS-1842-JS-31-AUG-1842-f}%
				When
					\marginnote{f}%
				I do the best I can--- when I am accomplishing the greatest good, then the most evils are got up against me. I would to God that you would be wise. I now counsel you, if you know any thing, hold your tongues, and the least harm will be done.

					\hypertarget{EMS-1842-JS-31-AUG-1842-g}%
				The
					\marginnote{g}%
				Female Relief Society has taken the the most active part in my welfare--- against my enemies--- in petitioning to the Governor--- These measures were all necessary--- Do you not see that I foresaw what was coming beforehand, by the spirit of prophesy?--- All had an influence in my redemption from the hand of my enemies.

				If these measures had not been taken, more serious consequences would have resulted.

					\hypertarget{EMS-1842-JS-31-AUG-1842-h}%
				I
					\marginnote{h}%
				have come here to bless you. The Society has done well--- their principles are to practice holiness--- God loves you and your prayers in my behalf shall avail much--- Let them not cease to ascend to God in my behalf. The enemy will never get weary--- I expect he will array every thing against me--- I expect a tremendous warfare. He that will war the christian warfare will have the angels of devils and all the infernal powers of darkness continually array'd against him.
					\hypertarget{EMS-1842-JS-31-AUG-1842-i}%
				When
					\marginnote{i}%
				wicked and corrupt men oppose, it is a criterion to judge if a man is warring the christian warfare. When all men speak evil of you, blessed are ye \&c. Shall a man be considered bad, when men speak evil of him? No: If a man stands and opposes the world of sin, he may expect all things array'd against him.

					\hypertarget{EMS-1842-JS-31-AUG-1842-j}%
				But
					\marginnote{j}%
				it will be but a little season and all these afflictions will be turn'd away from us inasmuch as we are faithful and are not overcome by these evils. By seeing the blessings of the endowment rolling on, and the kingdom increasing and spreading from sea to sea; we will rejoice that we were not overcome by these foolish things.

				Prest. S. then remark'd that a few things had been manifested to him in his absence, respecting the baptisms for the dead, which he should communicate next sabbath if nothing should occur to prevent.

					\hypertarget{EMS-1842-JS-31-AUG-1842-k}%
				Pres\supers{t.}
					\marginnote{k}%
				S. then address'd the throne of Grace.

				Pres\supers{t.} Emma Smith then rose and presented the following names, which were unanimously receiv'd to wit...

				Pres\supers{t.} J. S. remark'd that sis. Repshar had long since been advised to return to her husband--- has been ascertain'd by good evidence that she left her husband without cause--- that he is a moral man and a gentleman--- she has got into a way of having revelations, but not the rev. of God--- if she will go home we will pray for her, but if not our prayers will do no good.

					\hypertarget{EMS-1842-JS-31-AUG-1842-l}%
				Pres\supers{t.}
					\marginnote{l}%
				S. said he had one remark to make respecting the baptism for the dead--- to suffice for the time being, until he has opportunity to discuss the subject to greater length--- that is, all persons baptiz'd for the dead must have a Recorder present, that he may be an eye-witness to testify of it. It will be necessary in the grand Council, that these things be testified--- let it be attended to from this time, but if there is any lack it may be at the expense of our friends--- they may not come forth \&c.

				Prayer by br. [Erastus] Derby.

				Meeting adjourn'd...
				
		% -----
		\subsubsection{Letter to ``All the Saints in Nauvoo,'' 1 September 1842 [D\&C 127]}
		% -----
			\label{EMS-1842-JS-1-SEP-1842}
			% \label{EMS-1842-JS-DAY-3LETTERMONTH-YEAR}
		
			\begin{description}
				\item[Print Sources]
			\end{description}

			\begin{tabular}{p{0.5in}p{1in}p{0.5in}p{1in}}
				JSP	 	& ---		& JSR 		& --- \\
				DC	 	& ---		& HC 		& --- \\
				HDDC 	& ---		& TCHC 		& --- \\
			\end{tabular}
			% HDDC = woodford's DC diss; JSR = Marquardt's JS Revelations; TCHC = Vogel HC
			
			\begin{description}
				\item[Online Access] \href{https://www.josephsmithpapers.org/paper-summary/letter-to-all-the-saints-in-nauvoo-1-september-1842-dc-127/1}{Here}
				% \item[Analysis] \hyperref[XX]{Here}
			\end{description}
			
			\begin{description}
				\item[Background]
			\end{description}
			
				% It might also be useful to give a two sentence context of the period: e.g., "This speech was given in Nauvoo to a small congregation one month prior to the destruction of the Nauvoo Expositor. These and other tensions may have been on the prophet's mind when he gave the speech."
			
			\begin{description}
				\item[Original Source]
			\end{description}
			
				% brief description of original source material, with reference to JSP or somewhere else for more info
				% Background material: it might be helpful to have a note that says whether a given document was presented as a speech, was written in a journal, etc.
				% Also a comment here about how reliable the source might be, or what shape it is in, if there are large omissions or parts that are hard to understand, and so on.
			
			\begin{description}
				\item[Text]
			\end{description}
			
				\begin{tightright}
					September 1\supersu{st}\supers{.} 1842
				\end{tightright}

				To all the saints in Nauvoo

				Forasmuch as the Lord has revealed unto me that my enemies both of Mo [Missouri] \& this State were again on the pursuit of me, and inasmuch as they pursue me without cause and have not the least shadow or coloring of justice or right on their side in the getting up of their prosecutions against me; and inasmuch as their pretensions are all founded in falsehood of the blackest die. I have thought it expedient and wisdom in me to leave the place for a short season for my own safety and the safety of this people. I would say to all those with whom I have business that I have left my affairs with agents and clerks who will transact all business in a prompt and proper manner and will see that all my debts are cancelled in due time, by turning out property or otherwise as the case may require, or as the circumstances may admit of. When I learn that the storm is fully blown over then I will return to you again. And as for the perils which I am called to pass through they seem but a small thing to me, as the envy and wrath of man has been my common lot all the days of my life and for what cause it seems mysterious, unless I was ordained from before the foundation of the world for some good end, or bad as you may choose to call it. Judge ye for yourselves, God knoweth all these things wether it be good or bad, but nevertheless deep water is what I am wont to swim in, it all has become a second nature to me and I feel like Paul to glory in tribulation for unto this day has the God of my fathers delivered me out of them all and will deliver me from henceforth for behold and lo I shall triumph over all my enemies for the Lord God hath spoken it.

				Let all the saints rejoice therefore and be exceeding glad for Israels God is their God and he will meet [mete] out a just recompense of reward upon the heads of all your oppressors. And again verily thus saith the Lord let the work of my Temple and all the works which I have appointed unto you be continued on and not cease; and let your diligence and your perseverance and patience and your works be redoubled, and you shall in no wise lose your reward saith the Lord of Hosts. And if they persecute you so persecuted they the prophets and righteous men that were before you; for all this there is a reward in heaven.

				And again I give unto you a word in relation to the Baptism for your dead. Verily thus saith the Lord unto you concerning your dead when any of you are baptised for your dead let there be a recorder, and let him be eyewitness of your baptisms; let him hear with his ears that he may testify of a truth, saith the Lord; that in all your recordings it may be recorded in Heaven, that whatsoever you bind on earth may be bound in heaven; whatsoever you loose on earth may be loosed in heaven; for I am about to restore many things to the Earth, pertaining to the Priesthood saith the Lord of Hosts. And again let all the Records be had in order, that they may be put in the archives of my Holy Temple to be held in remembrance from generation to generation saith the Lord of Hosts

				I will say to all the saints that I desired with exceeding great desire to have addressed them from the stand on the subject of Baptism for the dead on the following sabbath but inasmuch as it is out of my power to do so I will write the word of the Lord from time to time on that subject and send it you by mail as well as many other things.---

				I now close my letter for the present for the want of more time, for the enemy is on the alert and as the saviour said the prince of this world cometh but he hath nothing in me. Behold my prayer to God is that you all may be saved and I subscribe myself your servant in the Lord, prophet and seer of the Church of Jesus Christ of Latter Day Saints

				Joseph Smith [\textit{1/2 page blank}]

				Mr W[illiam] Clayton

				Nauvoo

				Hancock Co

				Ill

				In Favor
				
		% -----
		\subsubsection{Proverbs, between circa 11 August and circa 4 September 1842}
		% -----
			\label{EMS-1842-JS-CIR-11-AUG-CIR-4-SEP-1842}
			% \label{EMS-1842-JS-DAY-3LETTERMONTH-YEAR}
		
			\begin{description}
				\item[Print Sources]
			\end{description}

			\begin{tabular}{p{0.5in}p{1in}p{0.5in}p{1in}}
				JSP	 	& ---		& JSR 		& --- \\
				DC	 	& ---		& HC 		& --- \\
				HDDC 	& ---		& TCHC 		& --- \\
			\end{tabular}
			% HDDC = woodford's DC diss; JSR = Marquardt's JS Revelations; TCHC = Vogel HC
			
			\begin{description}
				\item[Online Access] \href{https://www.josephsmithpapers.org/paper-summary/proverbs-between-circa-11-august-and-circa-4-september-1842/1}{Here}
				% \item[Analysis] \hyperref[XX]{Here}
			\end{description}
			
			\begin{description}
				\item[Background]
			\end{description}
			
				% It might also be useful to give a two sentence context of the period: e.g., "This speech was given in Nauvoo to a small congregation one month prior to the destruction of the Nauvoo Expositor. These and other tensions may have been on the prophet's mind when he gave the speech."
			
			\begin{description}
				\item[Original Source]
			\end{description}
			
				% brief description of original source material, with reference to JSP or somewhere else for more info
				% Background material: it might be helpful to have a note that says whether a given document was presented as a speech, was written in a journal, etc.
				% Also a comment here about how reliable the source might be, or what shape it is in, if there are large omissions or parts that are hard to understand, and so on.
			
			\begin{description}
				\item[Text]
			\end{description}
			
				\begin{tightcenter}
						\hypertarget{EMS-1842-JS-CIR-11-AUG-CIR-4-SEP-1842-a}%
					Proverbs
						\marginnote{a}%
					of Joseph
				\end{tightcenter}

				1st Never exact of a friend in adversity what you would requirre in prosperity

				2d If a man prove himsellfe to be honest in his deal, \& an ememy come upon him wickedly, through fraud or false pretences and because he is Stronger than he maketh him his prisener and Spoil him with his Goods never Say unto that man in the day of his adversity pay him what thou owest, for if though [thou] doest it though [thou] addest a deeper wound and condemnation Shall come upon the[e]. and the richus [riches] shall be Justifyed in the days of thine adversity if they mock at the[e]

					\hypertarget{EMS-1842-JS-CIR-11-AUG-CIR-4-SEP-1842-b}%
				3rd
					\marginnote{b}%
				Never afflict thy Soul for what an Enemy hath put it out of thy power to do, if thy Desieres are ever so Just

				4\supersu{th}\supers{.} Let thy hand never fale to hand out that that though [thou] owest which it is yet within thy Grasp to do So, but when thy Stock fails say to thy heart--- be Strong and to thine anxieties cease. for Man what is he, he is but dong [dung] upon the Earth and all though he demand of the the cattle of an thousan[d] hills he cannot posess himselfe of his own Life God made him and thee and gave all things in common

					\hypertarget{EMS-1842-JS-CIR-11-AUG-CIR-4-SEP-1842-c}%
				5\supersu{th}\supers{.}
					\marginnote{c}%
				There is one thing under the Sun which I have learned and that is that the ritcheousness of man is Sinn, be cause it exacteth over Much \sout{neverless} <nevertheless> the ritcheousness of God is Just because it exacteth nothing at all, but Sendeth the rain on the Just and the unjust Sead [seed] time and harvest for all of which Man is ungreatefull [\textit{1/2 page blank}]
				
		% -----
		\subsubsection{Letter to the Church, 7 September 1842 [D\&C 128]}
		% -----
			\label{EMS-1842-JS-7-SEP-1842}
			% \label{EMS-1842-JS-DAY-3LETTERMONTH-YEAR}
		
			\begin{description}
				\item[Print Sources]
			\end{description}

			\begin{tabular}{p{0.5in}p{1in}p{0.5in}p{1in}}
				JSP	 	& ---		& JSR 		& --- \\
				DC	 	& ---		& HC 		& --- \\
				HDDC 	& ---		& TCHC 		& --- \\
			\end{tabular}
			% HDDC = woodford's DC diss; JSR = Marquardt's JS Revelations; TCHC = Vogel HC
			
			\begin{description}
				\item[Online Access] \href{https://www.josephsmithpapers.org/paper-summary/letter-to-the-church-7-september-1842-dc-128/1}{Here}
				% \item[Analysis] \hyperref[XX]{Here}
			\end{description}
			
			\begin{description}
				\item[Background]
			\end{description}
			
				% It might also be useful to give a two sentence context of the period: e.g., "This speech was given in Nauvoo to a small congregation one month prior to the destruction of the Nauvoo Expositor. These and other tensions may have been on the prophet's mind when he gave the speech."
			
			\begin{description}
				\item[Original Source]
			\end{description}
			
				% brief description of original source material, with reference to JSP or somewhere else for more info
				% Background material: it might be helpful to have a note that says whether a given document was presented as a speech, was written in a journal, etc.
				% Also a comment here about how reliable the source might be, or what shape it is in, if there are large omissions or parts that are hard to understand, and so on.
			
			\begin{description}
				\item[Text]
			\end{description}
			
				\begin{tightright}
					Journeying Sep\supersu{r}\supers{.} 6\supersu{th}\supers{.} [7th] 1842
				\end{tightright}

				To the Church of Jesus Christ of Latter Day Saints Sendeth Greeting--- As I stated to you in my letter before I left my place that I would write to you from time to time and give you information in relation to many subjects. I now resume the subject of the Baptism for the dead as that subject seems to occupy my mind and press itself upon my feelings the strongest since I have been pursued by my enemies. I wrote a few words of Revelation to you, concerning a Recorder. \sout{a} I have \sout{received} <had> a few additional \sout{feelings} <views> in relation to this matter which I now Certify. i.e It was declared in my former letter that there should be a Recorder who should be eye-witness, and also to hear with his ears, that he might, make a Record of a truth before the Lord. Now in relation to this matter, it would be very difficult for one Recorder to be present at all times and to do all the business. To obviate this difficulty there can be a Recorder appointed in each ward of the City who is well qualified for taking accurate minutes and let him be very particular and precise in making his Record, in taking the whole proceeding certifying in his Record that he saw with his eyes, and heard with his ears, giving the date, and names \&c, and the history of the whole transaction, nameing also some three individuals that are present if there be any present who can at any time when called upon certify to the same that in the mouth of two or three witnesses every word may be established. Then let there be a general Recorder to whom these other records can be handed being attended with certificates over their own signatures Certifying that the Record \uline{which} they have made is true Then the General Church Recorder can enter the Record on the general Church Book with the certificates and all the attending witnesses, with his own statement that he verily believes the above statement and records to be true from his knowledge of the general characters and appointment of those men by the Church. And when this is done on the general Church Book the Record shall be just as Holy and shall answer the ordinance just the same as if he had seen with his eyes and heard with his ears and made a record of the same on the general Book You may think this order of things to be very particular but let me tell you, that they are only to answer the will of God by conforming to the ordinance and preparation that the Lord ordained and prepared before the foundation of the world for the salvation of the dead who should die without a knowledge of the Gospel. And further, I want you to remember that John the Revelator was contemplating this very subject in relation to the dead when he declared as you will find recorded in Revelations Chap 20 v 12 And I saw the dead, small and great, \sout{small an} stand before God: and the books were opened: and another Book was opened, which is the book of life; and the dead were judged out of those things which were written in the books, according to their works You will discover in this quotation that the books were opened and another book was opened which was the book of life but the dead were judged out of those things which were written in the books according to their works, consequently the books spoken off must be the books which contained the record of their \sout{worth} works and refers to the records which are kept on the earth. And the book which was the book of life is the record which is kept in heaven; the principal agreeing precisely with the doctrine which is commanded you in the revelation contained in the letters which I wrote you previous to my leaving my place ``that in all your recordings it may be recorded in Heaven''. Now the nature of this ordinance consists in the power of the priesthood by the Revelations of Jesus Christ wherein it is granted that whatsoever you bind on earth shall be bound in heaven, and whatsoever you loose on earth shall be loosed in heaven: or in other words, taking a different view of the translation; whatsoever you record on earth shall be recorded in Heaven; and whatsoever you do not record on earth, shall not be recorded in Heaven, for out of the books shall \sout{they be} <your dead be> judged according to their works \sout{wither} wether they themselves have attended to the ordinances in their own propria persona, or by the means of their own agents according to the ordinance which God has prepared for their salvation from before the foundation of the world, according to the records which they have kept concerning their dead. It may seem to some to be a very bold doctrine that we talk of; a power which records, or binds on earth, and binds in heaven. Nevertheless, in all ages of the world whenever the Lord has given a dispensatin of the priesthood to any man, by actual Revelation or any set of men this power has always been given; hence, whatsoever those men did, in authority in the name of the Lord, and did, it truly, and faithfully, and kept a proper and faithful record of the same, it became a law on earth and in Heaven and could not be annulled according to the decrees of the great Jehovah. This is a faithful saying!, who can hear it? And again for a precedent Matthew chapter 16 verses 18 \& 19 And I say also unto thee, That thou art Peter: and upon this rock I will build my Church; and the Gates of Hell shall not prevail against it. And I will give unto thee the keys of the Kingdom of Heaven; and whatsoever thou shalt bind on earth shall be bound in heaven; and whatsoever thou shalt loose on earth shall be loosed in Heaven. Now the great and grand secret of the whole matter and the sum and bonnum of the whole subject that is lying before us consists in obtaining the powers of the Holy priesthood. For him, to whom these keys are given there is no difficulty in obtaining a knowledge of facts in relation to the salvation of the children of men, both as well for the \sout{death} dead as for the living. Herein is glory, and honor, and immortality and eternal life. The ordinance of Baptism by water, to be immersed therein in order to answer to the likeness of the dead, that one principal might accord with the others. To be immersed in the water and come forth out of the water is in the likeness of the Resurrection of the dead in coming forth out of their graves; hence, this ordinance was instituted to form a relationship with the ordinance of Baptism for the dead, being in likeness of the dead. Consequently the Baptismal Font was instituted as a simile of the grave, and was commanded to be in a place underneath where the living are wont to assemble to shew forth the living and the dead; and that all things may have their likeness, and that they may accord one with another; that which is earthly, conforming to that which is heavenly as Paul hath declared 1 Corinthians Chapter 15 verses 46, 47 \& 48 Howbeit that was not first which is spiritual, but that which is natural, and afterward that which is spiritual. The first man is of the earth, earthy: the second man is the Lord from heaven. As is the earthy, such are they also that are earthy: and as is the heavenly, such are they also that are heavenly. And as are the records on the earth in relation to your dead which are truly made out so also are the records in Heaven This therefore is the sealing and binding power, and in one sense of the word the keys of the kingdom which consists in the key of knowledge. And now my dearly and beloved brethren and sisters, let me assure you that there are principals in relation to the dead and the living that cannot be lightly passed over, as pertaining to our salvation; For their salvation is necessary and essential to our salvation \sout{see} as Paul says concerning the fathers ``That they without us can not be made perfect''; Neither can we without our dead be made perfect. And now in relation, to the baptism for the dead I will give you another quotation of Paul 1 Corinthians 15 chap--- verses 29 ``Else what shall they do which are baptized for the dead, if the dead rise not at all? Why are they then baptized for the dead.'' And again in connexion with this quotation I will give you a quotation from one of the Prophets which had his eye fixid on the restoration of the priesthood the glories to be revealed in the last days, and in an especial <manner> this most glorious of all subjects belonging to the everlasting gospel viz the baptism for the dead, for Malachi says last Chapter verses 5\supersu{th} \& 6. [``]Behold I will send you Elijah the Prophet before the coming of the great and dreadful day of the Lord: and he shall turn the heart of the fathers to the children, and the heart of the children to their fathers, lest I come and smite the earth with a curse''. I might have rendered a plainer translation to this, but it is sufficiently plain, to suit my purpose as it stands. It is sufficient to know in this case that the earth will be smitten with a curse, unless there is a welding link of some kind or other, between the fathers and the children, upon some subject or other. And behold, what is that subject. It is the baptism for the dead. For we without them cannot be made perfect; neither can they without us be made perfect. Neither can they, or us, be made perfect without those who have died in the gospel also; for it is necessary in the ushering in of the dispensation of the fulness of times; which dispensation is now beginning to usher in that a whole, and compleat, and perfect union, and welding together of dispensations and keys and powers and glories should take place, and be revealed from the days of Adam even to the present time; and not only this, but \sout{that} those things which never have been revealed from the foundation of the world but have been kept hid from the wise and prudent shall be revealed unto babes and sucklings in this the dispensation of the fulness of times. Now what do we hear in the gospel which we have received? a voice of gladness! a voice of mercy from heaven! \& a voice of truth out of the earth, glad tidings for the dead; A voice of gladness for the living and the dead; glad tidings of great joy; How beautiful upon the mountains are the feet of those that bring glad tidings of good things; and that say unto Zion, behold! thy God reigneth as the dews of Carmel so shall the knowledge of God descend upon them. And again, what do we hear? Glad tidings from Cumorah! Moroni, An Angel from heaven, declaring the fulfilment of the prophets--- the book to be revealed. A voice of the Lord in the wilderness of Fayette, Senneca County, declaring the three witnesses to bear record of the book. \sout{A voice} The voice of Michael on the banks of the Susquehanna detecting the devil when he appeared as an Angel of light. The voice of Peter, James, and John in the wilderness between <Harmony> Susquehanna <County> and Colesville, Broom County; on the Susquehanna river, declaring themselves as possessing the keys of the kingdom, and of the dispensation of the fulness of times. And again, the voice of God in the chamber of old Father Whitmerss [Peter Whitmer Sr.'s] in Fayette, Senneca County and at sundry times, and in \sout{sundry} <divers> places, through[ou]t all the travels, and tribulations, of this Church of Jesus Christ of Latter Day Saints. And the voice of Michael the Archangel, the voice of Gabriel, and of Raphael, and of divers Angels from Michael or Adam, down to the present time; all declaring each one their dispensation, their rights, their keys, their honors, their majesty \& glory, and the power of their priesthood; giving line upon line; precept upon precept; here a little and there a little. Giving us consolation by holding forth that which is to come and confirming our hope.

				Brethren, shall we not go on in so great a cause Go forward and not backward. Courage brethren; and on---on to the victory. Let your hearts rejoice and be exceeding glad. Let the earth break forth into singing. Let the dead speak forth anthems of eternal praise to the king Immanuel; who hath ordained before the world was that which would enable us to redeem them out of their prisons; for the prisoner shall go free. Let the mountains shout for joy and all ye valleys cry aloud; and all ye Seas and dry lands tell the wonders of your eternal king. And ye rivers, and brooks, and rills, flow down with gladness. Let the woods, and all the trees of the field praise the Lord; and ye solid rocks, \sout{leap} weep for joy. \sout{And let all creation} And let the Sun, Moon and <the morning> Stars, \sout{strike hands} <sing> together and let all the sons of God shout for joy. And let the eternal creations declare his name for-ever and ever. And again I say, how glorious is the voice we hear, from heaven proclaiming in our ears, glory, and salvation, and honor, and immortality and eternal life; kingdoms, principalities and powers. Behold, the great day of the Lord is at hand, and who can abide the day of his coming and who can stand when he appeareth for he is like a refiners fire and like fullers soap; and he shall sit as a refiner and purifier of silver, and he shall purify the sons of Levi and purge them as Gold \& silver, that they may offer unto the Lord an offering in righteousness Let us therefore, as a church and a people, and as Latter Day saints, offer unto the Lord an offering in righteousness. And let us, present in his Holy Temple when it is finished, a book, containing the Records of our dead, which shall be worthy of all acceptation.

				Brethren, I have many things to say to you on the subject; but shall now close for the present and continue on the subject another time

				I am as ever your <humble> servant and never deviating friend

				\begin{tightright}
					\uline{\textbf{Joseph Smith}} [\textit{3/4 page blank}]
				\end{tightright}

				Will\supersu{m}\supers{.} Clayton

				Nauvoo

				Hancock Co

				\uline{Ill---}
				
		% -----
		\subsubsection{Letter to James Arlington Bennet, 8 September 1842}
		% -----
			\label{EMS-1842-JS-8-SEP-1842}
			% \label{EMS-1842-JS-DAY-3LETTERMONTH-YEAR}
		
			\begin{description}
				\item[Print Sources]
			\end{description}

			\begin{tabular}{p{0.5in}p{1in}p{0.5in}p{1in}}
				JSP	 	& ---		& JSR 		& --- \\
				DC	 	& ---		& HC 		& --- \\
				HDDC 	& ---		& TCHC 		& --- \\
			\end{tabular}
			% HDDC = woodford's DC diss; JSR = Marquardt's JS Revelations; TCHC = Vogel HC
			
			\begin{description}
				\item[Online Access] \href{https://www.josephsmithpapers.org/paper-summary/letter-to-james-arlington-bennet-8-september-1842/1}{Here}
				% \item[Analysis] \hyperref[XX]{Here}
			\end{description}
			
			\begin{description}
				\item[Background]
			\end{description}
			
				% It might also be useful to give a two sentence context of the period: e.g., "This speech was given in Nauvoo to a small congregation one month prior to the destruction of the Nauvoo Expositor. These and other tensions may have been on the prophet's mind when he gave the speech."
			
			\begin{description}
				\item[Original Source]
			\end{description}
			
				% brief description of original source material, with reference to JSP or somewhere else for more info
				% Background material: it might be helpful to have a note that says whether a given document was presented as a speech, was written in a journal, etc.
				% Also a comment here about how reliable the source might be, or what shape it is in, if there are large omissions or parts that are hard to understand, and so on.
			
			\begin{description}
				\item[Text]
			\end{description}
			
					\hypertarget{EMS-1842-JS-8-SEP-1842-a}%
				<Nauvoo
					\marginnote{a}%
				Sep. 2\supers{d} 1842>

				\begin{tightright}
					Nauvoo Sep\supersu{r}\supers{.} 8\supersu{th}\supers{.} 1842
				\end{tightright}

				Dear Sir---

				I have just received your very consoling letter dated August 16\supersu{th}\supers{.} 1842; which, I think, is the first letter you ever addressed to me. In which you speak of the arrival of Dr W[illard] Richards, and of his person, very respectfully.
					\hypertarget{EMS-1842-JS-8-SEP-1842-b}%
				In
					\marginnote{b}%
				this, I rejoice for I am as warm a friend to D\supersu{r}\supers{.} Richards, as he possibly can be to me. And in relation to his almost making a Mormon of yourself, it puts me in mind of the saying of Paul in his reply to Agrippa Acts chapter 26 verse 29 ``I would to God that not only thou, but also all that hear me this day, were both almost and altogether such as I am, except these bonds.'' And I will here remark, my Dear Sir; that Mormonism, is the pure doctrine of Jesus Christ, of which I myself am not ashamed. You speak also, of Elder [Lucian R.] Foster, president of the Church in New York, in high terms; and of D\supersu{r}\supers{.} [John M.] Bernhisel of New York.
					\hypertarget{EMS-1842-JS-8-SEP-1842-c}%
				These
					\marginnote{c}%
				men I am acquainted, with, by information, and it warms my heart to know that you speak well of them; and as you say could be willing to associate with them forever, if you never join [\textit{illegible}] their church, or acknowledged their faith.
					\hypertarget{EMS-1842-JS-8-SEP-1842-d}%
				This,
					\marginnote{d}%
				is a good principle for when we see virtuous qualities in men, we should always acknowledge them, let their understanding be what it may, in relation to creeds and doctrine; for all men are, or ought to be, free; possessing inalienable rights, and the high and noble qualifications of the laws of nature, and of self preservation; to think, and act, and say as they please while they maintain a due respect, to the rights, and privileges of all other creatures; infringing upon none.
					\hypertarget{EMS-1842-JS-8-SEP-1842-e}%
				This
					\marginnote{e}%
				doctrine, I do most heartily subscribe to, and practise; the testimony of mean men to the contrary, notwithstanding. But Sir, I will assure you, that my soul, soars far above all the mean, and grovelling dispositions of men, that are disposed to abuse me and my character; I therefore shall not dwell upon that subject.
					\hypertarget{EMS-1842-JS-8-SEP-1842-f}%
				In
					\marginnote{f}%
				relation to those men \sout{the} you speak of, referred to above; I will only say, that there are thousands of such men in this church; who, \sout{to be} <if a man is> found worthy to associate with, will call down the envy of a mean world, because of their high and noble demeanor; and it is with unspeakable delight, that I contemplate them as my friends and brethren. I love them with a perfect love; and I hope they love me, and have no reason to doubt but they do.

					\hypertarget{EMS-1842-JS-8-SEP-1842-g}%
				The
					\marginnote{g}%
				next in consideration, is John C. Bennett. I was, his friend, I am yet, his friend, as I feel myself bound to be a friend to all the sons of Adam, wether they are just or unjust; they have a degree of my compassion \& sympathy. If he is my enemy, it is his own fault; and the responsibility rests upon his own head, and instead of \sout{railing} <reigning> his character before you;
					\hypertarget{EMS-1842-JS-8-SEP-1842-h}%
				suffice
					\marginnote{h}%
				it to say that his own conduct wherever he goes, will be sufficient to recommend him to an enlightened public, wether for a bad man, or a good man. Therefore, whosever will associate themselves with him, may be assured that I will not persecute them; but I do not wish their association; and what I have said may suffice on that subject, so far as his character is concerned.

					\hypertarget{EMS-1842-JS-8-SEP-1842-i}%
				Now
					\marginnote{i}%
				in relation to his book, that he may write, I will venture a prophecy, that whoever has any hand in the matter, will find themselves in a poor fix, in relation to the money matters. And as to my having any fears of the influence that he may have against me, or any other man, or set of men may have is the most foreign from my heart;
					\hypertarget{EMS-1842-JS-8-SEP-1842-j}%
				for
					\marginnote{j}%
				I never knew what it was, as yet, to fear the face of clay, or the influence of man. My fear, Sir, is before God. I fear to offend him, and strive to keep his commandments. I am realy glad that you did not join John C. in relation to his book; from the assurances which I have, that it will prove a curse to all those who touch it.

					\hypertarget{EMS-1842-JS-8-SEP-1842-k}%
				In
					\marginnote{k}%
				relation to the honors that you speak of, both for yourself, and for Mr J. G. [James Gordon] Bennett of the Herald, you are <both> strangers to me; and as John C. Bennett kept all his letters which he received from you, entirely to himself; and there was no corresponce between you and me that I knew of, I had no opportunity to share very largely, in the getting up of any of those matters. I could not, as I had not sufficient knowledge to enable me to do so.
					\hypertarget{EMS-1842-JS-8-SEP-1842-l}%
				The
					\marginnote{l}%
				whole therefore, was at the instigation of John C. Bennett, and a quiet submission on the part of the rest, out of the best of feelings. But as for myself it was all done at a time when I was overwhelmed with a great many business cares, as well as the care of all the churches. I must be excused therefore for any wrongs that may have taken place, in relation to this matter. And so far as I obtain a knowledge of that which is right, shall meet with my hearty approval.
					\hypertarget{EMS-1842-JS-8-SEP-1842-m}%
				I
					\marginnote{m}%
				feel to tender you my most hearty \& sincere thanks for every expression of kindness you have tendered toward me or my brethren, and would beg the privilege of obtruding myself a little while upon your patience in offering a short relation of my circumstances. I am at this time persecuted the worst of any man on the earth; as well as this people, here in this place; and all our sacred rights, are trampled under the feet of the mob. I am now hunted as an hart by the mob, under the pretense, or shadow of law, to cover their abominable deeds.
					\hypertarget{EMS-1842-JS-8-SEP-1842-n}%
				An
					\marginnote{n}%
				unhallowed demand has been made from the Governor of Missouri on oath of Governor [Lilburn W.] Boggs, that I made an attempt to assassinate him, on the night of the sixth of May, when on that day <and on the seventh> it is well known <\sout{that} I was attending the officers drill and answered to my name when the roll was called> by the thousands that assembled here in Nauvoo, that I was at my post in reviewing the Nauvoo Legion in the presence of twelve thousand people; and the Gov. of the State of Illinois <notwithstanding his being> knowing to all these facts, <\uline{yet} he> immediately granted a writ, and by an unhallowed usurption, has taken away our chartered rights, and denied the right of Habeus Corpus;
					\hypertarget{EMS-1842-JS-8-SEP-1842-o}%
				and
					\marginnote{o}%
				<I am informed> has now about thirty of the most blood-thirsty kind of men <\sout{from Missouri} now on their way to> \sout{in} this place in search for me, threatening death, and destruction, and extermination upon all the Mormons; and searching my house almost continually from day to day, menacing, and threatening, and intimidating an innocent wife and children; and insulting them in a most diabolical manner, threatening their lives \&c, if I am not to be found; with a gang of Missourians with them, saying, they will have me dead or alive, and if alive, they will carry me to Missouri in chains; and when there, they will kill me at all hazards.
					\hypertarget{EMS-1842-JS-8-SEP-1842-p}%
				And
					\marginnote{p}%
				all this is backt up, and urged on, by the Gov. of this State, with all the rage of a demon, putting at defiance the constitution of this State--- our chartered rights, and the constitution of the United States; for not as yet, have they done one thing that was in accordance to them; while all the citizens of this city, \uline{en} \uline{masse}, have petitioned the Governor with remonstrances, and overtures, that would have melted the heart of an adamantine, to no effect.
					\hypertarget{EMS-1842-JS-8-SEP-1842-q}%
				And
					\marginnote{q}%
				at the same time, if any of us open our mouths, to plead our own cause, in the defence of law and justice, \sout{with} <we> are instantly threatened with Militia \& extermination Great God! When shall the oppressor cease to prey and glut itself upon innocent blood. Where is patriotism? Where is liberty? Where is the boast of this proud, and haughty nation? O humanity! where hast thou fled? hast thou fled <forever.>

					\hypertarget{EMS-1842-JS-8-SEP-1842-r}%
				I
					\marginnote{r}%
				now appeal to you, sir, inasmuch as you have subscribed yourself our friend; will you lift your voice, and your arm, with indignation, against such unhallowed oppression? I must say sir, that my bosom swells, with unutterable anguish, when I contemplate the scenes of horror that we have passed through in the State of Missouri and then look, and behold, and see the storm, and cloud, gathering ten times blacker--- ready to burst upon the heads of this innocent people. Would to God that I were able to throw off the yoke.
					\hypertarget{EMS-1842-JS-8-SEP-1842-s}%
				Shall
					\marginnote{s}%
				we bow down and be slaves? Is there no friends of humanity, in a nation that boasts itself so much? Will not the Nation rise up and defend us? If they will not defend us, will they not grant, to lend a voice of indignation, against such unhallowed oppression? Must the tens of thousands bow down to slavery and degradation? Let the pride of the nation arise, and wrench these shackles from the feet of their fellow citizens; and their quiet and peaceable, and innocent and loyal subjects.
					\hypertarget{EMS-1842-JS-8-SEP-1842-t}%
				But
					\marginnote{t}%
				I must forbear, for I cannot express my feelings. The Legion, would all willingly die in the defence of their rights; but what would this accomplish? I have kept down their indignation, and kept a quiet submission on all hands, and am determined to do so, at all hazards. Our enemies shall not have it to say, that we rebel against government, or commit treason; however much they may lift their hands in oppression, and tyranny, when it comes in the form of Government.
					\hypertarget{EMS-1842-JS-8-SEP-1842-u}%
				We
					\marginnote{u}%
				tamely submit, although it lead us to the slaughter, and to beggary; but our blood be upon their garments, and those who look tamely on, and boast of patriotism shall not without their condemnation. And if men are such fools, as to let once, the precedent be established, and through their prejudices, give assent, to such abomination, then let the oppressors hand lay heavily throughout the world, untill all flesh shall feel it together; and untill they may know, that the Almighty takes cognizance of such things;
					\hypertarget{EMS-1842-JS-8-SEP-1842-v}%
				and
					\marginnote{v}%
				then shall church rise up against church; and party against party; mob against mob; oppressor against oppressor; army against army; and kingdom against kingdom; and people against people; and kindred against kindred. And where, sir, will be your safety, or the safety of your children. If my children can be led to the slaughter with impunity, by the hands of murderous rebels; will they not lead yours to the slaughter, with the same impunity? 
					\hypertarget{EMS-1842-JS-8-SEP-1842-w}%
				Ought
					\marginnote{w}%
				not then, this oppression Sir, to be checked in the bud; and to be look'd down, with just indignattion by an enlightened world; before the \sout{flame} <fire> become unextinguishable, and, the fire devours the stubble. But again I say, I must forbear; and leave this painful subject. I wish you would write to me in answer to this and let me know your views. On my part, I am ready to be offered up a sacrifice, in that way that can bring to pass, the greatest benifit, and good, to those who must necessarily be interested, in this important matter.
					\hypertarget{EMS-1842-JS-8-SEP-1842-x}%
				\phantom{ }
					\marginnote{x}%
				\sout{I have dictated this letter, while my clerk is writing for me;} and I would to God that you could know all my feelings on this subject and the real facts in relation to this people, and their unrelenting persecution. And if any man, feels an interest in the welfare of their fellow-beings, and would think of saying or doing any thing in this matter, I would suggest the propriety of a committee of wise men being sent, to ascertain the justice or injustice of our cause to get in possession of all the facts; and then make report to an enlightened world, wether we individually, or collectively, are deserving such high-handed treatment.

					\hypertarget{EMS-1842-JS-8-SEP-1842-y}%
				In
					\marginnote{y}%
				relation to the books that you sent here, John C. Bennett put them into my store, to be sold on commission saying, that when I were able, the money must be remitted to yourself. Nothing was said about any consecration to the Temple.

				Another calamity has befallen us; our Post Office in this place is exceedingly corrupt. It is with great difficulty that we can get our letters to or from our friends. Our letters are broken open and robbed of their contents--- Our papers that we send to our sub[s]cribers, are embezzeled, and burned or wasted.
					\hypertarget{EMS-1842-JS-8-SEP-1842-z}%
				We
					\marginnote{z}%
				get no money from our subcribers, and very little information from abroad; and what little we do get, we get by private means, in consequence of these things. And I am sorry to say, that this robbing of the Post Office--- of money was carried on by John C. Bennett, and since he left here, it is carried on by the means of his confederate

				I now subscribe myself your friend, and a patriot and lover of my country, pleading at their feet for protection, and deliverance by the justice of their constitutions. I add no more. Your most obedient servant.

				\begin{tightright}
					Joseph Smith.
				\end{tightright}

				P. S. I have dictated this letter while my clerk is writing for me
				
		% -----
		\subsubsection{Discourse, 25 September 1842}
		% -----
			\label{EMS-1842-JS-25-SEP-1842}
			% \label{EMS-1842-JS-DAY-3LETTERMONTH-YEAR}
			
			% --
			\paragraph{As Reported by David Moore}
			% --
				\label{EMS-1842-JS-25-SEP-1842-DM}
		
				\begin{description}
					\item[Print Sources]
				\end{description}

				\begin{tabular}{p{0.5in}p{1in}p{0.5in}p{1in}}
					JSP	 	& ---		& JSR 		& --- \\
					DC	 	& ---		& HC 		& --- \\
					HDDC 	& ---		& TCHC 		& --- \\
				\end{tabular}
				% HDDC = woodford's DC diss; JSR = Marquardt's JS Revelations; TCHC = Vogel HC
				
				\begin{description}
					\item[Online Access] \href{https://www.josephsmithpapers.org/paper-summary/discourse-25-september-1842-as-reported-by-david-moore/1}{Here}
					% \item[Analysis] \hyperref[XX]{Here}
				\end{description}
				
				\begin{description}
					\item[Background]
				\end{description}
				
					% It might also be useful to give a two sentence context of the period: e.g., "This speech was given in Nauvoo to a small congregation one month prior to the destruction of the Nauvoo Expositor. These and other tensions may have been on the prophet's mind when he gave the speech."
				
				\begin{description}
					\item[Original Source]
				\end{description}
				
					% brief description of original source material, with reference to JSP or somewhere else for more info
					% Background material: it might be helpful to have a note that says whether a given document was presented as a speech, was written in a journal, etc.
					% Also a comment here about how reliable the source might be, or what shape it is in, if there are large omissions or parts that are hard to understand, and so on.
				
				\begin{description}
					\item[Text]
				\end{description}
				
					Meeting being now open he [Joseph Smith] arose and said that he felt to prophesy some this morning as he had of late been obliged to keep out of sight on account of his unrelenting persecutors, he therefore felt to say to the people, that inasmuch as they would keep the Commandment of God, they should never be drove from their habitations in Nauvoo, but that he would not promise that they would not be coaxed to leave.
					
		% -----
		\subsubsection{Notice, 15 November 1842}
		% -----
			\label{EMS-1842-JS-15-NOV-1842}
			% \label{EMS-1842-JS-DAY-3LETTERMONTH-YEAR}
		
			\begin{description}
				\item[Print Sources]
			\end{description}

			\begin{tabular}{p{0.5in}p{1in}p{0.5in}p{1in}}
				JSP	 	& ---		& JSR 		& --- \\
				DC	 	& ---		& HC 		& --- \\
				HDDC 	& ---		& TCHC 		& --- \\
			\end{tabular}
			% HDDC = woodford's DC diss; JSR = Marquardt's JS Revelations; TCHC = Vogel HC
			
			\begin{description}
				\item[Online Access] \href{https://www.josephsmithpapers.org/paper-summary/notice-15-november-1842/1}{Here}
				% \item[Analysis] \hyperref[XX]{Here}
			\end{description}
			
			\begin{description}
				\item[Background]
			\end{description}
			
				% It might also be useful to give a two sentence context of the period: e.g., "This speech was given in Nauvoo to a small congregation one month prior to the destruction of the Nauvoo Expositor. These and other tensions may have been on the prophet's mind when he gave the speech."
			
			\begin{description}
				\item[Original Source]
			\end{description}
			
				% brief description of original source material, with reference to JSP or somewhere else for more info
				% Background material: it might be helpful to have a note that says whether a given document was presented as a speech, was written in a journal, etc.
				% Also a comment here about how reliable the source might be, or what shape it is in, if there are large omissions or parts that are hard to understand, and so on.
			
			\begin{description}
				\item[Text]
			\end{description}
			
				\begin{tightcenter}
					VALEDICTORY.
				\end{tightcenter}
				
				I beg leave to inform the subscribers of the Times and Seasons that it is impossible for me to fulfil the arduous duties of the editorial department any longer. The multiplicity of other business that daily devolves upon me, renders it impossible for me to do justice to a paper so widely circulated as the Times and Seasons. I have appointed Elder John Taylor, who is less encumbered and fully competent to assume the responsiblities of that office, and I doubt not but that he will give satisfaction to the patrons of the paper. As this number commences a new volume, it also commences his editorial career.
				
				\begin{tightright}
					JOSEPH SMITH.
				\end{tightright}
		
		% -----
		\subsubsection{Letter to ``Hands in the Stone Shop,'' 21 December 1842}
		% -----
			\label{EMS-1842-JS-21-DEC-1842}
			% \label{EMS-1842-JS-DAY-3LETTERMONTH-YEAR}
		
			\begin{description}
				\item[Print Sources]
			\end{description}

			\begin{tabular}{p{0.5in}p{1in}p{0.5in}p{1in}}
				JSP	 	& ---		& JSR 		& --- \\
				DC	 	& ---		& HC 		& --- \\
				HDDC 	& ---		& TCHC 		& --- \\
			\end{tabular}
			% HDDC = woodford's DC diss; JSR = Marquardt's JS Revelations; TCHC = Vogel HC
			
			\begin{description}
				\item[Online Access] \href{https://www.josephsmithpapers.org/paper-summary/letter-to-hands-in-the-stone-shop-21-december-1842/1}{Here}
				% \item[Analysis] \hyperref[XX]{Here}
			\end{description}
			
			\begin{description}
				\item[Background]
			\end{description}
			
				% It might also be useful to give a two sentence context of the period: e.g., "This speech was given in Nauvoo to a small congregation one month prior to the destruction of the Nauvoo Expositor. These and other tensions may have been on the prophet's mind when he gave the speech."
			
			\begin{description}
				\item[Original Source]
			\end{description}
			
				% brief description of original source material, with reference to JSP or somewhere else for more info
				% Background material: it might be helpful to have a note that says whether a given document was presented as a speech, was written in a journal, etc.
				% Also a comment here about how reliable the source might be, or what shape it is in, if there are large omissions or parts that are hard to understand, and so on.
			
			\begin{description}
				\item[Text]
			\end{description}
			
				\begin{tightright}
						\hypertarget{EMS-1842-JS-21-DEC-1842-a}%
					Nauvoo
						\marginnote{a}%
					Dec\supersu{r}\supers{.} 21--- 1842
				\end{tightright}

				To the hands in the Stone Shop.

				Whereas, an appeal has been made to me as Sole Trustee in Trust for the Church of Jesus Christ of Latter Day Saints for a decision in relation to sundry matters in regulating the stone cutting \&c. \&c. by the Temple Committee. I have duly considered their complaints, and heard all their arguments in relation to the matter, and am satisfied that a proper deference has not been paid to their high standing by some or many of the hands in the stone Shop. And further, that their policy in relation to the Pork, and Beef, and Provisions, is for the furtherance of the building of the Temple, in the ultimatum thereof.
					\hypertarget{EMS-1842-JS-21-DEC-1842-b}%
				These
					\marginnote{b}%
				are therefore to advise you, to submit patiently to their economy, and instructions; and that we, with one accord--- with united feelings, submit patiently to the yoke that is laid upon us, and thereby secure the best interests, to the Temple of the most High God, that our limited circumstances can possibly admit of; and then having done all on our part, that great Eloheem, who has commanded us to build an house shall abundantly bless us, and reward us for all our pains.

				\begin{tightright}
					I am sirs, your sincere friend and brother and fellow sufferer in the bonds of the great work.

					Joseph Smith, Sole Trustee in Trust for the church
				\end{tightright}